%% OVERFORE-Subdocumento4.tex — generado desde OVERFORE-Subdocumento4.docx
%% Formato: clase overfore, modo subdoc. El formato lo pone la clase; aquí sólo va el contenido.
\documentclass[subdoc]{../../comun/overfore}
\numsubdoc{4}
\nombresubdoc{Arquitectura lógica y física de la solución}
\informe{Informe Preparatorio N° 2}
\fecha{12 de octubre de 2026}
\externos{A4,S3,T10}
\begin{document}

\capitulo{Introducción a la Arquitectura lógica y física de la solución}
Este capítulo traduce el Esquema de Solución del Subdocumento 3 en la arquitectura que se construye y se instala. La sección 4.1 presenta la arquitectura lógica: las ocho capas, los nueve módulos y las dos plataformas con los mismos nombres del Subdocumento 3, la integración con los diez sistemas existentes, la seguridad y la operación sin enlace de la sección 4.1.1 justifica las tecnologías de software de las capas con su alternativa, su versión y su fin de soporte, los componentes HiGHS, LightGBM y Datafaker se declaran en los Subdocumentos 9 y 13, y ocho registros de decisión de arquitectura (ADR-01 a ADR-08) documentan las decisiones que condicionan el diseño. La sección 4.2 asigna cada componente a la nube, a la sala técnica de la faena o a un gabinete de borde, describe los servicios de Azure, los ambientes, las redes, la continuidad y el dimensionamiento, la sección 4.2.1 resume los implementos que se detallan en el Formulario T-11. La sección 4.3 presenta los centros de datos: la sala técnica principal de la faena, la región primaria Azure Chile Central y la región secundaria Azure Brazil South.

El modelo y la gestión de los datos que guardan los almacenes de la capa 6 se desarrollan en el Subdocumento 5. Los códigos RF, RNF, SUP, RES y RP que se citan corresponden al catálogo de requerimientos, supuestos y restricciones del subdocumento 2 y el Formulario T-12 indica para cada requerimiento la sección donde se acredita.

\seccion{Arquitectura lógica}
La arquitectura lógica de SITCO Aranda se organiza en ocho capas. Los canales acceden a los nueve módulos por puertas de enlace. Los módulos conservan sus propios datos y se comunican mediante interfaces y eventos, mientras el Integrador de Sistemas los conecta con las aplicaciones que la compañía mantiene. Seguridad y observabilidad se aplican a todas las capas. La elección tecnológica y el estilo de servicios modulares permiten ejecutar las funciones críticas también en el Nodo de Faena cuando se pierde el enlace con Antofagasta.

\subseccion{Especificaciones Tecnologías de Software a utilizar}
Las tecnologías de software de SITCO Aranda se eligieron tras comparar cada opción con una alternativa, según su vigencia, operabilidad y portabilidad durante los 56 meses del contrato. Se prefieren servicios administrados para reducir la carga del equipo de TI, y los módulos que operan sin enlace usan el mismo contenedor en la nube y en el Nodo de Faena. La vigencia se controla con fechas de fin de soporte publicadas y actualizaciones planificadas antes de su vencimiento. Estas decisiones se aplican a las capas de presentación, borde y puertas de enlace, servicios de negocio, integración, datos, seguridad y observabilidad, además de los sistemas operativos.

\subsubseccion{Tecnologías de las capas 1 a 3: presentación, borde y puertas de enlace}
El portal y los equipos de terreno comparten interfaz, pero llegan a los servicios por rutas distintas. El tráfico público pasa por un borde protegido, las solicitudes operacionales deben alcanzar el Nodo incluso durante un corte. En la Tabla \ref{tab:selecciones-principales-capas-1-3} se distingue qué queda en el equipo, qué se resuelve en el acceso público y qué atiende el Nodo.

\begin{tabla}{I{0.171}I{0.220}I{0.306}I{0.303}}{Selecciones principales de las capas 1 a 3.}{tab:selecciones-principales-capas-1-3}
  {\cab{Capa} & \cab{Tecnología elegida} & \cab{Uso en la solución} & \cab{Alternativa evaluada y criterio de decisión}}
  {Fuente: elaboración propia con información de Angular (s.f.), Microsoft (s.f.-c) y Zetetic (s.f.).}
 & Angular 22 & Presenta la misma interfaz en portales y equipos de terreno & React añade librerías que el equipo tendría que mantener \\
 & .NET 10 LTS con Microsoft Edge WebView2 & Ejecuta la interfaz compartida en las aplicaciones instaladas & Aplicación nativa completa: duplica las pantallas \\
\multirow{-3}{=}{1. Presentación} & SQLite con SQLCipher 4 & SQLite guarda los registros pendientes y SQLCipher los cifra en la tableta & Almacenamiento del navegador: puede ser desalojado \\
 & Azure Front Door Premium con WAF y Azure DDoS Network Protection & Filtra solicitudes web y tráfico de ataque antes de los servicios & Proxy inverso propio: recibe el ataque dentro de la red \\
\multirow{-2}{=}{2. Borde y exposición} & Microsoft Entra Private Access & Restringe el Portal Web SITCO al personal autorizado & Red privada virtual (VPN): abre la red, no solo el portal. \\
 & Azure API Management & Publica interfaces externas con cuotas y catálogo & Desarrollo propio de cuotas y catálogo: exige mantenimiento \\
\multirow{-2}{=}{3. Puertas de enlace} & Envoy 1.39 & Recibe solicitudes operacionales durante un corte & Azure API Management: depende de la nube y del enlace \\
\end{tabla}
La tableta puede conservar un hecho antes de entregarlo al Nodo. Desde allí, la puerta local mantiene la recepción aunque la entrada pública quede fuera de alcance. El intercambio con la nube se retrasa, la captura en terreno continúa.

\subsubseccion{Tecnologías de las capas 4 a 6: servicios de negocio, integración y datos}
Los cinco módulos que operan en el Nodo necesitan versiones compatibles con la nube y una forma de conservar los eventos mientras falta el enlace. La integración también debe recibir señales del historiador, enviar avisos y comprobar contratos antes de desplegar cambios. Esas funciones se reúnen en la Tabla \ref{tab:tecnologias-capas-4-5-alternativa-evaluada}.

\begin{tabla}{I{0.143}I{0.285}I{0.275}I{0.297}}{Tecnologías de las capas 4 y 5 y alternativa evaluada.}{tab:tecnologias-capas-4-5-alternativa-evaluada}
  {\cab{Capa} & \cab{Tecnología y versión} & \cab{Uso en la solución} & \cab{Alternativa evaluada y criterio de decisión}}
  {Fuente: elaboración propia con información de Eclipse Adoptium (s.f.), Spring (s.f.), Kubernetes (s.f.) y Microsoft (s.f.-a, s.f.-h).}
 & Java 25 LTS (Eclipse Temurin) y Spring Boot 4.1 & Java ejecuta las reglas y Spring Boot estructura los nueve módulos & .NET 10 en el servidor: pierde soporte antes de terminar el contrato \\
 & Kubernetes 1.36: Azure Kubernetes Service (AKS) en la nube y K3s, distribución liviana de Kubernetes, en la faena & AKS ejecuta contenedores en nube y K3s, en faena & Máquinas virtuales por servicio: exigen preparar cada sitio por separado \\
 & Azure Managed Redis & Conserva el estado compartido fuera de cada servicio & Azure Cache for Redis: se retira durante el contrato \\
 & springdoc-openapi 3.1 & Genera desde el código los contratos de las interfaces & OpenAPI escrito a mano: puede diferir del servicio real \\
\multirow{-5}{=}{4. Servicios de negocio} & Springwolf 2.8 & Genera desde el código la descripción de los eventos & AsyncAPI escrito a mano: puede quedar desactualizado \\
 & Azure Service Bus Premium & Distribuye eventos entre módulos y separa los fallidos & Azure Event Hubs con interfaz Kafka: sin cola de fallidos nativa \\
 & Bandeja de salida sobre PostgreSQL 18 & Retiene eventos en faena mientras falta el enlace & Broker de mensajería local: suma otro motor que operar \\
 & PTC Kepware Server 7.1 & Lee señales del historiador mediante OPC UA & Cliente OPC UA propio: exige mantener el protocolo \\
 & Azure Communication Services y un agregador de SMS con cobertura en Chile & Envían correo, mensajería y SMS de las notificaciones & Servidor de correo propio: no cubre SMS ni mensajería \\
 & Pact JVM 4.7 & Prueba que los consumidores acepten las interfaces nueva & Comparación del contrato generado con el aprobado: no prueba a los consumidores \\
\multirow{-6}{=}{5. Integración y eventos} & Pact Broker 2.121 sobre PostgreSQL 18 & Registra las versiones de contrato en nube y Nodo & Registro manual: no bloquea versiones incompatibles \\
\end{tabla}
El dato y su evento se registran en una misma transacción local. Cuando regresa el enlace, el evento sale al bus, que conserva el orden por agregado y aparta los mensajes fallidos. Antes de desplegar una versión se prueban los consumidores, porque el Nodo puede seguir usando la anterior durante un corte.

Las operaciones diarias, las mediciones y los indicadores se escriben y consultan de maneras distintas. La Tabla \ref{tab:tecnologias-capa-6-alternativa-evaluada} separa esas tres cargas que el análisis queda fuera de los almacenes que reciben los hechos.

\begin{tabla}{I{0.088}I{0.305}I{0.294}I{0.313}}{Tecnologías de la capa 6 y alternativa evaluada.}{tab:tecnologias-capa-6-alternativa-evaluada}
  {\cab{Capa} & \cab{Tecnología y versión} & \cab{Uso en la solución} & \cab{Alternativa evaluada y criterio de decisión}}
  {Fuente: elaboración propia con información de PostgreSQL Global Development Group (s.f.), ClickHouse (s.f.) y OpenSearch Project (s.f.).}
 & PostgreSQL 18 & Guarda transacciones y auditoría de los módulos & Motor distinto por módulo: multiplica la operación \\
 & TimescaleDB 2.30 sobre PostgreSQL 18 & Guarda mediciones y telemetría ordenadas por tiempo & InfluxDB: añade otro motor que administrar \\
 & ClickHouse 26.8 LTS & Calcula indicadores sin cargar el registro operacional & Réplica de lectura de PostgreSQL: no ofrece almacenamiento columnar \\
 & OpenSearch 3 & Permite buscar entre módulos con resultados filtrados por rol & Búsqueda de texto de PostgreSQL: queda dentro de cada almacén \\
\multirow{-5}{=}{6. Datos} & Azure Blob Storage, con versionado e inmutabilidad por política de retención bloqueada & Conserva documentos y fotografías según su retención & Archivos dentro de la base de datos: mezclan tamaños y ciclos de vida distintos \\
\end{tabla}
La capa de datos usa un solo motor, PostgreSQL, para lo transaccional, las series de tiempo, la auditoría y la bandeja de salida, y sólo la analítica y la búsqueda agregan motores propios, porque su patrón de consulta no lo resuelve una base relacional.

\subsubseccion{Tecnologías de las capas 7 y 8: seguridad y observabilidad}
La pérdida del enlace separa las tareas de seguridad y supervisión. Las identidades y alertas se administran centralmente, mientras el Nodo mantiene la protección entre servicios y permite revisar su estado en faena. La Tabla \ref{tab:tecnologias-capas-7-8-alternativa-evaluada} reúne los controles centrales con la supervisión que queda disponible localmente.

\begin{tabla}{I{0.195}I{0.246}I{0.284}I{0.275}}{Tecnologías de las capas 7 y 8 y alternativa evaluada.}{tab:tecnologias-capas-7-8-alternativa-evaluada}
  {\cab{Capa} & \cab{Tecnología y versión} & \cab{Uso en la solución} & \cab{Alternativa evaluada y criterio de decisión}}
  {Fuente:elaboración propia con información de Microsoft (s.f.-e, s.f.-i), Envoy Project (s.f.), Prometheus (s.f.) y Grafana Labs (s.f.).}
 & Microsoft Entra ID y Microsoft Entra External ID & Entra ID identifica al personal y External ID, a contratistas y compradores & Directorio propio: perdería la federación corporativa \\
 & Azure Key Vault Premium & Custodia certificados y claves de cifrado mediante claves protegidas por HSM & Bóveda propia: exigiría administrar la custodia de claves \\
 & Malla de servicios basada en Istio administrada por AKS, y Envoy en la faena & Cifran la comunicación entre servicios en ambos entornos & SPIFFE con SPIRE propio: añade otra plataforma de identidad \\
 & Microsoft Sentinel & Reúne eventos y alertas de seguridad para su análisis & Correlación propia de eventos: exige un equipo de seguridad que el cliente no tiene \\
 & Microsoft Intune & Administra tabletas y consolas y distribuye sus configuraciones & Gestión manual por equipo: exige intervenir dispositivos dispersos \\
\multirow{-6}{=}{7. Seguridad} & Firma electrónica avanzada de un prestador acreditado & Firma actas, registros exigidos y reportes desde sus flujos & Firma con dispositivo físico por persona: saca la firma del proceso digital \\
 & OpenTelemetry y Azure Monitor & Reúnen métricas, registros y trazas con un identificador común & Agentes propietarios por producto: separan la información en varias consolas \\
\multirow{-2}{=}{8. Observabilidad} & Prometheus 3.13 LTS y Grafana 13.2 & Prometheus guarda métricas locales y Grafana las muestra durante un corte & Solo Azure Monitor: deja sin consulta local \\
\end{tabla}
La seguridad y la observabilidad usan servicios administrados en la nube y sólo tres componentes locales: Envoy, el colector local de OpenTelemetry y Prometheus con Grafana, justificados por la operación sin enlace. En la nube, la verificación mutua entre servicios la ejecuta el complemento de malla que administra AKS, cuyas revisiones se actualizan con la versión de Kubernetes (Microsoft, s.f.-e), en la faena la ejecuta Envoy con certificados emitidos por la autoridad certificadora custodiada en Key Vault. La gestión de dispositivos usa Microsoft Intune, que administra Windows 11 IoT Enterprise LTSC (Microsoft, s.f.-i), y la firma electrónica avanzada la entrega un prestador acreditado conforme a la Ley N° 19.799 (Ley N° 19.799, 2002).

\subsubseccion{Sistemas operativos y entorno de ejecución}
La vigencia de la plataforma depende también del sistema operativo de cada dominio de ejecución, como muestra la Tabla~\ref{tab:sistema-operativo-cada-dominio-ejecucion-fin}.

\begin{tabla}{I{0.210}I{0.179}I{0.299}I{0.313}}{Sistema operativo de cada dominio de ejecución y fin de soporte.}{tab:sistema-operativo-cada-dominio-ejecucion-fin}
  {\cab{Dominio de ejecución} & \cab{Sistema operativo} & \cab{Fin de soporte} & \cab{Justificación}}
  {Fuente: elaboración propia a partir de los ciclos de vida publicados por Canonical (s.f.) y Microsoft (s.f.-k, s.f.-l).}
Nodos de AKS en la nube & Azure Linux administrado por el servicio & Se actualiza con cada versión de Kubernetes & No queda a cargo del equipo de operación \\
Servidores del Nodo de Faena y de los gabinetes & Ubuntu Server 24.04 LTS con Ubuntu Pro & Estándar hasta mayo de 2029, seguridad hasta abril de 2034 & Cubre el contrato completo con una versión \\
Máquinas virtuales del conector del historiador & Windows Server 2025 & General hasta noviembre de 2029 está extendido hasta noviembre de 2034 & Kepware Server requiere Windows \\
Tabletas y equipos de la Consola de Portería & Windows 11 IoT Enterprise LTSC 2024 & General hasta octubre de 2029, extendido hasta octubre de 2034 & Canal de largo plazo sin cambio de versión durante el contrato \\
Botonera de Pala & Firmware firmado del fabricante & Soporte exigido en la especificación de compra hasta agosto de 2031 & No aloja lógica de la aplicación y se actualiza desde el Nodo de Faena \\
\end{tabla}
Los sistemas operativos bajo responsabilidad del adjudicatario cubren el contrato con una sola versión, y que el de la nube lo actualiza el propio servicio. Los servidores de la faena ejecutan K3s sobre Ubuntu, y el conector del historiador, Kepware Server, sobre Windows Server 2025 (Canonical, s.f.; Microsoft, s.f.-k, s.f.-l). Su emplazamiento en la zona desmilitarizada industrial, separado de los servidores del Nodo de Faena conforme a IEC 62443, se detalla en la arquitectura física.

\subsubseccion{Vigencia y plan de actualización}
Ninguna versión de Angular, Spring Boot, Kubernetes, ClickHouse, Envoy o .NET cubre por sí sola los 56 meses del contrato. Cada componente declara su versión, el fin de soporte publicado por el fabricante y la regla con que se actualiza, como muestra la Tabla~\ref{tab:version-fin-soporte-plan-actualizacion-componentes}.

\begin{tabla}{I{0.328}I{0.240}I{0.432}}{Versión, fin de soporte y plan de actualización de los componentes de software.}{tab:version-fin-soporte-plan-actualizacion-componentes}
  {\cab{Componente} & \cab{Fin de soporte} & \cab{Plan durante el contrato}}
  {Fuente: elaboración propia a partir de las políticas de soporte publicadas por cada fabricante.}
Angular 22 & Noviembre de 2027 & Cada versión mayor, antes del fin de soporte de la anterior \\
Java 25 LTS (Eclipse Temurin) & Septiembre de 2031 & Parches trimestrales que cubre el contrato \\
Spring Boot 4.1 & Julio de 2027 & Cada versión menor, dentro de seis meses \\
.NET 10 LTS del contenedor de terreno & Noviembre de 2028 & Siguiente LTS, antes del vencimiento \\
Kubernetes 1.36 (AKS y K3s) & 30 de junio de 2027, con soporte estándar. & Versión soportada, validada antes en preproducción \\
PostgreSQL 18 y TimescaleDB 2.30 & Noviembre de 2030 & Versión mayor en la detención de septiembre de 2030 y TimescaleDB la acompaña \\
ClickHouse 26.8 LTS & Agosto de 2027 & Cada LTS semestral \\
Envoy 1.39 & Julio de 2027 & Al menos una versión al año \\
Prometheus 3.13 LTS y Grafana 13.2 & Julio de 2027 y mayo de 2027 & Cada LTS de Prometheus y menor vigente de Grafana \\
OpenSearch 3 & Al menos un año tras la versión 4 & Migración a la 4 en ese plazo \\
SQLCipher 4 & Rama vigente, sin fecha de término & Con cada versión de la App SITCO Terreno \\
Kepware Server 7.1 & Junio de 2028 & Cada versión, antes de que venza la anterior \\
Servicios de Azure y Microsoft Intune & Microsoft avisa con al menos 12 meses de anticipación antes de retirar un servicio. & Revisión trimestral de los avisos \\
springdoc-openapi 3.1, Springwolf 2.8, Pact JVM 4.7 y Pact Broker 2.121 & Sin fecha publicada & Con cada menor de Spring Boot (módulo spring 7) y Pact Broker, cada trimestre \\
\end{tabla}
La tabla muestra, con las fechas que publica cada fabricante, que sólo Java 25 LTS cubre el contrato con una versión, los demás componentes cambian de versión dentro de los 56 meses con una regla fija, y los servicios de Azure e Intune se controlan con la revisión trimestral de sus avisos de retiro. Los cuatro componentes de contratos son de código abierto y no publican fecha de término, su reemplazo no cambia los contratos aprobados. Pact JVM, springdoc-openapi y Springwolf no mencionan Java 25 en su documentación, aunque Spring Boot 4.1 lo admite. Los tres se validan juntos en un piloto con Java 25 antes de la puesta en marcha. El Pact Broker prueba en su integración continua las versiones vigentes de PostgreSQL.

\subseccion{Estilo arquitectónico y decisiones de diseño}
El estilo de la arquitectura se deriva de tres condiciones del caso. La faena opera las 24 horas con dos turnos de 12 horas, el enlace de fibra con Antofagasta depende de un único proveedor y registra entre tres y seis cortes al año, de cuatro a veinte horas cada uno, y el área de tecnologías de información del cliente tiene once personas, cuatro de ellas en faena. De esas condiciones se obtienen un principio de ubicación de las funciones, un estilo de descomposición y un registro de decisiones fechadas.

\subsubseccion{Criterio de continuidad en la faena y consolidación en la nube}
Cada función se emplaza según su criticidad operacional, el origen de sus datos y su destinatario. Las funciones que sostienen la operación y reciben datos generados en la faena se ejecutan también en el Nodo de Faena; las destinadas a consolidación, usuarios externos o sistemas corporativos se ejecutan en la nube.

Este criterio mantiene instancias locales de Operación Mina, Trazabilidad y Embarque, Conciliación Metalúrgica, Acreditación y Acceso, y Seguridad y Salud Ocupacional. Los demás módulos operan en la nube, mientras Administración y Servicios Comunes replica al Nodo las identidades, parámetros y servicios necesarios durante un corte. Así, el Nodo ejecuta sólo las funciones críticas y no una réplica completa de la plataforma.

\subsubseccion{Servicios modulares frente a microservicios de grano fino}
Se compararon tres estilos: monolito modular, nueve servicios modulares uno por módulo y microservicios de grano fino. La evaluación considera operabilidad, continuidad sin enlace, despliegue independiente, escalamiento y evolución por etapas. Los pesos provienen de las condiciones del caso y la valoración utiliza una escala de 1, muy desfavorable, a 5, muy favorable. La Tabla \ref{tab:matriz-ponderada-estilos-descomposicion} presenta la matriz ponderada de los tres estilos.

\begin{tabla}{I{0.475}I{0.133}I{0.158}I{0.233}}{Matriz ponderada de estilos de descomposición.}{tab:matriz-ponderada-estilos-descomposicion}
  {\cab{Criterio y peso} & \cab{Monolito modular} & \cab{Servicios modulares} & \cab{Microservicios de grano fino}}
  {Fuente: elaboración propia.}
Operabilidad por un área de TI de 11 personas (30 \%) & 5 & 4 & 1 \\
Ejecución de cinco módulos en la faena sin enlace (25 \%) & 2 & 5 & 3 \\
Despliegue independiente de lo crítico (20 \%) & 1 & 4 & 5 \\
Escalamiento en el cambio de turno (15 \%) & 2 & 4 & 5 \\
Incorporación de la Etapa 2 sin rehacer la Etapa 1 (10 \%) & 3 & 5 & 4 \\
Puntaje ponderado & 2,80 & 4,35 & 3,20 \\
\end{tabla}
Los servicios modulares obtienen el mayor puntaje ponderado, con 4,35 frente a 2,80 del monolito modular y 3,20 de los microservicios de grano fino. El monolito simplifica la operación, pero limita el despliegue independiente y obliga a replicar más funciones en la faena. Los microservicios ofrecen mayor granularidad de escalamiento, pero agregan una complejidad operacional que no se justifica para la carga prevista.

Se adopta, por tanto, un servicio desplegable por cada módulo del Esquema de Solución, con almacén y ciclo de versión propios e integración mediante eventos e interfaces versionadas. Cada servicio se organiza internamente por agregados y puede dividirse posteriormente sin modificar sus contratos externos.

\subsubseccion{Registro de decisiones de arquitectura}
Las decisiones que condicionan la arquitectura se documentan como registros de decisión arquitectónica, con contexto, decisión, alternativas evaluadas, fundamento, consecuencias y referencia. Cada registro tiene fecha y estado, y un cambio posterior se documenta como un registro nuevo que reemplaza al anterior sin borrarlo. Las ocho decisiones vigentes de la arquitectura lógica se presentan en las tablas siguientes; las dos de la recuperación, ADR-09 y ADR-10, se registran en la sección del data center secundario. La Tabla \ref{tab:registro-decision-adr-01-servicios-modulares} registra la primera decisión, los servicios modulares.

\begin{tabla}{I{0.196}I{0.804}}{Registro de decisión ADR-01: servicios modulares.}{tab:registro-decision-adr-01-servicios-modulares}
  {\cab{Campo del registro} & \cab{ADR-01 · Servicios modulares}}
  {Fuente: elaboración propia, con una matriz de decisión ponderada por criterio.}
Contexto & Nueve módulos con distinto calendario de entrega, cinco de ellos también en la faena, y un área de TI de once personas que operará la solución. \\
Decisión & Construir nueve servicios modulares, uno por módulo, que se despliegan por separado. \\
Alternativas evaluadas & (a) Monolito modular de despliegue único, (b) microservicios de grano fino, uno por agregado raíz. \\
Fundamento & Obtiene 4,35 de 5 en la matriz ponderada, frente a 2,80 del monolito y 3,20 de los microservicios. \\
Consecuencias & Exige contratos versionados entre módulos, y los cinco módulos que corren también en la faena deben desplegarse en la nube y en el Nodo con versiones compatibles. \\
Referencia & RES-10, RNF-OPE-001, RP-006 y RNF-DIS-008, Evans (2003). \\
Fecha y estado & 5 de octubre de 2026 y aceptada \\
\end{tabla}
La decisión se sostiene en la operación posterior al contrato y en la ejecución en la faena, no en la carga, que no exige escalar función por función. La Tabla \ref{tab:registro-decision-adr-02-nodo-faena} registra la decisión de ejecutar cinco módulos en el Nodo de Faena.

\begin{tabla}{I{0.196}I{0.804}}{Registro de decisión ADR-02: Nodo de Faena con cinco módulos.}{tab:registro-decision-adr-02-nodo-faena}
  {\cab{Campo del registro} & \cab{ADR-02 · Nodo de Faena con cinco módulos}}
  {Fuente: elaboración propia, con una matriz de decisión ponderada por criterio.}
Contexto & El enlace con Antofagasta registra tres a seis cortes al año, de cuatro a veinte horas cada uno, y la faena debe operar al menos 48 horas sin enlace. \\
Decisión & Ejecutar cinco módulos en el Nodo de Faena, con un solo sitio escritor por agregado: los hechos de faena se escriben en el Nodo y los maestros y parámetros, en la nube. \\
Alternativas evaluadas & (a) Nodo que sólo almacena y reenvía, (b) réplica completa con escritura activa en ambos sitios. \\
Fundamento & La opción (a) no valida la portería ni calcula el balance del turno sin enlace, y la (b) permite que un mismo dato se escriba en los dos sitios durante un corte y quede con dos versiones. \\
Consecuencias & Los conflictos se resuelven con reglas fijas por tipo de dato y cada resolución queda en bitácora. \\
Referencia & RES-02, RNF-DIS-001, RNF-DIS-003 y RF-CNX-002. \\
Fecha y estado & 5 de octubre de 2026 y aceptada \\
\end{tabla}
Se puede concluir de la tabla que el sitio escritor único es lo que permite operar sin enlace sin crear dos versiones de un mismo dato. La Tabla \ref{tab:registro-decision-adr-03-integracion-eventos} registra la integración por eventos.

\begin{tabla}{I{0.236}I{0.764}}{Registro de decisión ADR-03: integración por eventos.}{tab:registro-decision-adr-03-integracion-eventos}
  {\cab{Campo del registro} & \cab{ADR-03 · Integración por eventos}}
  {Fuente: elaboración propia, con una matriz de decisión ponderada por criterio.}
Contexto & Un mismo hecho alimenta a varios módulos y debe llegar sin pérdida ni duplicados aunque ocurra durante un corte. \\
Decisión & Transportar los hechos por un bus de eventos con entrega al menos una vez, orden por agregado y deduplicación en el consumidor, alimentado desde la faena por una bandeja de salida transaccional. \\
Alternativas evaluadas & (a) Llamadas síncronas punto a punto, (b) bus de particiones con interfaz Kafka. \\
Fundamento & La opción (a) falla durante un corte y la (b) no entrega cola de fallidos por suscripción ni deduplicación nativa. \\
Consecuencias & Todo evento lleva un identificador generado en origen y todo consumidor es idempotente. \\
Referencia & RNF-TRV-02.06 y RNF-TRV-02.07, Microsoft (s.f.-g), Richardson (2018). \\
Fecha y estado & 5 de octubre de 2026 y aceptada \\
\end{tabla}
A partir de lo anterior, la integración por eventos convierte el corte del enlace en un retraso de entrega y no en una pérdida de datos. La Tabla \ref{tab:registro-decision-adr-04-conector-sistema} registra el uso de un conector por sistema existente.

\begin{tabla}{I{0.196}I{0.804}}{Registro de decisión ADR-04: un conector por sistema existente.}{tab:registro-decision-adr-04-conector-sistema}
  {\cab{Campo del registro} & \cab{ADR-04 · Un conector por sistema existente}}
  {Fuente: elaboración propia, con una matriz de decisión ponderada por criterio.}
Contexto & Nueve de los diez sistemas existentes se mantienen (las planillas sólo se migran), sus interfaces no están documentadas, y SAP y el despacho alimentan a varios módulos. \\
Decisión & Conectar cada sistema mediante un conector propio, con capa anticorrupción, dentro del Integrador de Sistemas. \\
Alternativas evaluadas & (a) Conexión directa de cada módulo con cada sistema, (b) plataforma comercial de integración con orquestación central. \\
Fundamento & Con un conector por sistema un cambio externo se resuelve en un solo lugar, y la opción (b) agrega una plataforma que exige especialistas. \\
Consecuencias & El levantamiento, cerrado antes del hito H2 del mes 4 (SUP-16), define la interfaz real de cada conector. \\
Referencia & SUP-16, RNF-INT-002, RNF-INT-003, RNF-INT-004 y RNF-TRV-02.08. \\
Fecha y estado & 5 de octubre de 2026 y aceptada \\
\end{tabla}
Se deduce que el conector aislado protege a los módulos de la falta de inventario técnico de la compañía: una interfaz distinta a la supuesta cambia un conector y no un módulo. La Tabla \ref{tab:registro-decision-adr-05-lectura-historiador} registra la lectura del historiador por diodo de datos.

\begin{tabla}{I{0.196}I{0.804}}{Registro de decisión ADR-05: lectura del historiador por diodo de datos.}{tab:registro-decision-adr-05-lectura-historiador}
  {\cab{Campo del registro} & \cab{ADR-05 · Lectura del historiador por diodo de datos}}
  {Fuente: elaboración propia, con una matriz de decisión ponderada por criterio.}
Contexto & La red de control de planta es de sólo lectura y el historiador tiene 14.500 etiquetas. \\
Decisión & Leer, con un conector OPC UA restringido a Read y Subscribe, un servidor espejo en la zona desmilitarizada industrial, alimentado por el historiador a través de un diodo de datos. \\
Alternativas evaluadas & (a) Cliente OPC UA directo al historiador detrás de un cortafuegos, (b) exportación programada de archivos. \\
Fundamento & El diodo impide físicamente el flujo hacia la red de control, lo que una regla de cortafuegos no garantiza, y la opción (b) no entrega el tonelaje dentro del turno. \\
Consecuencias & Se incorporan continuamente 5.000 señales de alimentación, flujos y detenciones, con una muestra cada cinco segundos; las demás etiquetas se consultan bajo demanda desde el espejo industrial. \\
Referencia & RES-03, RNF-INT-001, RNF-SEG-004 y SUP-19, International Society of Automation (2010), OPC Foundation (s.f.). \\
Fecha y estado & 5 de octubre de 2026 y aceptada \\
\end{tabla}
La tabla muestra que la barrera contra la escritura en la red de control no depende de una configuración que alguien pueda revertir. La Tabla \ref{tab:registro-decision-adr-06-canales-terreno} registra la decisión sobre los canales de terreno.

\begin{tabla}{I{0.196}I{0.804}}{Registro de decisión ADR-06: canales de terreno.}{tab:registro-decision-adr-06-canales-terreno}
  {\cab{Campo del registro} & \cab{ADR-06 · Canales de terreno}}
  {Fuente: elaboración propia, con una matriz de decisión ponderada por criterio.}
Contexto & Los dispositivos instalados antes en la faena se abandonaron porque no funcionaban con guantes ni se leían al sol, y el operador de pala no puede sumar pasos al ciclo. \\
Decisión & Botonera de Pala sin pantalla en la cabina y App SITCO Terreno híbrida, con núcleo web y contenedor nativo Windows, para la supervisión y las perforadoras sin instrumentación. \\
Alternativas evaluadas & (a) Tableta táctil en la cabina de la pala, (b) aplicación web progresiva, (c) aplicación nativa completa. \\
Fundamento & La botonera reemplaza una llamada por radio por un toque, y el contenedor entrega almacenamiento persistente, lectura de la credencial de proximidad y claves en hardware que el navegador no garantiza. \\
Consecuencias & Windows 11 IoT Enterprise pasa a ser condición de compra de las tabletas, y el RNF-POR-004 se atiende certificando la App SITCO Terreno en la versión LTSC vigente de Windows IoT Enterprise y en las dos versiones LTSC anteriores, con las actualizaciones mensuales de seguridad del canal LTSC como política de actualización (RT-17.03). \\
Referencia & RES-08, RNF-USA-001, RNF-USA-008, RNF-DES-016, SUP-09 y SUP-17, Microsoft (s.f.-c). \\
Fecha y estado & 5 de octubre de 2026 y aceptada \\
\end{tabla}
De lo anterior, se observa que la decisión separa dos usuarios con exigencias opuestas: el operador que no puede mirar una pantalla y el supervisor que necesita fotografía, firma y 48 horas sin señal. La Tabla \ref{tab:registro-decision-adr-07-identidad-poblacion} registra la identidad por población.

\begin{tabla}{I{0.196}I{0.804}}{Registro de decisión ADR-07: identidad por población.}{tab:registro-decision-adr-07-identidad-poblacion}
  {\cab{Campo del registro} & \cab{ADR-07 · Identidad por población}}
  {Fuente: elaboración propia, con una matriz de decisión ponderada por criterio.}
Contexto & Usan la solución 1.780 personas propias, hasta 6.700 contratistas durante una detención, unas noventa empresas, los compradores, el operador portuario y los laboratorios externos. \\
Decisión & Separar tres planos de identidad: directorio corporativo federado para el personal propio, identidad externa para empresas y compradores, y credencial con número de identificación personal (PIN) en terreno; en la botonera, el segundo factor es el contraste de la credencial con el operador que el despacho asigna a la pala, y el toque no abre sesión. \\
Alternativas evaluadas & (a) Todas las identidades en el directorio corporativo, (b) biometría obligatoria en todo acceso. \\
Fundamento & La opción (a) obliga al área de TI a administrar miles de identidades temporales, y la (b) depende de un acuerdo con el sindicato. \\
Consecuencias & Cada empresa contratista administra a sus trabajadores desde el Portal de Contratistas, y la botonera no usa PIN: su segundo factor es el contraste con el despacho, y el toque sólo firma la autoría de la declaración. \\
Referencia & RES-09, RNF-SEG-003, SUP-04, SUP-08 y SUP-13. \\
Fecha y estado & 5 de octubre de 2026, aceptada \\
\end{tabla}
Se puede observar que el doble factor se cumple en las tres poblaciones, incluida la botonera, sin exigir contraseñas alfanuméricas en terreno. En la botonera, el segundo factor es distinto: el toque de credencial no abre una sesión ni da acceso a datos, sólo firma la autoría de la declaración, y esa autoría se contrasta con el operador que el despacho tiene asignado a la pala, que actúa como segundo factor. La Tabla \ref{tab:registro-decision-adr-08-almacen-propio} registra el almacén propio por módulo y el modelo analítico separado.

\begin{tabla}{I{0.196}I{0.804}}{Registro de decisión ADR-08: almacén propio por módulo y modelo analítico separado.}{tab:registro-decision-adr-08-almacen-propio}
  {\cab{Campo del registro} & \cab{ADR-08 · Almacén propio por módulo y modelo analítico separado}}
  {Fuente: elaboración propia, con una matriz de decisión ponderada por criterio.}
Contexto & Las consultas de gestión y del CIO no pueden degradar el registro de terreno ni la portería. \\
Decisión & Dar a cada módulo su propio almacén transaccional y construir los indicadores sobre un modelo analítico alimentado por eventos. \\
Alternativas evaluadas & (a) Base de datos compartida por todos los módulos, (b) réplica de lectura de la base transaccional. \\
Fundamento & Con la opción (a) una consulta pesada compite con el registro, y la (b) no entrega el modelo dimensional ni la búsqueda entre módulos. \\
Consecuencias & Los indicadores tienen una latencia de hasta 15 minutos en turno y de 4 horas en gestión, y el cumplimiento del plan se publica al inicio del turno siguiente (RNF-DES-008). \\
Referencia & RNF-DES-010 y RNF-DES-017. \\
Fecha y estado & 5 de octubre de 2026 y aceptada \\
\end{tabla}
Del ADR-08 se concluye que el aislamiento entre el registro y la consulta es una decisión de datos y no sólo de capacidad de cómputo.

Cinco decisiones, ADR-02, ADR-03, ADR-05, ADR-06 y ADR-07, responden a condiciones de la faena que ninguna tecnología elimina: el enlace frágil, la red de control intocable, los guantes y el sol, y la mayoría contratista. Las otras tres ordenan la solución para que el cliente pueda operar y para que ninguna carga degrade el registro de terreno.

\subsubseccion{Patrones y principios de diseño}
La arquitectura se organiza con las vistas de ISO/IEC/IEEE 42010, orientadas a preocupaciones específicas del cliente (International Organization for Standardization, 2022): la vista lógica se desarrolla en 4.1.3 a 4.1.6. Los datos se describen en 4.1.4.6 y el Subdocumento 5. La integración, la seguridad, la operación sin enlace y el despliegue y la capacidad se describen en 4.1.7, 4.1.8, 4.1.9 y 4.2, respectivamente. La solución se apoya en ISA-95 para definir el intercambio entre SAP, SITCO y control de proceso (International Society of Automation, 2010), en confianza cero para seguridad (Rose et al., 2020), y en diseño guiado por el dominio para delimitar los módulos (Evans, 2003).

Los patrones de diseño resuelven los riesgos recurrentes: bandeja de salida transaccional en el Nodo de Faena y gabinetes, para guardar el hecho y su evento en una sola transacción (4.1.7.3; Richardson, 2018) consumidores y claves idempotentes, para evitar efectos duplicados ante reenvíos (4.1.7.2) y capas anticorrupción en cada conector del Integrador de Sistemas, para aislar los modelos de SAP, despacho e historiador (4.1.7.1; Evans, 2003). Cada conector aplica tiempos de espera, reintentos con retroceso exponencial, cortacircuito y mamparo de aislamiento, de modo que la falla de un sistema existente no se propague (4.1.7.1; Nygard, 2018).

Los módulos intercambian eventos de dominio mediante publicación y suscripción, con orden por agregado y cola de fallidos (4.1.7.3). Cada módulo conserva su propio almacén de las vistas integradas de Equipos y Mantenimiento, e Indicadores y Centro Integrado, son vistas materializadas alimentadas por eventos (ADR-08, 4.1.6.2 y 4.1.6.4). Dos puertas de enlace separan la exposición para portales e integraciones de los canales operacionales de terreno (4.1.4.3), y las colas nivelan la carga para absorber saturaciones sin rechazar registros de terreno (4.2.4).

Cada servicio modular se implementa con puertos y adaptadores y aplica SOLID (Martin, 2017). Cada servicio tiene una única responsabilidad, y su razón de cambio son las reglas de su proceso, mientras la identidad y la traducción de sistemas externos permanecen aisladas. Se amplía mediante módulos, conectores y parámetros versionados, sin modificar lo existente. Permite sustituir conectores que cumplen el mismo contrato de correlación y antigüedad del dato, incluso entre interfaces y exportaciones programadas (SUP-16). Además, expone sólo las interfaces necesarias para cada consumidor (Tabla \ref{tab:interfaces-sincronas-eventos-publicados-datos-consumidos}), y hace depender la lógica de negocio de contratos OpenAPI y AsyncAPI y de puertos abstractos de mensajería y almacenamiento. Así, un mismo módulo puede operar en nube mediante el bus de eventos o en el Nodo de Faena mediante la bandeja de salida local, cambiando únicamente el adaptador.

\subseccion{Vista general en capas}
Esta sección presenta las ocho capas en una vista general, la correspondencia de los componentes del Esquema de Solución con cada una, su ubicación en la faena, el Nodo de Faena y la nube, y los ambientes lógicos en que se despliegan.

\subsubseccion{Las ocho capas y la regla de dependencia entre ellas}
La Figura~\ref{fig:vista-general-sitco-aranda-ocho-capas} presenta la vista general. Los canales de terreno aparecen a la izquierda y los canales web a la derecha, porque recorren caminos distintos hasta los servicios de negocio. En ella, los portales externos entran por la entrada pública y el Portal Web SITCO, por el acceso de confianza cero.

\figura{15.62}{fig-4-1-vista-general-sitco-aranda-ocho}{Vista general de SITCO Aranda en ocho capas.}{fig:vista-general-sitco-aranda-ocho-capas}{Fuente: elaboración propia.}
Cada capa depende sólo de capas inferiores o de una transversal, y ninguna interfaz permite llegar a un almacén o a un sistema existente sin pasar por su responsable. Los canales llegan a los módulos a través de una puerta, los módulos llegan a los sistemas existentes a través de la capa 5 y cada almacén transaccional lo escribe sólo su módulo dueño. El único salto previsto es el de los canales de terreno, que no atraviesan el borde para que el registro en terreno no compita con el tráfico de los portales.

\subsubseccion{Correspondencia de los canales, módulos y plataformas con las capas}
El Esquema de Solución organiza el software en nueve módulos, siete de negocio y dos transversales, que las personas usan a través de seis canales y que se apoyan en dos plataformas sin funciones de negocio propias: el Integrador de Sistemas y el Nodo de Faena, ambos del componente software. En la vista lógica, el Nodo de Faena se ubica en la capa 5 por el software que ejecuta: los cinco módulos locales, los conectores con los sistemas de la faena y la copia de identidades. La arquitectura lógica conserva esos nombres y agrega servicios técnicos sin funciones de negocio para completar las capas que el esquema no detalla. Esos servicios técnicos llevan nombres fijos, que se usan igual en esta sección, en el modelo de amenazas y en la arquitectura física. La Tabla~\ref{tab:correspondencia-canales-modulos-plataformas-esquema-solucion} presenta la correspondencia.

\begin{tabla}{I{0.195}I{0.402}I{0.402}}{Correspondencia de los canales, módulos y plataformas del Esquema de Solución con las ocho capas.}{tab:correspondencia-canales-modulos-plataformas-esquema-solucion}
  {\cab{Capa} & \cab{Canales, módulos y plataformas del Esquema de Solución} & \cab{Servicios técnicos que la completan}}
  {Fuente: elaboración propia a partir del Esquema de Solución del Subdocumento 3.}
1. Presentación & App SITCO Terreno, Botonera de Pala, Portal Web SITCO, Portal de Contratistas, Portal para Compradores y Consola de Portería & No aplica \\
2. Borde y exposición & No aplica & Entrada pública y acceso de confianza cero \\
3. Puertas de enlace & No aplica & Puerta de exposición y puerta operacional \\
4. Servicios de negocio & Siete módulos de negocio y dos módulos transversales & Caché de sesión y de proceso \\
5. Integración y eventos & Integrador de Sistemas y Nodo de Faena & Bus de eventos \\
6. Datos & No aplica & Almacén transaccional de cada módulo, series de tiempo, almacén documental, copia inmutable de auditoría, modelo analítico, índice de búsqueda y repositorio de consulta histórica \\
7. Seguridad transversal & No aplica & Proveedor de identidad, bóveda de claves, malla de servicios, registro y detección de seguridad, y gestión de dispositivos \\
8. Observabilidad transversal & No aplica & Colectores, consola operativa única y alertas \\
\end{tabla}
En la tabla se aprecia que los seis canales, los nueve módulos y las dos plataformas del esquema quedan en las capas 1, 4 y 5 con el mismo nombre, y que las capas 2, 3, 6, 7 y 8 se completan con servicios técnicos que no contienen reglas de negocio. Las capas 4 y 5 también se completan con servicios técnicos: la caché de sesión y de proceso y el bus de eventos. Esta separación permite que la Etapa 2 agregue módulos y conectores sobre las mismas capas, sin modificar el borde, las puertas ni la seguridad construidos en la Etapa 1.

La tabla presenta, para cada módulo o plataforma, los requerimientos que esta sección desarrolla, su distribución por etapa está en el Subdocumento 3, sección 3.2.4.2. La Tabla \ref{tab:requerimientos-funcionales-modulo-plataforma-etapa} muestra cuántos requerimientos funcionales tiene cada módulo o plataforma por etapa.

\begin{tabla}{I{0.297}I{0.125}I{0.125}I{0.453}}{Requerimientos funcionales por módulo o plataforma y etapa.}{tab:requerimientos-funcionales-modulo-plataforma-etapa}
  {\cab{Módulo o plataforma} & \cab{RF de la Etapa 1} & \cab{RF de la Etapa 2} & \cab{Requerimientos desarrollados en la arquitectura lógica}}
  {Fuente: elaboración propia a partir de Overfore-Subdocumento2-Anexos y del Formulario T-12.}
Operación Mina & 20 & 0 & RF-TRZ-001, RF-TER-002, RF-TER-003, RF-TER-006, RF-TER-008 \\
Trazabilidad y Embarque & 16 & 1 & RF-TRZ-005, RF-TRZ-014, RF-TRZ-017, RF-TRZ-019 \\
Conciliación Metalúrgica & 15 & 0 & RF-CON-001 a RF-CON-006, RF-CON-015 \\
Equipos y Mantenimiento & 0 & 10 & RF-EQU-001, RF-EQU-005, RF-PLA-006 \\
Acreditación y Acceso & 16 & 2 & RF-ACR-009, RF-ACR-015, RF-ACR-016 \\
Seguridad y Salud Ocupacional & 0 & 7 & RF-CUM-005 \\
Cumplimiento Ambiental & 0 & 10 & RF-CUM-009, RF-CUM-014 \\
Indicadores y Centro Integrado & 1 & 4 & RF-ANA-001, RF-ANA-002, RF-CON-013, RF-CON-014 \\
Administración y Servicios Comunes & 19 & 0 & RF-ADM-005, RF-ADM-012, RF-DAT-001, RF-NOT-001 \\
Integrador de Sistemas y Nodo de Faena & 5 & 0 & RF-CNX-001 a RF-CNX-005 \\
Total & 92 & 34 & No aplica \\
\end{tabla}
Los 126 requerimientos funcionales quedan asignados a un módulo o plataforma con el mismo nombre que en el Subdocumento 3. El Anexo B de OVERFORE-Subdocumento4-Anexos lista los 126 códigos por módulo, de modo que la matriz se puede revisar completa sin recurrir al Formulario T-12. Los seis canales de uso no tienen requerimientos funcionales propios, porque sólo presentan lo que resuelve el módulo al que sirven con sus exigencias que son no funcionales, como la operación con guantes o la accesibilidad.

El catálogo asigna los cinco requerimientos de conectividad en conjunto al Integrador de Sistemas y al Nodo de Faena. La arquitectura los reparte según qué componente ejecuta la función. El Nodo de Faena es responsable de los cinco: sincroniza lo acumulado y aplica el orden y la deduplicación al volver el enlace que guarda la lista de habilitaciones con que el agente del gabinete de portería valida sin enlace, recibe los registros cifrados de los dispositivos y les entrega la fuente horaria común y la sesión local en el relevo, administra el firmware de las botoneras y coordina con Microsoft Intune el bloqueo remoto en 15 minutos o menos, y registra cada corte con su inicio, su término y el volumen acumulado. El Integrador de Sistemas participa sólo en dos: al volver el enlace, sus conectores de nube vacían la cola de envíos pendientes hacia SAP y la autoridad, y el conector del control de acceso, que corre en el gabinete de portería, entrega la lectura de la credencial y ejecuta la apertura del torniquete durante el corte.

\subsubseccion{Ubicación de las capas en la faena, el Nodo de Faena y la nube}
La Tabla~\ref{tab:ubicacion-cada-capa-faena-nodo-faena} ubica cada capa en las zonas de la vista general del Esquema de Solución (Subdocumento 3, Figura \refext{S3}{fig:vista-general-sitco-aranda}):

\begin{tabla}{I{0.195}I{0.268}I{0.256}I{0.281}}{Ubicación de cada capa en la faena, el Nodo de Faena y la nube.}{tab:ubicacion-cada-capa-faena-nodo-faena}
  {\cab{Capa} & \cab{Dispositivos y gabinetes de faena} & \cab{Nodo de Faena} & \cab{Nube}}
  {Fuente: elaboración propia a partir del Subdocumento 3, Figura \refext{S3}{fig:vista-general-sitco-aranda}.}
1. Presentación & App SITCO Terreno, Botonera de Pala y Consola de Portería & Instancia local del Portal Web SITCO para los cinco módulos & Portal Web SITCO, Portal de Contratistas y Portal para Compradores \\
2. Borde y exposición & No aplica & No aplica & Entrada pública y acceso de confianza cero \\
3. Puertas de enlace & Puerta operacional en cada gabinete & Puerta operacional & Puerta de exposición \\
4. Servicios de negocio & Validación de portería en el gabinete de portería & Cinco módulos & Nueve módulos \\
5. Integración y eventos & Agente de cada gabinete con su cola local & Bandeja de salida, conectores locales y copia de identidades & Bus de eventos y conectores de nube \\
6. Datos & Registro local cifrado de cada dispositivo & Almacenes de los cinco módulos & Todos los almacenes \\
7. Seguridad transversal & Claves del dispositivo en su módulo de plataforma segura (TPM) & Copia de políticas e identidades & Proveedor de identidad y bóveda de claves \\
8. Observabilidad transversal & Colector local & Colector local y tableros de faena & Consola operativa única \\
\end{tabla}
Las capas 1, 3, 4, 5, 6, 7 y 8 tienen componentes en la faena para sostener un camino completo desde el canal hasta el dato durante un corte. La capa 2 existe solo en la nube, porque la faena no publica servicios hacia Internet. La asignación componente por componente, incluidos los servicios de nube, equipos y dispositivos, se presenta en 4.2 y en el Anexo B Mapeo lógico y físico.

\subsubseccion{Ambientes lógicos y promoción entre ellos}
La solución opera Desarrollo, QA, Preproducción, Producción y Recuperación ante Desastres. Todos comparten la misma arquitectura lógica y contratos, se diferencian en escala, datos y permisos. Una misma imagen firmada avanza sin recompilar de Desarrollo a QA, Preproducción y Producción, con configuración y secretos separados. Los ambientes no productivos usan datos sintéticos o enmascarados. Recuperación no forma parte de la promoción: recibe la réplica de Producción y asume los servicios ante una contingencia. La ubicación, las pruebas y la cadena de entrega se detallan en 4.2.3.1 y en el Subdocumento 6.

\subseccion{Descripción de las capas}
Cada capa se describe por su responsabilidad, los componentes responsables de cumplirla, las interfaces que ofrece y consume, y lo que queda fuera de ella. Para seguir el hilo entre una capa y la siguiente se usa un mismo caso a lo largo de la sección: la declaración de material que hace el operador de pala al cargar un camión. Ese hecho nace en la capa 1, entra por la capa 3, lo registra Operación Mina en la capa 4, viaja por la capa 5 y queda en la capa 6, mientras las capas 7 y 8 lo protegen y lo siguen en todo el recorrido.

\subsubseccion{Capa 1: Presentación, con los seis canales de uso}
La capa de presentación reúne los seis canales por los que las personas usan la solución. Su responsabilidad es capturar el hecho o la consulta, mostrar el resultado y, en terreno, guardar el registro cuando no hay señal. Cada canal está a cargo de una población de usuarios, y los seis se agrupan en dos familias según el camino que siguen hasta los servicios de negocio, como muestra la Figura \ref{fig:canales-capa-presentacion-rutas-hacia-puertas}.

\begin{horizontal}
\figura{22.44}{fig-4-2-canales-capa-presentacion-rutas-hacia}{Canales de la capa de presentación y sus rutas hacia las puertas de enlace.}{fig:canales-capa-presentacion-rutas-hacia-puertas}{Fuente: elaboración propia.}
\end{horizontal}
La Figura~\ref{fig:canales-capa-presentacion-rutas-hacia-puertas} muestra que los tres canales de terreno llegan a la puerta operacional por la red privada de telefonía móvil LTE o por la fibra de portería, y que durante un corte la instancia local del Portal Web SITCO también llega a ella, los tres canales web pasan por el borde antes de llegar a la puerta de exposición. Los tres canales de terreno guardan registros sin enlace. La instancia local del Portal Web SITCO opera también sin enlace en el Nodo de Faena. El operador de pala usa un botón y el de camión sigue usando el despacho de flota, de modo que ningún operador de equipo recibe una pantalla nueva; las tabletas pertenecen a los puestos de supervisión, que registran en el lugar lo que hoy se anota en planillas al final del turno.

Los canales no contienen reglas de negocio ni acceden a un almacén: los de terreno llaman a la puerta operacional y los web a la puerta de exposición, siempre por interfaces versionadas. El registro local de los dispositivos de terreno es una cola de hechos firmados a la espera de sincronización, no una copia de la base de datos. La capa cumple WCAG 2.2 nivel AA en todas sus interfaces y presenta en español el 100 \% de sus mensajes.

\subsubseccion{Capa 2: Borde y exposición}
La capa de borde es la única puerta desde Internet hacia la solución. Existe sólo en la nube, porque la faena no publica servicios hacia Internet, y tiene dos entradas con responsables distintos. La Figura \ref{fig:entradas-capa-borde-exposicion} muestra ambas entradas, los controles que aplica cada una y el destino de cada tráfico.

\begin{horizontal}
\figura{22.44}{fig-4-3-entradas-capa-borde-exposicion}{Entradas de la capa de borde y exposición.}{fig:entradas-capa-borde-exposicion}{Fuente: elaboración propia.}
\end{horizontal}
Cada población que llega desde Internet tiene una sola entrada. La entrada pública atiende a contratistas, compradores e integraciones autorizadas, incluidos los laboratorios externos, y aplica en serie la red de distribución de contenidos, el cortafuegos de aplicaciones web con reglas gestionadas y propias, la protección contra denegación de servicio distribuida en las capas 3, 4 y 7 del modelo OSI y la terminación TLS 1.3. Entrega el tráfico al Portal de Contratistas, al Portal para Compradores y a la puerta de exposición, con lo que el 100 \% de los servicios alcanzables desde fuera se publica sólo a través de esta capa. La protección contra bots usa un reto progresivo compatible con WCAG 2.2 AA.

La figura marca también el tráfico que la capa no admite. Los canales de terreno llegan a la puerta operacional por la red local, y ningún viaje ni paso de portería atraviesa el borde público. Por eso la declaración del operador de pala no pasa por esta capa, y los tiempos de registro en terreno se mantienen aunque los portales operen a 1,5 veces su peak. La capa tampoco decide autorizaciones de negocio, que corresponden a la puerta de exposición y a los módulos.

\subsubseccion{Capa 3: Puertas de enlace de servicios}
En la capa de puertas convergen los dos caminos de la capa 1. Tiene dos puertas porque sus consumidores tienen perfiles incompatibles: los portales y las integraciones necesitan control por consumidor, y los dispositivos de terreno necesitan que ningún mensaje se pierda. La Figura~\ref{fig:consumidores-controles-puerta-exposicion-puerta-operacional} muestra qué consumidores atiende cada puerta y qué hace cada una con un mensaje que valida y con uno que no valida.

\begin{horizontal}
\csname ov@hoja\endcsname{330.2mm}{215.9mm}
\expandafter\def\csname ov@altofig\endcsname{0.88}
\figura{27.52}{fig-4-4-consumidores-controles-puerta-exposicion-puerta}{Consumidores y controles de la puerta de exposición y de la puerta operacional.}{fig:consumidores-controles-puerta-exposicion-puerta-operacional}{Fuente: elaboración propia.}
\end{horizontal}
Ambas puertas aplican la misma secuencia de controles sobre el 100 \% de las llamadas: autenticación, autorización, cuotas, límites de tasa, validación de esquema e inspección de carga útil. Se diferencian en lo que hacen después. La puerta de exposición, en la nube, entrega lo válido a los módulos de la nube y rechaza con error explícito lo que no valida. La puerta operacional, en el Nodo de Faena y en los cuatro gabinetes, entrega lo válido a los módulos del Nodo y a los agentes de gabinete, y guarda lo que no valida en la cola de fallidos con su carga íntegra, para revisarlo sin perder el hecho. En el caso de la declaración de material, la puerta operacional verifica el certificado de la Botonera de Pala y la credencial del operador, contrasta esa credencial con el operador que el despacho tiene asignado a la pala, valida el mensaje y lo entrega a Operación Mina en el Nodo. La Tabla \ref{tab:diferencias-entre-puerta-exposicion-puerta-operacional} resume las diferencias entre ambas puertas.

\begin{tabla}{I{0.241}I{0.322}I{0.437}}{Diferencias entre la puerta de exposición y la puerta operacional.}{tab:diferencias-entre-puerta-exposicion-puerta-operacional}
  {\cab{Atributo} & \cab{Puerta de exposición} & \cab{Puerta operacional}}
  {Fuente: elaboración propia con exigencias de Compañía Minera Altos de Aranda (2026b, cap. 15, p. 25)}
Dónde se ejecuta & Nube & Nodo de Faena y cuatro gabinetes \\
Consumidores & Portales, Portal Web SITCO e integraciones autorizadas & App SITCO Terreno, Botonera de Pala, Consola de Portería, instancia local del Portal Web SITCO y agentes de gabinete \\
Autenticación & OAuth 2.1 con credenciales de cliente o TLS mutuo, más la sesión de la persona & Certificado del dispositivo o puesto inventariado más la credencial de la persona \\
Mensaje que no valida & Rechazo con error explícito & Cola de fallidos con su carga íntegra \\
Prioridad & Cuota por consumidor & Prioridad del tráfico operacional \\
Alcance desde Internet & Sí, detrás del borde & No, sólo abre conexiones salientes \\
\end{tabla}
Analizando la tabla, se puede observar que la puerta de exposición prioriza controlar a cada consumidor y la operacional prioriza no perder ningún hecho. Como la puerta operacional no acepta conexiones entrantes desde Internet, sólo es alcanzable desde la red local de cada sitio: la red de la faena para el Nodo y sus gabinetes, y la red del acopio para el gabinete de puerto.

\subsubseccion{Capa 4: Servicios de negocio, con siete módulos de negocio y dos transversales}
\enlargethispage{2\baselineskip}La capa de servicios de negocio es la que decide. Contiene los siete módulos de negocio y los dos transversales del Esquema de Solución, cada uno como un servicio modular, y su responsabilidad es aplicar las reglas del negocio: registrar los hechos, validarlos, calcular balances e indicadores y decidir la habilitación de una persona. La Figura \ref{fig:modulos-capa-servicios-negocio-nube-nodo} muestra cómo se agrupan los módulos, cuáles tienen además una instancia en el Nodo de Faena y con qué capas se relacionan.

\begin{horizontal}
\csname ov@hoja\endcsname{330.2mm}{215.9mm}
\expandafter\def\csname ov@altofig\endcsname{0.88}
\figura{27.52}{fig-4-5-modulos-capa-servicios-negocio-nube}{Módulos de la capa de servicios de negocio en la nube y en el Nodo de Faena.}{fig:modulos-capa-servicios-negocio-nube-nodo}{Fuente: elaboración propia.}
\end{horizontal}
Los nueve módulos se reparten en cuatro grupos, y cada grupo responde por una parte del problema. La cadena del mineral enlaza tres módulos por eventos: Operación Mina publica el viaje, Trazabilidad y Embarque lo sigue hasta formar el lote y Conciliación Metalúrgica cuadra las cantidades de cada tramo. El grupo de activos contiene Equipos y Mantenimiento, que entra en la Etapa 2. En cumplimiento y personas, Acreditación y Acceso entrega a Operación Mina la habilitación de cada persona, y Seguridad y Salud Ocupacional y Cumplimiento Ambiental entran en la Etapa 2. Los transversales, Indicadores y Centro Integrado y Administración y Servicios Comunes, sirven a todos los demás.

En la faena, los cinco módulos definidos en 4.1.2.1 tienen una instancia en el Nodo de Faena. El gabinete de portería ejecuta la validación de portería y el gabinete de puerto registra los hechos del acopio. Cada agregado tiene un solo sitio escritor. Los hechos nacidos en la faena, como el viaje, la capa de stock, el lote de planta, el despacho a puerto, el paso de portería y el registro de seguridad, se escriben en el Nodo y suben a la nube, los hechos del acopio de puerto se escriben en la nube desde el gabinete de puerto, y los maestros y parámetros se escriben en la nube y bajan al Nodo. La declaración del operador de pala la registra así la instancia de Operación Mina del Nodo, que la asocia al viaje aunque no haya enlace.

Hacia arriba, los servicios ofrecen interfaces síncronas a las puertas de la capa 3. Hacia abajo, publican y consumen eventos en el bus de la capa 5, y cada uno lee y escribe sólo su almacén en la capa 6. Ningún servicio guarda estado en memoria: el estado de sesión y de proceso vive en una caché externa de alta disponibilidad, lo que permite agregar o retirar instancias sin perder transacciones. Los umbrales, plazos, tolerancias y catálogos son parámetros versionados de Administración y Servicios Comunes y no constantes de código, de modo que un período cerrado se recalcula con la regla vigente en ese período.

\subsubseccion{Capa 5: Integración y eventos, con el Integrador de Sistemas y el Nodo de Faena}
La capa de integración y eventos mueve el dato sin decidir sobre él. Sus responsables son el bus de eventos, el Integrador de Sistemas y el Nodo de Faena, que además tiene un agente en cada gabinete. Las Figuras \ref{fig:componentes-capa-integracion-eventos-nube-faena} y \ref{fig:componentes-capa-integracion-eventos-nube-faena-b} muestran dónde corre cada componente y qué sistema atiende cada conector.

\begin{horizontal}
\expandafter\def\csname ov@altofig\endcsname{0.80}
\begingroup
\def\thefigure{4.6a}
\def\theHfigure{4.6a}
\figura{22.44}{s4-revision-6a}{Integración y eventos en la nube.}{fig:componentes-capa-integracion-eventos-nube-faena}{Fuente: elaboración propia.}
\endgroup
El bus distribuye los eventos y cada conector conserva el límite del sistema que integra.

\clearpage
\addtocounter{figure}{-1}
\begingroup
\def\thefigure{4.6b}
\def\theHfigure{4.6b}
\figura{22.44}{s4-revision-6b}{Integración y eventos en la faena.}{fig:componentes-capa-integracion-eventos-nube-faena-b}{Fuente: elaboración propia.}
\endgroup
El Nodo conserva los hechos durante el corte y los entrega a la nube cuando vuelve el enlace.

\end{horizontal}
El bus transporta eventos entre módulos con entrega al menos una vez, orden por agregado y cola de fallidos. El Integrador aísla cada sistema existente mediante un conector y una capa anticorrupción. El Nodo de Faena ejecuta los conectores locales y conserva los eventos en una bandeja de salida durante los cortes. El detalle de conectores, contratos y resiliencia se presenta en 4.1.7.

El historiador tiene una ruta propia. La red de control de planta entrega sus datos a través de un diodo de datos, que sólo permite la salida, a un servidor espejo en la zona desmilitarizada industrial, y el conector lee ese espejo con OPC UA restringido a Read y Subscribe (ADR-05). Así, ningún componente de SITCO puede escribir en la red de control.

El Nodo de Faena ejecuta los cinco módulos locales y cumple las cinco responsabilidades de conectividad descritas anteriormente. Intercambia lecturas y órdenes con los conectores locales y recibe de ellos los eventos. Hacia la nube, la bandeja de salida publica en el bus cuando hay enlace, y la nube le envía los maestros, el plan y las identidades. En cada gabinete de borde corre un agente del Nodo con su cola local: en portería, la lista de personas habilitadas; en planta, el espejo del historiador y el caudalímetro de alimentación de agua; en relaves, la instrumentación ambiental y geotécnica, y en puerto, el registro del acopio y del embarque. Los caudalímetros del campamento y del punto de carga de los camiones aljibe se leen por la red de la faena desde el Nodo de Faena.

En el caso de la declaración de material, el viaje que registra Operación Mina en el Nodo queda en la bandeja de salida y se publica en el bus cuando hay enlace; si el enlace está cortado, espera allí sin que la pala se detenga. La capa ofrece a los módulos la publicación y suscripción de eventos y el acceso a los sistemas existentes, y consume las interfaces de esos sistemas a través de su capa anticorrupción. Excluye las decisiones de negocio: el Integrador traduce y el Nodo ejecuta los módulos, pero ninguno de los dos contiene reglas propias.

\subsubseccion{Capa 6: Datos}
En la capa de datos termina el recorrido del hecho. La capa guarda los datos en almacenes lógicos separados por un motivo verificable, y cada almacén tiene un responsable. La Figura \ref{fig:almacenes-capa-datos-unidades-trazabilidad} muestra los almacenes de cada módulo, los compartidos y los de la faena, junto con las unidades de trazabilidad que conectan los datos del ciclo del mineral.

\begin{horizontal}
\figura{22.44}{fig-4-7-almacenes-capa-datos-unidades-trazabilidad}{Almacenes de la capa de datos y unidades de trazabilidad.}{fig:almacenes-capa-datos-unidades-trazabilidad}{Fuente: elaboración propia.}
\end{horizontal}
Cada módulo es dueño exclusivo de su almacén y comparte información mediante eventos. Los almacenes comunes atienden series de tiempo, documentos, analítica, búsqueda y consulta histórica. El Nodo y los dispositivos mantienen copias locales cifradas para operar sin enlace. De la figura se concluye que la trazabilidad no depende de un almacén central: cada unidad queda en el almacén de su módulo y la cadena se reconstruye con los eventos, por lo que la faena sigue registrando aunque no haya enlace.

La declaración del operador de pala queda guardada en el almacén transaccional de Operación Mina como parte del viaje, la primera unidad de trazabilidad del ciclo del mineral. La cadena de unidades se muestra en la Figura \ref{fig:almacenes-capa-datos-unidades-trazabilidad} y su modelo se desarrolla en el Subdocumento 5. La capa ofrece a los servicios transacciones con garantías de atomicidad y aislamiento, cifrado en reposo y respaldo, y excluye cualquier consulta de un módulo al almacén de otro. La selección de motores, la consistencia por dominio, la retención y la migración de los datos se desarrollan en el Subdocumento 5.

\subsubseccion{Capa 7: Seguridad transversal}
Las capas 7 y 8 no forman un paso del recorrido, porque se aplican sobre las capas 1 a 6. La capa de seguridad sigue el modelo de confianza cero, y la red de la faena no concede confianza por sí misma. Las Figuras \ref{fig:componentes-capa-seguridad-transversal-excepcion-durante} y \ref{fig:componentes-capa-seguridad-transversal-excepcion-durante-b} muestran los componentes responsables en la nube y en la faena, la población que atiende cada identidad y la forma en que opera la capa durante un corte.

\begin{horizontal}
\expandafter\def\csname ov@altofig\endcsname{0.80}
\begingroup
\def\thefigure{4.8a}
\def\theHfigure{4.8a}
\figura{22.44}{s4-revision-8a}{Seguridad transversal en la nube.}{fig:componentes-capa-seguridad-transversal-excepcion-durante}{Fuente: elaboración propia.}
\endgroup
Las identidades, las claves y la supervisión central se administran desde la nube.

\clearpage
\addtocounter{figure}{-1}
\begingroup
\def\thefigure{4.8b}
\def\theHfigure{4.8b}
\figura{22.44}{s4-revision-8b}{Seguridad local y corte de enlace.}{fig:componentes-capa-seguridad-transversal-excepcion-durante-b}{Fuente: elaboración propia.}
\endgroup
La copia controlada de identidades permite continuar la operación y acumular los eventos locales.

\end{horizontal}
La seguridad aplica confianza cero, identidad por población, autorización por rol y atributos, claves protegidas y TLS mutuo. Durante un corte, el Nodo usa una copia temporal de políticas e identidades. De las figuras se concluye que un corte cambia dónde se toma la decisión de acceso, no la regla: el Nodo aplica las mismas políticas con la copia temporal y acumula los eventos de acceso para revisarlos en la nube al volver el enlace. Los controles y amenazas se desarrollan en 4.1.8.

\subsubseccion{Capa 8: Observabilidad transversal}
La capa de observabilidad sigue el hecho por todas las capas. Reúne métricas, registros y trazas unidos por un identificador de transacción en el 100 \% de las operaciones, instrumentados con OpenTelemetry en la nube, el Nodo de Faena y los gabinetes. La Figura \ref{fig:flujo-telemetria-capa-observabilidad-transversal} muestra cómo viaja la telemetría desde la faena hasta la consola y qué alertas produce.

\begin{horizontal}
\figura{22.44}{fig-4-9-flujo-telemetria-capa-observabilidad-transversal}{Flujo de la telemetría en la capa de observabilidad transversal.}{fig:flujo-telemetria-capa-observabilidad-transversal}{Fuente: elaboración propia.}
\end{horizontal}
Los componentes emiten métricas, registros y trazas correlacionados. Los colectores locales conservan telemetría durante los cortes y Azure Monitor consolida la operación al restablecerse el enlace. La consola alerta por estado, integración y retraso de sincronización. De la figura se concluye que un corte no deja puntos ciegos: lo ocurrido en la faena se consolida después, y el retraso de sincronización es una alerta propia, de modo que la operación sabe cuánto falta por sincronizar.

\subseccion{Canales de terreno}
Los canales de terreno condicionan la arquitectura más que cualquier otro componente, porque en ellos nace el dato. Esta sección desarrolla la decisión ADR-06: las exigencias que la acotan, las alternativas comparadas y la solución adoptada.

Para la supervisión se compararon tres naturalezas de aplicación sobre la misma tableta rugerizada, como muestra la Tabla~\ref{tab:comparacion-aplicacion-web-progresiva-nativa-hibrida}.

\begin{tabla}{I{0.191}I{0.349}I{0.148}I{0.313}}{Comparación de aplicación web progresiva, nativa e híbrida para la supervisión en terreno.}{tab:comparacion-aplicacion-web-progresiva-nativa-hibrida}
  {\cab{Criterio} & \cab{Web progresiva} & \cab{Nativa} & \cab{Híbrida}}
  {Fuente: elaboración propia.}
Registro durante 48 horas sin señal & Insuficiente: el sistema puede desalojar el almacenamiento del navegador & Suficiente & Suficiente: el contenedor gestiona un almacenamiento persistente \\
Lectura de la credencial de proximidad & Insuficiente: el navegador no accede al lector del equipo & Suficiente & Suficiente: el contenedor expone el lector al núcleo web \\
Claves de firma en hardware & Insuficiente & Suficiente & Suficiente, en el módulo TPM \\
Bases de código que mantener & Una & Dos, web y nativa & Una, más una capa nativa acotada \\
\end{tabla}
La aplicación web progresiva falla en persistencia y periféricos, y que la nativa obliga a mantener una segunda base de código además de la web de los portales. La híbrida cumple las tres exigencias técnicas con una sola base de interfaz.

\subsubseccion{Decisión: Botonera de Pala y App SITCO Terreno}
Se selecciona la Botonera de Pala para registrar la declaración del operador desde la cabina y enviarla al Nodo por la red LTE privada. Si falta la identificación durante un relevo, el sistema registra el hecho y avisa al supervisor sin detener la operación. La botonera conserva declaraciones cifradas al menos 48 horas, actualiza la luz de destino en un máximo de dos segundos y permite bloquear el dispositivo remotamente en 15 minutos o menos. La credencial utilizada corresponde a SUP-17.

La App SITCO Terreno y la Consola de Portería comparten el núcleo de captura. En terreno, la aplicación guarda registros cifrados hasta 48 horas, protege las claves con TPM, permite validar al supervisor con PIN y reanuda la sincronización cuando vuelve la conexión. La consola agrega la comunicación con torniquetes y lectores.

Ambas interfaces usan objetivos táctiles de al menos 44 × 44 píxeles, alto contraste, íconos con texto, un máximo de tres interacciones por función y confirmación visible del estado de sincronización (RNF-USA-012, RNF-USA-013 y RNF-USA-014). Estos criterios se validan durante la marcha blanca con operadores de ambos turnos.

\subseccion{Módulos y límites de contexto}
Cada módulo es un límite de contexto: tiene su propio lenguaje, sus agregados raíz y su almacén, y sólo se relaciona con los demás mediante eventos publicados o interfaces versionadas (Evans, 2003). Ningún módulo lee el almacén de otro, si necesita un dato, se suscribe al evento que lo publica. Los módulos se agrupan según el problema que resuelven. Lo que cada módulo hace para quienes lo usan está en el Subdocumento 3 (secciones 3.4.2 y 3.4.3), aquí se describe lo que cada uno registra, publica y consume, y dónde se ejecuta.

\subsubseccion{Cadena del mineral: Operación Mina, Trazabilidad y Embarque, y Conciliación Metalúrgica}
La cadena del mineral comprende el registro de su origen, su paso por el stock y la planta, y la conformación del lote que se embarca. Operación Mina registra los viajes y la productividad, Trazabilidad y Embarque conserva la relación entre los viajes, las capas de stock y los lotes, y Conciliación Metalúrgica contrasta sus tonelajes y leyes. La Figura \ref{fig:cadena-mineral} presenta el recorrido general y sitúa las responsabilidades de los tres módulos, y las Figuras \ref{fig:cadena-mineral-a} y \ref{fig:cadena-mineral-b} lo detallan por tramos.

\begin{horizontal}
\csname ov@hoja\endcsname{330.2mm}{215.9mm}
\expandafter\def\csname ov@altofig\endcsname{0.80}
\figura{27.52}{s4-revision-10}{Cadena del mineral: vista general.}{fig:cadena-mineral}{Fuente: elaboración propia a partir de Compañía Minera Altos de Aranda (2026b, cap. 4, pp. 7-9).}

\clearpage
\addtocounter{figure}{-1}
\begingroup
\def\thefigure{4.10a}
\def\theHfigure{4.10a}
\figura{27.52}{s4-revision-10a}{Cadena del mineral: plan, carguío y producción.}{fig:cadena-mineral-a}{Fuente: elaboración propia a partir de Compañía Minera Altos de Aranda (2026b, cap. 4, pp. 7-9).}
\endgroup
El viaje mantiene su origen al pasar por el stock o alimentar directamente la planta.

\clearpage
\addtocounter{figure}{-1}
\begingroup
\def\thefigure{4.10b}
\def\theHfigure{4.10b}
\figura{27.52}{s4-revision-10b}{Cadena del mineral: balance, despacho y cierre.}{fig:cadena-mineral-b}{Fuente: elaboración propia a partir de Compañía Minera Altos de Aranda (2026b, cap. 4, pp. 7-9).}
\endgroup
El despacho y la conciliación siguen sus ramas y se reúnen al completar el cierre.
\end{horizontal}
El viaje es la primera unidad de trazabilidad. El material descargado se incorpora a la capa del día del stock o alimenta directamente la planta. Cuando se retoma del stock, conserva el vínculo con sus capas de origen. Esa relación continúa en el lote de planta, el concentrado, el camión a puerto y el lote comercial. Para el despacho a puerto, Trazabilidad y Embarque envía los datos de la guía a SAP S/4HANA, que la emite con su folio y QR. SITCO guarda el folio recibido y registra el despacho. En paralelo, se calculan el cumplimiento del plan y los balances metalúrgicos, sin que el cumplimiento del turno espere los análisis pendientes del laboratorio.

\subsubseccion{Operación Mina}
Operación Mina escribe en la nube la versión del plan, que recibe del conector de planificación y copia al Nodo de Faena como dato maestro, y escribe en el Nodo los hechos que nacen en la faena: la declaración de material que llega de la Botonera de Pala por la puerta operacional, el viaje y el ciclo que entrega el conector del despacho, y la perforación y la tronadura registradas en la App SITCO Terreno. Asocia cada cambio de destino del despacho a la declaración previa del operador e identifica los viajes cargados en el contacto entre dos polígonos.

En las detenciones no programadas de mina también se registra la causa raíz, quién la determinó y cuándo. Si pasan 72 horas desde la detención sin esa causa, el sistema genera una alerta.

También clasifica las demoras provenientes del despacho mediante la taxonomía única acordada con mina y planta, dejando al supervisor del turno sólo las que no puede resolver. La Figura~\ref{fig:entradas-responsabilidades-intercambios-operacion-mina} reúne estas responsabilidades y muestra el viaje que publica para Trazabilidad y Embarque, junto con el balance que recibe de Conciliación Metalúrgica.

\begin{horizontal}
\figura{22.44}{fig-4-11-entradas-responsabilidades-intercambios-operacion-mina}{Entradas, responsabilidades e intercambios de Operación Mina.}{fig:entradas-responsabilidades-intercambios-operacion-mina}{Fuente: elaboración propia.}
\end{horizontal}
El módulo conserva dos grupos de información: el origen del mineral, asociado a viajes, polígonos, perforación y tronadura, y la productividad, expresada en demoras, cumplimiento del plan y entrega de turno digital. Publica el cumplimiento al inicio del turno siguiente, con los registros operacionales disponibles (RNF-DES-008 y RNF-DES-010). El balance recibido de Conciliación Metalúrgica complementa esa información sin condicionar su publicación.

Los pozos de las siete perforadoras sin instrumentación se registran en la App SITCO Terreno, y los de las dos instrumentadas llegan en forma automática por el conector del despacho de flota (RF-TER-006),y, desde la Etapa 2, se complementan con el conector de telemetría de su fabricante (SUP-10). La tronadura y la liberación del sector se registran en la App SITCO Terreno y se notifican al despacho y a Operación Mina sin recurrir a la radio. Geología descarga los pozos y su índice de dureza desde el Portal Web SITCO en formato estructurado.

\subsubseccion{Trazabilidad y Embarque}
Trazabilidad y Embarque consume el viaje que publica Operación Mina y escribe en el Nodo de Faena la capa de stock, el lote de planta y el despacho a puerto, en la nube escribe el lote comercial y el embarque, a partir de las lecturas que guarda el gabinete de puerto. La regla de mezcla de los stocks y su verificación mensual contra el laboratorio (SUP-03) se describen en el Subdocumento 3 (sección 3.4.2.2).

Conforma el lote de planta de cada turno, prepara los datos de la guía de despacho de cada camión a puerto, incluido el código del lote, que SAP S/4HANA emite con su folio e imprime como QR, asocia el camión al lote comercial en el acopio y registra el embarque con su acta firmada. Retiene el lote con análisis rechazado hasta contar con uno válido (SUP-06) y mantiene un vector de elementos que incluye el arsénico, para estimar la penalización antes de conformar un lote (SUP-23). La Figura~\ref{fig:cadena-custodia-controles-trazabilidad-embarque} muestra las unidades que enlaza y los controles de calidad que aplica a esa cadena.

\begin{horizontal}
\figura{22.44}{fig-4-12-cadena-custodia-controles-trazabilidad-embarque}{Cadena de custodia y controles de Trazabilidad y Embarque.}{fig:cadena-custodia-controles-trazabilidad-embarque}{Fuente: elaboración propia.}
\end{horizontal}
La cadena de custodia conserva la relación entre el viaje, la capa de stock, el lote de planta, el concentrado, el camión a puerto, el lote comercial y el embarque. La lectura del QR en el acopio identifica la guía emitida por el SAP y permite registrar la recepción y la asociación al lote comercial. Desde la Etapa 2, el comprador consulta la trazabilidad autorizada mediante su portal.

El módulo también registra la liquidación que informa el comprador, vincula cada componente al lote y la entrega a SAP sin calcularla. Esta cadena debe permitir reconstruir el origen de un lote embarcado en menos de dos horas, frente a las semanas que toma hoy.

\subsubseccion{Conciliación Metalúrgica}
El módulo recibe el tonelaje de la balanza de alimentación por el conector del historiador, los pesajes por los conectores del despacho y de la báscula, y los resultados por el conector del laboratorio.

El módulo envía al sistema de laboratorio la identificación de cada muestra identificador, punto de toma, fecha y hora y lote de turno y recibe los resultados asociados sin transcripción. Conciliación conserva en la nube las solicitudes y resoluciones de arbitraje, vinculadas a los análisis originales. La liquidación informada por el comprador continúa en Trazabilidad y Embarque.

Los contratos conservan los identificadores y certificados originales de cada análisis, incluido el dirimente. La resolución del arbitraje se entrega a Trazabilidad por la interfaz de Conciliación y como evento, sin lectura directa de su base ni cálculo del pago del comprador en SITCO. La Figura \ref{fig:fuentes-medicion-resultados-conciliacion-metalurgica} muestra las fuentes de medición y los resultados del módulo.

\begin{horizontal}
\figura{22.44}{fig-4-13-fuentes-medicion-resultados-conciliacion-metalurgica}{Fuentes de medición y resultados de Conciliación Metalúrgica.}{fig:fuentes-medicion-resultados-conciliacion-metalurgica}{Fuente: elaboración propia.}
\end{horizontal}
El balance del turno se calcula al cierre en el Nodo de Faena. Como el laboratorio entrega sus resultados entre 4 y 36 horas después de tomada la muestra, se publica como provisional, identifica los análisis pendientes y marca la ley estimada del plan como supuesto. Al recibirse cada análisis validado pendiente, el balance se recalcula automáticamente.

El balance mensual se cierra en la nube con factores versionados, y la conciliación aprobada se entrega a SAP por su conector.

\subsubseccion{Equipos y Mantenimiento}
El aviso de falla se registra en la App SITCO Terreno y Equipos y Mantenimiento lo crea en SAP mediante su conector. El módulo clasifica automáticamente las detenciones de planta con un motor de reglas deterministas y, para el residuo, con un clasificador basado en LightGBM sobre los datos de la telemetría, el despacho y el historiador, con la taxonomía única acordada con mina y planta, y el supervisor completa las causas que las reglas no logran resolver (SUP-05). El clasificador de detenciones es uno de dos componentes de inteligencia artificial, junto con la verificación facial de portería. Su riesgo de clasificación se trata con NIST AI RMF 1.0 (RT-18.04). El clasificador se incorpora sin cambiar los contratos: se construye en los meses 13 a 18 y se activa en la mina en el mes 22 y en la planta en el mes 27. Cada detención conserva su clasificación, diagnóstico y escalamientos, con el historial de cambios.

La telemetría se recibe en el estándar ISO 15143-3 cuando el fabricante lo soporta (International Organization for Standardization, 2020). El módulo se ejecuta sólo en la nube, porque su telemetría llega desde los portales de telemetría en la nube, fuera de la faena, y entra en la Etapa 2 (Subdocumento 3, sección 3.2.2.1, Tabla \refext{S3}{tab:asignacion-etapa}).

\subsubseccion{Cumplimiento y personas: Acreditación y Acceso, Seguridad y Salud Ocupacional, y Cumplimiento Ambiental}
Acreditación y Acceso habilita personas según la documentación del Portal de Contratistas (Subdocumento 3). En portería, los torniquetes consultan el módulo en hasta 1,5 segundos o usan la lista local del gabinete sin enlace que el guardia controla excepciones desde la Consola de Portería. La verificación facial uno a uno se incorporará en la Etapa 2, previo acuerdo sindical (SUP-13).

Seguridad y Salud Ocupacional recibe registros de la App SITCO Terreno, aun sin señal, y gestiona su firma electrónica avanzada mediante Administración y Servicios Comunes. El Nodo de Faena conserva cada alerta de fatiga y la escala al responsable de turno si no hay acción documentada en 30 minutos, sin duplicar el escalamiento. Si fallan los tres caminos, se envían al correo de respaldo del satélite geoestacionario y se registran con su referencia y hora originales.

Cumplimiento Ambiental integra 186 instrumentos, laboratorios externos, telemetría y SAP para calcular la huella de carbono mensual (Subdocumento 3, sección 3.4.2.7). Los tres módulos responden al tercer problema diagnosticado: mantienen registros firmados, fiscalizables y conservados durante diez años.

\subsubseccion{Módulos transversales: Indicadores y Centro Integrado, y Administración y Servicios Comunes}
Indicadores y Centro Integrado no origina datos: consume los eventos que publican los demás módulos y los organiza en un modelo analítico separado de los almacenes transaccionales. Entrega los tableros de turno y la vista del CIO en la Etapa 1, y los indicadores de gestión en la Etapa 2 (Subdocumento 3, sección 3.4.3). Desde cualquier indicador permite bajar hasta la transacción que lo originó, y cada indicador tiene una definición versionada. Si un indicador necesita un dato que ningún módulo publica, la corrección es publicar el evento y no leer el almacén de otro módulo.

Administración y Servicios Comunes presta a los demás módulos los servicios comunes descritos en el Subdocumento 3, como interfaces síncronas para firmar, notificar, buscar y parametrizar, y como eventos de parámetro aprobado, maestro actualizado y documento firmado (Tabla \ref{tab:interfaces-sincronas-eventos-publicados-datos-consumidos}). Es el sitio escritor de los maestros y los parámetros, que escribe en la nube y baja al Nodo de Faena junto con las identidades que los módulos locales necesitan durante un corte.

\subsubseccion{Interfaces que expone y consume cada módulo}
Cada módulo expone interfaces síncronas sólo para las consultas y órdenes que un canal necesita en el momento, y publica eventos para todo lo demás. La siguiente tabla resume las interfaces principales.

\begin{tabla}{I{0.195}I{0.271}I{0.267}I{0.267}}{Interfaces síncronas, eventos publicados y datos consumidos por cada módulo.}{tab:interfaces-sincronas-eventos-publicados-datos-consumidos}
  {\cab{Módulo} & \cab{Interfaz síncrona que expone} & \cab{Eventos que publica} & \cab{Eventos y datos que consume}}
  {Fuente: elaboración propia.}
Operación Mina & Registrar declaración y consultar plan y cumplimiento & Plan publicado, material declarado, viaje registrado, detención registrada y ciclo de carguío registrado & Plan, despacho con la posición de la pala, habilitación del operador y parámetros \\
Trazabilidad y Embarque & Consultar la cadena de un lote y preparar la guía para SAP & Stock actualizado, lote conformado o bloqueado, embarque registrado y liquidación recibida & Viajes, análisis validados, pesajes y datos del comprador \\
Conciliación Metalúrgica & Consultar y aprobar balances y solicitar arbitraje & Pesaje capturado, análisis validado, balance de turno calculado y conciliación mensual cerrada & Viajes, lotes, balanza de alimentación, báscula y laboratorios \\
Equipos y Mantenimiento & Consultar la vista única y registrar aviso de falla & Estado de equipo actualizado, aviso de falla, detención clasificada y alerta de mantenimiento & Telemetría, despacho, pautas de SAP y detenciones \\
Acreditación y Acceso & Consultar la habilitación en portería y el expediente con sus versiones documentales & Habilitación otorgada o bloqueada, paso de portería e ingreso excepcional & Maestro de personal, lecturas de torniquete y versiones documentales \\
Seguridad y Salud Ocupacional & Registrar incidente o inspección y cerrar caso de fatiga & Registro de seguridad creado, caso de fatiga abierto y alerta escalada & Alertas de fatiga, identidad del autor y base de tratamiento \\
Cumplimiento Ambiental & Consultar compromisos, reportes y huella & Medición recibida, huella calculada, reporte firmado y cumplimiento enviado & Instrumentos, laboratorios externos, combustible, viajes y embarques \\
Indicadores y Centro Integrado & Consultar indicadores y tableros & Indicador publicado y métrica de adopción calculada & Eventos de todos los módulos y datos analíticos de SAP \\
Administración y Servicios Comunes & Firmar, notificar, buscar y parametrizar & Parámetro aprobado, maestro actualizado, documento firmado y notificación entregada & Solicitudes de los demás módulos, identidades externas y plantillas \\
\end{tabla}
El resumen presentado muestra que las interfaces síncronas son pocas y de consulta u orden puntual, mientras el flujo del dato entre módulos ocurre por eventos. Esa proporción permite que un módulo siga operando cuando otro está degradado o, en la faena, cuando la nube no está disponible.

\subseccion{Arquitectura de integración}
La integración resuelve tres relaciones: la de SITCO con los sistemas que la compañía mantiene, la de los módulos entre sí y la de la faena con la nube. Las tres siguen el mismo principio: el dato viaja como un hecho con identidad propia, y ningún componente depende de que otro esté disponible en ese instante, salvo cuando el canal necesita la respuesta para continuar.

\subsubseccion{Integrador de Sistemas: un conector por cada uno de los nueve sistemas que se mantienen, y una banda de migración para las planillas}
Se dispone de nueve conectores operativos y de un conector de migración para la carga única de las planillas y carpetas compartidas (Subdocumento 5, sección 5.3), y cada conector incorpora una capa anticorrupción (ACL) que traduce los datos externos al modelo de SITCO y aísla los módulos de los cambios del sistema de origen. La ubicación de cada conector en la nube, en el Nodo de Faena o en un gabinete se muestra en las Figuras \ref{fig:conectores-ejecutados-nube-relacion-bus-eventos} y \ref{fig:conectores-faena-procesamiento-local-carga-unica}, y la Tabla \ref{tab:conectores-integrador-sistemas-etapa-lugar-ejecucion} detalla, para cada sistema, la etapa en que entra, dónde se ejecuta su conector, el modo de integración y el comportamiento de la solución cuando ese sistema no responde.

\begin{horizontal}
\figura{22.44}{fig-4-14-conectores-ejecutados-nube-relacion-bus}{Conectores ejecutados en la nube y relación con el bus de eventos.}{fig:conectores-ejecutados-nube-relacion-bus-eventos}{Fuente: elaboración propia.}
\end{horizontal}
En la nube se ejecutan los conectores de SAP, telemetría, fatiga y planificación minera. El bus transporta los eventos entre estos conectores y los módulos de SITCO. La versión vigente del plan minero y los datos de referencia se envían al Nodo de Faena para que la operación local conserve la información necesaria durante un corte de enlace.

Las Figuras \ref{fig:conectores-faena-procesamiento-local-carga-unica} y \ref{fig:conectores-faena-procesamiento-local-carga-unica-b} muestran cómo se atienden los sistemas ubicados en la faena y cómo se conserva la información antes de enviarla a la nube.

\begin{horizontal}
\expandafter\def\csname ov@altofig\endcsname{0.80}
\begingroup
\def\thefigure{4.15a}
\def\theHfigure{4.15a}
\figura{22.44}{s4-revision-15a}{Conectores operativos de faena.}{fig:conectores-faena-procesamiento-local-carga-unica}{Fuente: elaboración propia.}
\endgroup
Cada conector traduce el sistema existente y entrega los hechos al procesamiento local.

\clearpage
\addtocounter{figure}{-1}
\begingroup
\def\thefigure{4.15b}
\def\theHfigure{4.15b}
\figura{22.44}{s4-revision-15b}{Condiciones de operación y migración.}{fig:conectores-faena-procesamiento-local-carga-unica-b}{Fuente: elaboración propia.}
\endgroup
Las condiciones de acceso, continuidad y migración acompañan a los conectores operativos.

\end{horizontal}
Los conectores de despacho, laboratorio, báscula y control de acceso se ejecutan en el Nodo de Faena o en el gabinete correspondiente. El conector del historiador se ejecuta en la zona desmilitarizada industrial de planta. Esta ubicación permite mantener los intercambios locales durante una interrupción del enlace con la nube. La consulta de habilitación en portería se resuelve de forma síncrona en su gabinete. El conector del historiador consulta un espejo alimentado mediante un diodo de datos, que impide escribir hacia la red de control.

Los módulos locales guardan el hecho y su evento en una misma transacción. La bandeja de salida conserva los eventos pendientes y los publica en el bus al recuperarse el enlace. Las flechas entre ambas partes representan este envío y la recepción del plan vigente y los datos de referencia.

La banda de migración corresponde al décimo sistema: las planillas y carpetas compartidas. Su contenido se valida y se carga una sola vez, sin intercambio recurrente con el bus.

La Tabla \ref{tab:conectores-integrador-sistemas-etapa-lugar-ejecucion} detalla, para cada conector, la etapa, el lugar de ejecución, el modo de integración y el comportamiento ante falla.

\begin{tabla}{I{0.162}I{0.074}I{0.181}I{0.205}I{0.378}}{Conectores del Integrador de Sistemas: etapa, lugar de ejecución, modo y comportamiento ante falla.}{tab:conectores-integrador-sistemas-etapa-lugar-ejecucion}
  {\cab{Sistema existente} & \cab{Etapa} & \cab{Dónde se ejecuta} & \cab{Modo} & \cab{Comportamiento ante falla}}
  {Fuente: elaboración propia con datos de Compañía Minera Altos de Aranda (2026b, cap. 5, p. 10) y del subdocumento 3}
SAP S/4HANA & 1 y 2 & Nube & Asíncrono, por la interfaz de SAP & Cola persistente, la emisión contable se posterga \\
Despacho de flota & 1 & Nodo de Faena & Asíncrono, por eventos del despacho & Cola en el conector; los viajes se completan al volver el despacho \\
Gestión de laboratorio & 1 & Nodo de Faena & Asíncrono en ambos sentidos & Cortacircuito y cola para el lote queda retenido sin análisis válido \\
Portales de telemetría & 2 & Nube & Asíncrono, por lectura periódica de cada portal & La vista única indica la fuente que falta y sigue disponible \\
Control de fatiga y somnolencia & 2 & Nube & Asíncrono en ambos sentidos & La atención continúa en el Nodo y sin enlace, la alerta llega por el correo de respaldo (4.1.6.3) \\
Historiador de proceso & 1 & Zona desmilitarizada industrial de planta & Suscripción OPC UA de sólo lectura sobre el espejo & El historiador conserva las señales y el balance se recalcula \\
Planificación minera & 1 & Nube & Asíncrono, por publicación de versión & Rige la última versión recibida, marcada con su fecha \\
Software de la báscula & 1 & Nodo de Faena & Asíncrono, por lectura de pesajes & El factor semanal usa los pesajes disponibles \\
Control de acceso de portería & 1 y 2 & Gabinete de portería & Síncrono en el gabinete de portería & Ingreso manual autorizado y registrado en la Consola de Portería \\
Planillas y carpetas compartidas & 1 y 2 & Nube & Carga única en la migración & No aplica \\
\end{tabla}
La Etapa 1 habilita conectores para SAP, despacho, laboratorio, historiador, planificación, báscula y control de acceso, además de la carga única de planillas. La Etapa 2 incorpora telemetría y control de fatiga, y amplía los conectores de SAP con mantenimiento y costos y portería con verificación facial, sin modificar SAP ni la arquitectura. De los diez sistemas existentes, sólo portería opera en forma síncrona y dentro de la faena, los demás encolan o informan la antigüedad del dato. Se asume disponibilidad continua, salvo mantenciones de SAP y ventanas de telemetría y planificación, por confirmar en el levantamiento (SUP-10 y SUP-16).

El Integrador también atiende instrumentación ambiental, geotécnica y de caudal, laboratorios externos, notificaciones, radio y megafonía, y firma electrónica avanzada. En total, operan catorce integraciones: nueve sistemas existentes y estas cinco adicionales.

Todos los conectores aplican tiempos de espera de 1 segundo en consultas y 2 en escrituras, tres reintentos con retroceso exponencial desde 200 milisegundos y variación aleatoria de 20 \%, cortacircuito, mamparo de aislamiento y un máximo de 50 llamadas por segundo por dependencia. Ante una falla, encolan solicitudes y el módulo usa el último dato, marcado con su antigüedad. Toda entrada y salida conserva un identificador de correlación, cuando no existe interfaz, se utiliza una exportación automática programada (SUP-16).

La Tabla~\ref{tab:modo-ventana-comportamiento-ante-falla-integraciones} presenta los mismos datos para las cinco integraciones adicionales.

\begin{tabla}{I{0.207}I{0.281}I{0.188}I{0.324}}{Modo, ventana y comportamiento ante falla de las integraciones adicionales.}{tab:modo-ventana-comportamiento-ante-falla-integraciones}
  {\cab{Integración} & \cab{Modo} & \cab{Ventana de la contraparte} & \cab{Comportamiento ante falla}}
  {Fuente: elaboración propia.}
Instrumentación ambiental y caudalímetros & Asíncrono, por lectura en los gabinetes y en el Nodo de Faena & Intermitente en relaves, por radioenlace & El gabinete guarda las mediciones y las reenvía al volver el enlace \\
Laboratorios externos & Asíncrono, por la puerta de exposición & La que declare cada laboratorio & El informe se carga al volver la conexión y el reporte espera su dato \\
Radio y megafonía de faena & Síncrono desde el Nodo de Faena & Continua, igual que la operación & Rige el procedimiento de evacuación vigente de la faena \\
Mensajería de notificaciones & Asíncrono con reintento & La que garantice cada proveedor & Reintento sin duplicados y registro de cada entrega (RNF-DES-018) \\
Firma electrónica avanzada & Asíncrono, por la interfaz del prestador acreditado & La que garantice el prestador & El documento queda pendiente de firma y se emite al volver el servicio \\
\end{tabla}
Observando la tabla, se puede concluir que sólo la radio y megafonía es síncrona, y es la única que no admite espera, por eso su conector se ejecuta en el Nodo de Faena, mientras la firma deja el documento pendiente y lo emite al volver el servicio.

\subsubseccion{Servicios y llamadas síncronas}
Una interacción es síncrona sólo cuando el canal no puede continuar sin la respuesta y el umbral de tiempo lo exige y las demás son eventos. La Tabla~\ref{tab:llamadas-sincronas-entre-componentes-umbral} presenta las llamadas síncronas de la solución.

\begin{tabla}{I{0.229}I{0.335}I{0.208}I{0.229}}{Llamadas síncronas entre componentes y su umbral.}{tab:llamadas-sincronas-entre-componentes-umbral}
  {\cab{Consumidor} & \cab{Proveedor} & \cab{Propósito} & \cab{Umbral}}
  {Fuente: elaboración propia a partir de Overfore-Subdocumento2-Anexos.}
 & Operación Mina en el Nodo de Faena & Registrar la declaración y confirmarla con la luz & 2 segundos o menos en el percentil 95 (RNF-DES-001) \\
\multirow{-2}{=}{Botonera de Pala} & Copia de identidades del Nodo de Faena & Identificar al operador en el relevo & 2 segundos o menos en el percentil 95 (RNF-DES-004) \\
App SITCO Terreno & Módulo local en el Nodo de Faena o, en el acopio, agente del gabinete de puerto & Registrar el hecho & 2 segundos o menos en el percentil 95 \\
Torniquete, mediante el conector del control de acceso & Acreditación y Acceso en el gabinete de portería & Decidir si la persona está habilitada & 1,5 segundos o menos, también sin enlace (RNF-DES-003) \\
Conector de la báscula & Conciliación Metalúrgica en el Nodo de Faena & Registrar el pesaje & 3 segundos o menos (RNF-DES-002) \\
Portal Web SITCO & Vista materializada de Equipos y Mantenimiento & Mostrar la vista única del equipo & 5 segundos y 8 en modo degradado (RNF-DES-011) \\
\end{tabla}
Analizando la tabla, se observa que todas las llamadas síncronas de terreno se resuelven dentro del sitio donde ocurren, en la faena o en el acopio, sin un viaje de ida y vuelta a la nube, y que la vista única del equipo, que combina tres fuentes, no consulta en línea a ninguna de ellas. Toda escritura expuesta a reintentos exige una clave de idempotencia declarada por el cliente, con cero efectos duplicados ante un reenvío.

\subsubseccion{Mensajería y eventos de dominio}
Todo lo que no es una llamada síncrona viaja como evento de dominio por el bus de eventos. El bus entrega cada evento al menos una vez, conserva el orden dentro de cada agregado y descarta en el consumidor los duplicados, porque cada consumidor es idempotente. El orden por agregado se obtiene asignando el identificador del agregado como identificador de sesión del mensaje, y los mensajes que fallan de forma repetida pasan a una cola de fallidos por suscripción, sin perderse (Microsoft, s.f.-g).

En el Nodo de Faena y en los gabinetes, el hecho y el evento que lo anuncia se escriben en una sola transacción de la base local, en una bandeja de salida, y un proceso de envío los publica en el bus cuando hay enlace (Richardson, 2018). Así no existe un instante en que el hecho quede guardado sin su evento, ni un evento publicado sin su hecho. Durante un corte, la misma bandeja sirve de transporte entre los cinco módulos locales. Cada evento lleva los campos de la Tabla~\ref{tab:campos-obligatorios-evento-dominio}.

\begin{tabla}{I{0.184}I{0.449}I{0.368}}{Campos obligatorios de un evento de dominio.}{tab:campos-obligatorios-evento-dominio}
  {\cab{Campo del evento} & \cab{Contenido} & \cab{Para qué sirve}}
  {Fuente: elaboración propia.}
Identificador del evento & Generado en el dispositivo o en el Nodo al registrar el hecho & Deduplicar en el consumidor y en la sincronización \\
Agregado y secuencia & Identificador del agregado y número de secuencia & Conservar el orden y detectar huecos \\
Sello de tiempo de origen & Hora del hecho con la deriva registrada del reloj & Ordenar lo ocurrido durante un corte \\
Autor y sitio & Persona o sistema, y sitio de origen & Atribuir el hecho y aplicar la regla de conflicto \\
Versión de esquema & Versión semántica del contrato & Consumir versiones distintas sin romperse \\
Correlación y firma & Identificador de la operación y firma del registro & Seguir la operación y verificar su integridad \\
\end{tabla}
La tabla muestra que un evento puede llegar a la nube horas después de ocurrido sin ambigüedad sobre cuándo ocurrió, quién lo registró y si ya fue cargado. La fuente horaria común de la faena, el CIO y el puerto (SUP-25) es la que permite confiar en el sello de tiempo de origen.

\subsubseccion{Contratos}
Cada módulo mantiene sus propios datos y se integra sólo mediante contratos de interfaces y eventos (Tabla \ref{tab:interfaces-sincronas-eventos-publicados-datos-consumidos}). Los contratos clasifican la información, los datos confidenciales o personales sensibles se cifran, auditan y, cuando corresponde, se reemplazan por identificadores no personales para Indicadores y Centro Integrado. Las interfaces y eventos se generan y validan en cada compilación con OpenAPI 3.1 y AsyncAPI 3.0; las integraciones externas usan el portal de desarrolladores, datos sintéticos en pruebas y OAuth 2.1 en producción. Los conectores traducen los formatos de SAP, despacho y laboratorios al contrato interno.

\subsubseccion{Versionado}
Los contratos usan versionado semántico: agregar información opcional no rompe compatibilidad, quitar o cambiar un campo exige una versión mayor. Durante la transición se publican ambas versiones con el mismo identificador y cada consumidor procesa el hecho una sola vez. La versión anterior se retira tras seis meses de aviso y cuando ningún ambiente la use. Esta regla permite sincronizar Nodos que vuelven de un corte, incorporar eventos de la Etapa 2 sin modificar módulos existentes y adaptar cambios de formato externos mediante el conector. Pact Broker registra las versiones desplegadas en nube y Nodo; el Nodo sigue operando aunque no pueda acceder a ese registro (Anexo F).

\subsubseccion{Gobierno}
El Comité de Arquitectura, con participación del cliente, aprueba mensualmente los cambios incompatibles y los registra como ADR, el Arquitecto de Solución aprueba los cambios compatibles. Las correcciones de seguridad se resuelven según severidad y pueden bloquear temporalmente la interfaz afectada. Antes de desplegar una nueva versión, las pruebas de contrato verifican a sus consumidores y la validación se ejecuta también en el Nodo durante un corte. Las actas se entregan en dos días hábiles y el catálogo, los registros y las reglas quedan en poder del cliente al término del contrato.

\subsubseccion{Escenario de proceso: recorrido del dato en el ciclo del mineral}
El escenario sigue un dato del mineral desde el plan vigente hasta el embarque y la consulta del comprador. Se divide en cuatro partes cada una tiene su diagrama de secuencia y recibe lo que deja la anterior. Cada componente aparece con su sitio (nube, Nodo de Faena, gabinete de planta o gabinete de puerto), y los intermediarios que no cambian el resultado se indican en el texto de cada parte.

La primera parte explica cómo llega el destino planificado a la pala y cómo queda vinculado el cambio que decide el despacho. Antes de que empiece, el conector de planificación lee en la nube la salida del plan semanal. Operación Mina guarda la versión aprobada con fecha y hora y la copia al Nodo de Faena como dato maestro. La Figura \ref{fig:parte-1-recorrido-plan-vigente-declaracion} muestra la secuencia que sigue, desde el plan vigente hasta el viaje registrado.

\figura{16.05}{fig-4-16-parte-1-recorrido-plan-vigente}{Parte 1 del recorrido: plan vigente, declaración y viaje.}{fig:parte-1-recorrido-plan-vigente-declaracion}{Fuente: elaboración propia.}
Operación Mina envía el plan vigente a la botonera, que enciende la luz del destino en dos segundos o menos (mensajes 1 y 2). Cuando el operador ve que el material cambió y aprieta el botón, la botonera registra la declaración en Operación Mina. Esta la guarda con su evento y la botonera confirma con una luz en dos segundos o menos (mensajes 3 a 7). La llamada pasa por la puerta operacional del Nodo de Faena.

El evento “material declarado” llega al Integrador de Sistemas. El Integrador avisa al despachador en diez segundos o menos en el 95 \% de los casos con cobertura, sin radio, en el Portal Web SITCO de su puesto o en el propio despacho si su interfaz admite mensajes (mensajes 8 y 9). El despachador ordena la reasignación, y el despacho entrega el viaje con su pesaje a bordo, que el Integrador reenvía a Operación Mina (mensajes 10 a 12).

Operación Mina vincula el viaje con la declaración y con la versión del plan, sin que nadie escriba la razón del cambio (mensaje 13). Cuando la interfaz del despacho lo permite, Operación Mina también le envía los polígonos del plan vigente. Salvo la carga del plan, que ya está copiado en el Nodo antes de un corte, todo ocurre en la faena, así que un corte del enlace no interrumpe esta parte.

La segunda parte muestra cómo llegan las mediciones del turno a Conciliación Metalúrgica sin que nadie las transforme. Antes, el viaje descargado se suma a la capa del día del stock o pasa directo al chancador, según su destino. La Figura \ref{fig:parte-2-recorrido-mediciones-balance-turno} presenta las tres fuentes de tonelaje y el intercambio con el laboratorio.

\figura{13.09}{fig-4-17-parte-2-recorrido-mediciones-balance}{Parte 2 del recorrido: mediciones y balance del turno.}{fig:parte-2-recorrido-mediciones-balance-turno}{Fuente: elaboración propia.}
Conciliación recibe el pesaje a bordo de cada viaje (mensaje 1). Sólo los viajes que pasan por la báscula suman un pesaje de control. El Integrador lo lee del software de la báscula y lo registra en Conciliación en tres segundos o menos (mensajes 2 a 5). La lectura es asíncrona; sólo el registro es una llamada síncrona.

El tonelaje de la balanza de alimentación llega por la ruta del historiador, desde el conector del gabinete de planta (mensaje 6). Al cierre del turno, Conciliación calcula el balance provisional sin reemplazar ninguna medición original (mensaje 7). Luego envía al laboratorio las muestras del turno por el Integrador y recibe sus resultados sin transcripción (mensajes 8 a 11).

Los resultados llegan entre 4 y 36 horas después de la toma de muestra. Conciliación valida cada análisis, recalcula el balance (mensaje 12) y publica el evento “análisis validado”, con el que empieza la tercera parte. Operación Mina publica el cumplimiento del plan al inicio del turno siguiente, el balance lo complementa sin condicionar su publicación.

La tercera parte liga el lote de planta, el peso, el camión y la guía antes de que el despacho llegue a la nube. La Figura \ref{fig:parte-3-recorrido-camion-puerto-guia} empieza con el análisis validado y termina con el despacho guardado en la nube.

\figura{12.85}{fig-4-18-parte-3-recorrido-camion-puerto}{Parte 3 del recorrido: camión a puerto, guía y envío a la nube.}{fig:parte-3-recorrido-camion-puerto-guia}{Fuente: elaboración propia.}
Con el análisis validado, Trazabilidad y Embarque liga el concentrado del turno a su lote de planta (mensajes 1 y 2). El camión a puerto se pesa en la misma báscula, pero es otra medición. El Integrador la lee y la registra en Conciliación en tres segundos o menos, y Conciliación la envía a Trazabilidad (mensajes 3 a 7).

Trazabilidad liga el pesaje al camión y al lote y entrega los datos de la guía a SAP S/4HANA, que la emite con su folio y con el código del lote impreso como QR, SITCO guarda el folio (mensajes 8 y 9). La guía lleva el identificador del lote, el tonelaje húmedo y seco, la humedad, las leyes de cobre, oro, plata y arsénico, y el vínculo al origen.

El despacho a puerto viaja a la nube por la bandeja de salida y el bus de eventos, y Trazabilidad lo guarda allí (mensajes 10 y 11). Si el enlace está cortado, la bandeja conserva el evento y lo publica al volver.

La cuarta parte sigue el lote desde el acopio de Antofagasta hasta el embarque y la consulta del comprador. En el acopio, la tableta lee el código QR de la guía y registra cada hecho por la puerta operacional del gabinete de puerto. El gabinete guarda la lectura si falla el enlace del operador portuario (SUP-24). La Figura \ref{fig:parte-4-recorrido-acopio-embarque-consulta} empieza cuando el gabinete publica esos hechos.

\figura{16.05}{fig-4-19-parte-4-recorrido-acopio-embarque}{Parte 4 del recorrido: acopio, embarque y consulta.}{fig:parte-4-recorrido-acopio-embarque-consulta}{Fuente: elaboración propia.}
Trazabilidad recibe la recepción y asocia el camión al lote comercial que se está formando (mensajes 1 y 2). Cuando recibe la solicitud de programación, la acepta sólo si el lote tiene un análisis válido; si no, retiene el lote y pide una nueva toma de muestra (mensajes 3 y 4).

Al término de la carga, Trazabilidad conserva el hito y solicita la firma avanzada del acta a Administración y Servicios Comunes. Si el prestador no responde, el registro permanece con el cierre documental pendiente y se completa cuando vuelve el servicio. Con el acta firmada se publica el cierre del embarque.

Desde la Etapa 2, el comprador consulta el origen de un lote embarcado en el Portal para Compradores. El portal pide la cadena a Trazabilidad por la puerta de exposición (mensajes 10 y 11). El comprador recibe en dos segundos o menos un certificado derivado con sectores, fechas, stocks recorridos, análisis y cadena de custodia, sin la ley por polígono (mensajes 12 y 13, SUP-22).

\subsubseccion{Escenario de acceso en portería}
El ingreso a la faena comienza con una consulta sencilla para quien llega al torniquete: saber si su credencial sigue habilitada. El lector entrega la identificación al conector de control de acceso y Acreditación y Acceso responde con la información disponible en portería. La Figura \ref{fig:secuencia-validacion-credencial-porteria} muestra el recorrido de esa consulta y permite ver quién registra la causa cuando se niega el ingreso.

\begin{horizontal}
\figura{22.44}{fig-4-20-secuencia-validacion-credencial-porteria}{Secuencia de validación de la credencial en portería.}{fig:secuencia-validacion-credencial-porteria}{Fuente: elaboración propia.}
\end{horizontal}
El rechazo queda registrado con su causa y llega a la Consola de Portería antes de que el lector comunique la denegación. Cuando la credencial está habilitada, en la Etapa 1 el conector ordena abrir el torniquete y Acreditación y Acceso registra el ingreso. Desde la Etapa 2, y sólo si existe el acuerdo con el sindicato (SUP-13), la apertura queda pendiente de comprobar que quien presenta la credencial sea su titular.

La Figura \ref{fig:secuencia-verificacion-facial-porteria-etapa-2} muestra esa verificación de la Etapa 2 completa: cómo se compara el rostro con la plantilla cifrada y qué respuesta recibe cada dispositivo, hasta que el control de acceso ordena abrir el torniquete o comunica el rechazo. Las Figuras \ref{fig:secuencia-verificacion-facial-porteria-etapa-2-a}, \ref{fig:secuencia-verificacion-facial-porteria-etapa-2-b}, \ref{fig:secuencia-verificacion-facial-porteria-etapa-2-c} y \ref{fig:secuencia-verificacion-facial-porteria-etapa-2-d} la detallan: la captura y comparación, el rostro que no coincide y el rostro que coincide.

\begin{horizontal}
\csname ov@hoja\endcsname{330.2mm}{215.9mm}
\expandafter\def\csname ov@altofig\endcsname{0.88}
\figura{27.52}{s4-revision-21}{Verificación facial en portería: vista general.}{fig:secuencia-verificacion-facial-porteria-etapa-2}{Fuente: elaboración propia.}

\clearpage
\addtocounter{figure}{-1}
\begingroup
\def\thefigure{4.21a}
\def\theHfigure{4.21a}
\figura{27.52}{s4-revision-21a}{Verificación facial: captura y comparación del rostro.}{fig:secuencia-verificacion-facial-porteria-etapa-2-a}{Fuente: elaboración propia.}
\endgroup
El lector recibe la imagen y la compara con la plantilla cifrada.

\clearpage
\addtocounter{figure}{-1}
\begingroup
\def\thefigure{4.21b}
\def\theHfigure{4.21b}
\figura{27.52}{s4-revision-21b}{Verificación facial: el rostro no coincide.}{fig:secuencia-verificacion-facial-porteria-etapa-2-b}{Fuente: elaboración propia.}
\endgroup
Se registra el rechazo, se notifica a la portería y se comunica el acceso denegado.

\clearpage
\addtocounter{figure}{-1}
\begingroup
\def\thefigure{4.21c}
\def\theHfigure{4.21c}
\figura{27.52}{s4-revision-21c}{Verificación facial: el rostro coincide, apertura y registro del paso.}{fig:secuencia-verificacion-facial-porteria-etapa-2-c}{Fuente: elaboración propia.}
\endgroup
Con la coincidencia se abre el torniquete y se registra el paso concedido.

\clearpage
\addtocounter{figure}{-1}
\begingroup
\def\thefigure{4.21d}
\def\theHfigure{4.21d}
\figura{27.52}{s4-revision-21d}{Verificación facial: el rostro coincide, notificación y señal de acceso.}{fig:secuencia-verificacion-facial-porteria-etapa-2-d}{Fuente: elaboración propia.}
\endgroup
La consola confirma el aviso y se comunica el acceso concedido.
\end{horizontal}
Si el rostro no coincide, la consola recibe el aviso y el paso permanece cerrado. Si coincide, el conector ordena abrir el torniquete y Acreditación y Acceso guarda el ingreso. Así, el guardia puede consultar el resultado de cada intento sin tener que reconstruirlo a partir de la lectura del dispositivo.

\subsubseccion{Escenario Operación sin enlace}
Un corte del enlace con la nube no afecta del mismo modo a los tres canales de terreno. La Botonera de Pala y la App SITCO Terreno pueden seguir capturando registros para entregarlos al Nodo de Faena, si también pierden la señal local, los conservan en el dispositivo hasta recuperarla. En portería, la consola trabaja con la lista de habilitaciones disponible en su gabinete. La Figura \ref{fig:secuencia-operacion-durante-corte-enlace} muestra dónde continúa cada operación mientras no hay conexión con la nube. Las Figuras \ref{fig:secuencia-operacion-durante-corte-enlace-a} y \ref{fig:secuencia-operacion-durante-corte-enlace-b} separan lo que ocurre con enlace de lo que ocurre sin señal.

\begin{horizontal}
\csname ov@hoja\endcsname{330.2mm}{215.9mm}
\expandafter\def\csname ov@altofig\endcsname{0.80}
\figura{27.52}{s4-revision-22}{Operación durante un corte del enlace: vista general.}{fig:secuencia-operacion-durante-corte-enlace}{Fuente: elaboración propia.}

\clearpage
\addtocounter{figure}{-1}
\begingroup
\def\thefigure{4.22a}
\def\theHfigure{4.22a}
\figura{27.52}{s4-revision-22a}{Operación durante un corte: con enlace.}{fig:secuencia-operacion-durante-corte-enlace-a}{Fuente: elaboración propia.}
\endgroup
Con enlace al Nodo, el hecho se registra en el módulo local y se confirma al canal.

\clearpage
\addtocounter{figure}{-1}
\begingroup
\def\thefigure{4.22b}
\def\theHfigure{4.22b}
\figura{27.52}{s4-revision-22b}{Operación durante un corte: sin señal y reenvío.}{fig:secuencia-operacion-durante-corte-enlace-b}{Fuente: elaboración propia.}
\endgroup
Sin señal, la cola conserva el hecho cifrado y lo reenvía al volver la conexión.
\end{horizontal}
Los registros quedan en el dispositivo, el Nodo de Faena o el gabinete de portería, según dónde se originan y si hay señal local. Por eso, un corte del enlace con la nube no impide declarar material, registrar actividades en terreno ni controlar el ingreso a la faena. Cuando vuelve la conexión, los datos pendientes se envían a la nube.

\subseccion{Arquitectura de seguridad}
La seguridad se aplica sobre todas las capas y responde a cinco condiciones del caso: la mayoría de las personas con acceso a la faena son contratistas, la red de control de planta no se puede intervenir, la solución guarda datos de salud y biométricos, la operación sigue durante un corte y el área de TI del cliente no tiene especialistas de seguridad dedicados.

\subsubseccion{Modelo de confianza cero (Zero Trust) y su excepción}
La arquitectura adopta el modelo de confianza cero: verificación explícita de cada solicitud, privilegio mínimo y presunción de compromiso (Rose et al., 2020). Ninguna posición de red concede confianza. Una solicitud que proviene de la red de la faena se autentica igual que una que viene de Internet, los conectores se autentican como cualquier otro consumidor y los servicios se verifican entre sí con TLS mutuo.

Durante un corte, el punto central de decisión de políticas es inalcanzable, y verificar cada solicitud contra él detendría la mina. Por eso la decisión se delega en la copia local de políticas e identidades del Nodo de Faena, con vigencia máxima de 72 horas. La verificación se difiere y se repite al reconciliar: un hecho admitido durante el corte con una credencial que el maestro central había revocado se incorpora igual, porque ocurrió, pero queda marcado, alertado y trazable hasta la regla con que se admitió.

La revocación de una credencial se propaga en 24 horas o menos cuando hay enlace, y el bloqueo de un dispositivo, en 15 minutos o menos. Durante un corte, la baja no se propaga: se registra en la consola local que opera el área de seguridad del cliente, y el Nodo rechaza la credencial al vencer su copia de 72 horas.

\subsubseccion{Capa expuesta y superficie de exposición}
Todo servicio alcanzable desde fuera de la red de la compañía se publica a través de la capa de borde, y esa superficie está enumerada en la Tabla~\ref{tab:servicios-alcanzables-desde-fuera-red-compania}.

\begin{tabla}{I{0.171}I{0.292}I{0.278}I{0.258}}{Servicios alcanzables desde fuera de la red de la compañía.}{tab:servicios-alcanzables-desde-fuera-red-compania}
  {\cab{Servicio expuesto} & \cab{Quién accede} & \cab{Autenticación} & \cab{Control de borde}}
  {Fuente: elaboración propia.}
Portal de Contratistas & Empresas contratistas, laboratorios externos y operador portuario & Cuenta de empresa con doble factor & CDN, WAF, DDoS en capas 3, 4 y 7 del modelo OSI, y TLS 1.3 \\
Portal para Compradores & Compradores del concentrado & Cuenta nominada con doble factor y registro de cada consulta & Los mismos controles, más cuota por consumidor \\
Portal Web SITCO & Personal propio fuera de la faena & Identidad federada con doble factor y postura del dispositivo & Acceso de confianza cero, sin red privada virtual \\
Puerta de exposición & Integraciones autorizadas y laboratorios externos & OAuth 2.1 con credenciales de cliente o TLS mutuo & Validación de esquema e inspección de carga útil \\
\end{tabla}
Se observa que la superficie expuesta se reduce a cuatro servicios, cada uno con un control de borde y una autenticación propios, y que ningún almacén de datos, conector ni componente de la faena es alcanzable desde Internet. Ninguna integración usa una clave fija en la dirección web. La superficie declarada se contrasta con un escaneo externo de dominios, puertos y servicios antes de cada paso a producción.

\subsubseccion{Identidad por población: personal propio, operadores, contratistas y compradores}
Las poblaciones que usan SITCO tienen condiciones de acceso distintas, y un mecanismo único no las cubre. La Tabla~\ref{tab:mecanismo-identidad-cada-poblacion-usuaria} asocia cada población con su mecanismo de identidad.

\begin{tabla}{I{0.193}I{0.250}I{0.361}I{0.196}}{Mecanismo de identidad de cada población usuaria.}{tab:mecanismo-identidad-cada-poblacion-usuaria}
  {\cab{Población} & \cab{Cantidad} & \cab{Mecanismo de identidad} & \cab{Requerimiento}}
  {Fuente: elaboración propia a partir de Overfore-Subdocumento2-Anexos.}
Personal propio & 1.780 personas, más 115 del CIO & Directorio corporativo federado, doble factor y sin credenciales compartidas & RES-09 y RNF-SEG-003 \\
Supervisión con tableta & Supervisores de cada área y operadores de siete perforadoras & Credencial más PIN numérico & RES-09 y RF-CNX-003 \\
Operador de pala & Cinco palas por turno & Toque de credencial, contrastado con el operador asignado por el despacho, sin sesión ni acceso a datos & SUP-04, RNF-DES-004 y RES-09 \\
Operador de camión & 38 camiones & Ingreso al despacho de flota, no accede a SITCO & SUP-04 \\
Persona contratista en terreno & 2.400 permanentes y hasta 4.300 adicionales & Credencial más PIN numérico & SUP-08 \\
Empresa contratista en el portal & Unas noventa empresas & Identidad externa con doble factor y administración delegada & SUP-08 y RNF-SEG-003 \\
Compradores & Compradores del concentrado & Identidad externa con registro y recuperación autoservidos & RF-ACR-015 y RF-TRZ-016 \\
Operador del terminal portuario & Personal del recinto de terceros & Identidad externa autoservida, con acceso sólo a los eventos de embarque & RF-TRZ-014 \\
Laboratorios externos & Laboratorios de aire y de relaves & Identidad externa autoservida para cargar informes y credencial de cliente para la integración & RF-CUM-009 \\
Personal propio en la faena durante un corte & Hasta 445 personas por turno & Credencial más PIN numérico contra la copia de identidades del Nodo de Faena & RES-02 y RF-CNX-002 \\
\end{tabla}
Según los datos presentados, ninguna población de terreno usa contraseña alfanumérica: en terreno, quien entra a la App SITCO Terreno o a la Consola de Portería usa como doble factor la credencial física más el PIN. En la botonera, el segundo factor es el contraste con el despacho. El toque de credencial no abre una sesión ni da acceso a datos: sólo firma la autoría de la declaración, que la puerta operacional verifica con el certificado de la Botonera de Pala y contrasta con el operador que el despacho tiene asignado a esa pala, y toda discrepancia queda marcada para revisión en la reconciliación. El personal propio entra por Microsoft Entra ID federado con el directorio corporativo, y las empresas contratistas, los compradores, el operador portuario y los laboratorios externos por Microsoft Entra External ID, un plano separado. Durante un corte, la instancia local del Portal Web SITCO no alcanza al proveedor de identidad de la nube, por lo que el personal propio de la faena entra con su credencial y su PIN, validados contra la copia de identidades del Nodo de Faena, que vence a las 72 horas. Esa sesión local sólo alcanza a los cinco módulos del Nodo, y cada acción queda marcada para revisarla en la reconciliación, igual que un hecho de terreno. La elevación de privilegios de administración es temporal, requiere aprobación previa y graba la sesión.

\subsubseccion{Cifrado en tránsito y en reposo}
El cifrado cubre los tres lugares donde vive el dato: en tránsito, en los almacenes y en los dispositivos de terreno. La Tabla~\ref{tab:mecanismo-cifrado-custodia-claves-ambito} resume el mecanismo de cada ámbito.

\begin{tabla}{I{0.177}I{0.484}I{0.339}}{Mecanismo de cifrado y custodia de claves por ámbito.}{tab:mecanismo-cifrado-custodia-claves-ambito}
  {\cab{Ámbito} & \cab{Mecanismo} & \cab{Custodia de claves}}
  {Fuente: elaboración propia a partir de Overfore-Subdocumento2-Anexos.}
Tránsito externo & TLS 1.3, con TLS 1.0 y 1.1 rechazados & Certificados renovados en forma automática (RNF-SEG-008) \\
Tránsito entre servicios & TLS mutuo en la nube y en la faena & Autoridad certificadora en la bóveda, con rotación cada 90 días \\
Reposo en nube, faena y respaldos & Cifrado del 100 \% de los almacenes & Claves en módulo de seguridad de hardware (RNF-SEG-009) \\
Campos sensibles & Cifrado de campo de salud, exámenes preocupacionales, fatiga, alcohol y drogas, biometría, documento de identidad e identidad del operador en registros operacionales & Clave separada por categoría (RNF-SEG-001) \\
Dispositivos de terreno & Registro local cifrado y firmado & Claves en el módulo TPM, con borrado remoto (RNF-SEG-016) \\
\end{tabla}
La tabla muestra que acceder a la base de datos no basta para leer un dato sensible, porque su clave está separada y custodiada en hardware, y que la pérdida de una tableta no expone sus registros.

\subsubseccion{Controles: amenazas principales y casos de detección}
Las amenazas se modelan con STRIDE (suplantación, alteración, repudio, divulgación, denegación de servicio y elevación de privilegios) sobre los seis canales, los nueve módulos, las dos plataformas, los servicios técnicos y las catorce integraciones, y el modelo se revisa ante todo cambio de límite de contexto, de interfaz o de clasificación de información. La Tabla~\ref{tab:amenazas-mayor-consecuencia-componentes-expuestos-control} presenta las amenazas de mayor consecuencia en esta faena.

\begin{tabla}{I{0.296}I{0.307}I{0.397}}{Amenazas de mayor consecuencia, componentes expuestos y control.}{tab:amenazas-mayor-consecuencia-componentes-expuestos-control}
  {\cab{Amenaza} & \cab{Componentes expuestos} & \cab{Control}}
  {Fuente: elaboración propia.}
Préstamo de credencial entre personas & Acreditación y Acceso, Consola de Portería y Botonera de Pala & Detección por patrones de acceso para verificación facial desde la Etapa 2 \\
Alteración del registro de una tableta extraviada & App SITCO Terreno y Nodo de Faena & Registro firmado con clave del TPM para corrección sólo como registro nuevo \\
Corrección de un registro sin autor & Los nueve módulos & Auditoría inalterable con autor, fecha, motivo y valores anterior y posterior \\
Filtración de la ley por polígono & Portal para Compradores y Trazabilidad y Embarque & Filtro en el servicio y no en la pantalla en la publicación con aprobación comercial \\
Acceso indebido a datos de salud o fatiga & Seguridad y Salud Ocupacional y Acreditación y Acceso & Cifrado de campo y auditoría de cada consulta \\
Denegación de servicio sobre un portal & Portal de Contratistas y Portal para Compradores & Mitigación en el borde y rutas separadas de la ruta de terreno \\
Saturación del enlace al volver de un corte & Nodo de Faena y bus de eventos & Sincronización por lotes con límite de tasa y prioridad del tráfico en vivo \\
Alteración de un parámetro con una cuenta de administración & Administración y Servicios Comunes & Doble aprobación y elevación temporal con sesión grabada \\
Escritura en la red de control de proceso & Integrador de Sistemas & Conector sin operaciones de escritura y diodo de datos \\
\end{tabla}
La tabla detalla que la amenaza sobre la red de control se controla con una barrera física y otra de diseño, y no con una regla de configuración que pudiera revertirse. El modelo de amenazas se entrega y se versiona junto con el registro de decisiones.

Los eventos de seguridad se registran de forma centralizada e inalterable, con doce meses en línea y veinticuatro meses adicionales en archivo, y el Nodo de Faena los acumula durante un corte para reconciliarse después.

\subseccion{Operación sin enlace}
El modo sin enlace se activa ante la pérdida simultánea de las fibras y del respaldo satelital. Los dispositivos de terreno y el Nodo de Faena mantienen durante al menos 48 horas las funciones operacionales críticas, almacenando localmente los registros cifrados y firmados. La identidad y las habilitaciones se validan mediante una copia local controlada, mientras la nube continúa atendiendo las funciones externas con la última información sincronizada.

\subsubseccion{Funciones disponibles durante un corte}
Para cada función que debe seguir disponible durante un corte la Tabla \ref{tab:modulo-plataforma-sostiene-cada-funcion-durante} indica el módulo o la plataforma que la sostiene y el procedimiento con que se suple cuando no está disponible.

\begin{tabla}{I{0.231}I{0.385}I{0.385}}{Módulo o plataforma que sostiene cada función durante un corte.}{tab:modulo-plataforma-sostiene-cada-funcion-durante}
  {\cab{Función} & \cab{Módulo o plataforma que la sostiene} & \cab{Procedimiento}}
  {Fuente: elaboración propia.}
Registros de seguridad y acción ante alertas de fatiga & Seguridad y Salud Ocupacional, en el Nodo de Faena & Las alertas registradas se atienden y escalan localmente, y el caso de nube se consolida al restablecerse el enlace. \\
Acreditación de personas nuevas & Acreditación y Acceso en la nube; en el Nodo sólo se valida la habilitación vigente & Ingreso excepcional autorizado y registrado en la Consola de Portería \\
Vista del CIO con datos de faena & No disponible durante el corte (Indicadores y Centro Integrado opera sólo en la nube) & Coordinación por radio y teléfono satelital \\
Vista única del equipo con telemetría & No disponible durante el corte (Equipos y Mantenimiento opera sólo en la nube) & El supervisor usa el despacho para la App guarda los avisos de falla y las clasificaciones y los envía al volver el enlace. \\
Envíos a SAP y reportes a la autoridad & Nodo de Faena & Bandeja de salida con el orden de origen \\
Portal para Compradores & Trazabilidad y Embarque & Muestra el lote con la fecha del último dato recibido \\
\end{tabla}
La tabla describe que las funciones que siguen operando se apoyan en el Nodo de Faena o en el gabinete de portería, y que las que se suspenden tienen un procedimiento de suplencia y se completan al volver el enlace por la bandeja de salida. Además de estas funciones, el Nodo de Faena sostiene durante el corte dos servicios que el Esquema de Solución entrega a través de Administración y Servicios Comunes y de la copia de identidades. El primero es el anuncio de evacuaciones por el sistema de radio y megafonía de la faena, cuyo conector se ejecuta en el Nodo. El segundo es la instancia local del Portal Web SITCO para los cinco módulos locales.

\subsubseccion{Sincronización y reglas de resolución de conflictos}
Al volver el enlace, el Nodo de Faena envía lo acumulado en la bandeja de salida en el mismo orden en que ocurrió, por lotes y con límite de tasa, y da prioridad al tráfico en vivo sobre el acumulado para que la operación en curso no se degrade. Los duplicados se descartan por el identificador del evento, de modo que un reenvío no produce efectos dobles. Cuando dos registros se contradicen, la nube aplica las reglas fijas de la Tabla~\ref{tab:reglas-resolucion-conflictos-sincronizacion}.

\begin{tabla}{I{0.360}I{0.395}I{0.245}}{Reglas de resolución de conflictos en la sincronización.}{tab:reglas-resolucion-conflictos-sincronizacion}
  {\cab{Situación} & \cab{Regla} & \cab{Resultado registrado}}
  {Fuente: elaboración propia.}
Hecho de la operación registrado en la faena y en la nube & Prevalece el registro de la faena, porque el hecho ocurrió allí & Bitácora con ambos valores \\
Parámetro o dato maestro aprobado en la nube durante el corte & Prevalece la última versión aprobada en la nube, el Nodo la recibe al volver el enlace & Bitácora y aviso al dueño del maestro \\
Hecho registrado con una versión de parámetro o maestro ya reemplazada en la nube & Se conserva con la versión vigente al momento del hecho y se recalcula sólo si el dueño lo ordena & Bitácora y aviso al dueño del maestro \\
Hecho que cita un maestro inexistente en la nube & Queda en cuarentena hasta crear o corregir el maestro & Alerta al dueño del maestro \\
Hecho admitido con una credencial revocada durante el corte & Se incorpora y queda marcado & Alerta a seguridad \\
Evento duplicado & Se descarta por su identificador & Registro del descarte \\
\end{tabla}
La tabla permite reconocer que ninguna regla pierde información, cada decisión deja en la bitácora los valores en conflicto y la regla aplicada, y la resolución es determinista, por lo que dos reconciliaciones de los mismos datos dan el mismo resultado. El ciclo se cierra cuando el retraso de reconciliación de cada sitio llega a cero.

Para dimensionar la recuperación de un corte de 48 horas se considera un escenario conservador de aproximadamente 3,5 GB: 3,0 GB de series comprimidas y 0,5 GB de eventos y documentos. La sincronización debe completarse automáticamente en un máximo de cuatro horas desde el restablecimiento del enlace, sin degradar la operación en curso. El volumen y el tiempo de recuperación se verificarán mediante una prueba de desconexión y reconexión.

\seccion{Arquitectura física}
La arquitectura física asigna cada componente de la arquitectura lógica de SITCO Aranda a un servicio de nube, a un equipo de la faena o a un dispositivo de terreno, con su ubicación, su redundancia y su capacidad. El modelo es híbrido: la carga principal se ejecuta en Microsoft Azure y en la faena se despliega sólo lo que sostiene la operación cuando se corta el enlace con Antofagasta. La región primaria es Azure Chile Central, con tres zonas de disponibilidad, y la región secundaria es Azure Brazil South, en São Paulo. En la faena se instala una sala técnica nueva, que aloja el Nodo de Faena, y cuatro gabinetes de borde en portería, planta, depósito de relaves y acopio de puerto. En terreno se agregan pocos dispositivos, que usan la red LTE privada existente. La correspondencia de cada componente lógico con su emplazamiento físico se detalla en el Anexo B.

La arquitectura se representa en nueve zonas, de Z01 a Z09. La Figura \ref{fig:vista-general-arquitectura-fisica-sitco-aranda} presenta la vista general: reúne todos los sitios y los enlaces entre ellos. Las Figuras \ref{fig:vista-general-arquitectura-fisica-sitco-aranda-a} y \ref{fig:vista-general-arquitectura-fisica-sitco-aranda-b} la desglosan por zonas, y la Figura \ref{fig:vista-general-arquitectura-fisica-sitco-aranda-c} la resume sobre el plano de la faena.

\begin{horizontal}
\csname ov@hoja\endcsname{330.2mm}{215.9mm}
\expandafter\def\csname ov@altofig\endcsname{0.88}
\figura{27.52}{s4-revision-23}{Topología física: vista general de sitios y conexiones.}{fig:vista-general-arquitectura-fisica-sitco-aranda}{Fuente: elaboración propia.}

\clearpage
\addtocounter{figure}{-1}
\begingroup
\def\thefigure{4.23a}
\def\theHfigure{4.23a}
\figura{27.52}{s4-revision-23a}{Topología física: zonas Z01 a Z06.}{fig:vista-general-arquitectura-fisica-sitco-aranda-a}{Fuente: elaboración propia.}
\endgroup


\clearpage
\addtocounter{figure}{-1}
\begingroup
\def\thefigure{4.23b}
\def\theHfigure{4.23b}
\figura{27.52}{s4-revision-23b}{Topología física: zonas Z07 a Z09.}{fig:vista-general-arquitectura-fisica-sitco-aranda-b}{Fuente: elaboración propia.}
\endgroup


\clearpage
\addtocounter{figure}{-1}
\begingroup
\def\thefigure{4.23c}
\def\theHfigure{4.23c}
\figura{27.52}{s4-revision-23c}{Topología física: vista resumida de la faena.}{fig:vista-general-arquitectura-fisica-sitco-aranda-c}{Fuente: elaboración propia.}
\endgroup
\end{horizontal}
El conjunto de vistas distingue nube, faena, puerto, terreno y red de control. Las fibras A y B usan proveedores y trazados distintos, y el enlace satelital constituye el tercer camino. Los gabinetes y la red interna continúan operando durante un corte, el CIO y los usuarios externos consultan directamente la información ya sincronizada en la nube. Las Figuras \ref{fig:actores-caminos-acceso-sobre-topologia-fisica} y \ref{fig:actores-caminos-acceso-sobre-topologia-fisica-b} muestran los actores y sus caminos de acceso sobre la topología física.

La Tabla \ref{tab:zonas-laminas-fisicas} permite pasar de la vista general al detalle de cada zona.

\begin{tabla}{I{0.10}I{0.48}I{0.42}}{Correspondencia entre zonas y figuras de detalle.}{tab:zonas-laminas-fisicas}
{\cab{Zona} & \cab{Ubicación o función} & \cab{Figuras de detalle}}
{Fuente: elaboración propia.}
Z01 & Accesos externos y WAN & Figura \ref{fig:accesos-externos-enlaces-wan} \\
Z02 & Nube principal en Chile Central & Figura \ref{fig:nube-principal-chile-central-acceso-produccion}, Figura \ref{fig:nube-principal-chile-central-operacion-produccion} \\
Z03 & Recuperación en Brazil South & Figura \ref{fig:recuperacion-brazil-south} \\
Z04 & Sala técnica y Nodo de Faena & Figura \ref{fig:sala-tecnica-nodo-faena}, Figura \ref{fig:plano-distribucion-sala-tecnica-principal} \\
Z05 & Gabinetes de portería, planta y relaves & Figura \ref{fig:gabinetes-borde-porteria-planta-relaves} \\
Z06 & Acopio de puerto & Figura \ref{fig:gabinete-borde-acopio-puerto} \\
Z07 & Planta y red de control & Figura \ref{fig:planta-red-control-lectura-solo-sentido} \\
Z08 & Dispositivos de terreno y LTE privada & Figura \ref{fig:terreno-dispositivos-botonera-pala-tabletas-rugerizadas} \\
Z09 & Energía de respaldo & Figura \ref{fig:energia-respaldo-sala-tecnica-gabinetes} \\
\end{tabla}
Cada código conserva el mismo significado en la vista general, en los caminos de acceso y en el detalle. La Figura \ref{fig:entrega-software-cinco-ambientes-despliegue} muestra por separado la entrega de software y sus ambientes.

\begin{horizontal}
\csname ov@hoja\endcsname{330.2mm}{215.9mm}
\expandafter\def\csname ov@altofig\endcsname{0.80}
\begingroup
\def\thefigure{4.24a}
\def\theHfigure{4.24a}
\figura{27.52}{s4-revision-24a}{Caminos de acceso: fuera de la faena e instalaciones de la faena.}{fig:actores-caminos-acceso-sobre-topologia-fisica}{Fuente: elaboración propia.}
\endgroup
Las rutas externas llegan a los portales de la nube; las oficinas y la portería usan la instancia local cuando falta el enlace.

\clearpage
\addtocounter{figure}{-1}
\begingroup
\def\thefigure{4.24b}
\def\theHfigure{4.24b}
\figura{27.52}{s4-revision-24b}{Caminos de acceso: acopio de puerto y terreno.}{fig:actores-caminos-acceso-sobre-topologia-fisica-b}{Fuente: elaboración propia.}
\endgroup
El puerto usa su enlace con respaldo 4G; las botoneras y las tabletas llegan al Nodo por la red LTE privada existente.
\end{horizontal}
Las Figuras \ref{fig:actores-caminos-acceso-sobre-topologia-fisica} y \ref{fig:actores-caminos-acceso-sobre-topologia-fisica-b} muestran un punto de entrada por población. Los portales externos pasan por el borde de nube y la faena dispone de una instancia local del Portal Web SITCO. Los nueve módulos se ejecutan en Chile Central y los cinco críticos también en el Nodo. El mapeo completo se mantiene en el Anexo B.

\subseccion{Especificaciones Implementos a proveer (Hardware y Software)}
Las condiciones del sitio que acotan la elección son cinco. La primera es la altura: a 3.180 metros la densidad del aire es cerca de 28 \% menor que a nivel del mar, lo que reduce la capacidad de enfriamiento por aire, y la guía térmica de ASHRAE reduce la temperatura máxima admisible de entrada a los equipos en 1 °C por cada 300 metros sobre los 900 metros, es decir, en unos 7,6 °C en la faena (ASHRAE, 2021). La segunda es el polvo del rajo, que obliga a gabinetes sellados con filtros y a equipos de terreno con grado de protección IP65 o superior. La tercera es el frío de la madrugada, de hasta −6 °C, que reduce la autonomía de las baterías y endurece las pantallas táctiles. La cuarta es el uso con guantes y bajo sol directo, que descartó los dispositivos que la compañía instaló antes. La quinta es la distancia: la faena está a 145 km de Antofagasta, por lo que el compromiso de resolución de cuatro horas para el hardware crítico sólo se cumple con repuestos guardados en la propia faena. Cada implemento, con su cantidad, modelo, especificación y ubicación, se declara en el Formulario T-11. Sus fichas técnicas están en el Anexo A, y quién provee cada uno y la fuente de su especificación, en el Anexo C.

\subsubseccion{Implementos en la nube}
En la nube se privilegian servicios administrados de Azure en Chile Central y Brazil South. El cómputo, integración, datos, identidad, seguridad y observabilidad se contratan a nombre del cliente, sin exposición pública de los almacenes y con infraestructura definida como código.

La nube se contrata por consumo, sin hardware que reponer. La capacidad base de producción, que se usa las 24 horas, se reserva por tres años, y sólo el crecimiento en el cambio de turno y en los cierres mensuales se paga por uso. Cada recurso lleva etiquetas de ambiente, módulo y centro de costo, con presupuestos, alertas de desviación y un reporte mensual de consumo al cliente, conforme a las prácticas FinOps que exige el componente en nube.

\subsubseccion{Implementos en la faena}
En la faena se proveen la sala técnica, el clúster del Nodo y tres gabinetes de borde. El cuarto gabinete se instala en el acopio de puerto, en Antofagasta. Los equipos se seleccionan para altura, polvo, vibración, continuidad N+1 y reposición local. La red se estandariza para reducir la carga operacional del área de TI.

Todo el almacenamiento de la faena es de estado sólido, porque los discos mecánicos suelen certificarse sólo hasta unos 3.000 metros. Los servidores de la sala son equipos de rack estándar porque trabajan en un recinto climatizado, pero la orden de compra exige la certificación del fabricante para 3.180 metros o la reducción de temperatura que declare para esa altura. Los gabinetes, en cambio, trabajan con polvo y vibración, por lo que usan servidores de borde de chasis corto y rango térmico extendido, switches industriales sin ventiladores, todos con doble fuente. El gabinete de planta suma un par de cortafuegos industriales.

\subsubseccion{Implementos de terreno}
En terreno se agregan Botoneras de Pala y tabletas rugerizadas conectadas a la red LTE privada existente. Las cantidades se calculan por puestos simultáneos y consideran repuestos. La Figura \ref{fig:terreno-dispositivos-botonera-pala-tabletas-rugerizadas} muestra el alcance y los equipos que no se incorporan.

\begin{horizontal}
\csname ov@hoja\endcsname{330.2mm}{215.9mm}
\expandafter\def\csname ov@altofig\endcsname{0.88}
\figura{27.52}{fig-4-25-terreno-dispositivos-botonera-pala-tabletas}{Z08 · Terreno y dispositivos: Botonera de Pala, tabletas rugerizadas y red LTE privada existente.}{fig:terreno-dispositivos-botonera-pala-tabletas-rugerizadas}{Fuente: elaboración propia.}
\end{horizontal}
La figura muestra que en la zona Z08 se agregan sólo dos familias de dispositivos en los puestos de operación y que ambas llegan al Nodo de Faena por la red LTE privada existente, sin enlaces nuevos. La Botonera de Pala se dibuja a la izquierda porque es el primer punto donde nace el dato del mineral, en la cabina de la pala, y su flecha de eventos pasa por encima de las tabletas para llegar a la red LTE, porque no depende de ellas. Las tabletas envían formularios por la misma red. La red LTE entrega ambos flujos a la red de la faena y de ahí a la puerta operacional del Nodo, que queda fuera de la lámina, en la sala técnica. La nota inferior deja explícito lo que no se agrega: los camiones no reciben equipos porque el despacho de flota ya registra sus ciclos y su pesaje a bordo, y las perforadoras tampoco, porque los pozos de las siete que no tienen instrumentación se registran en tabletas del puesto.

Las cantidades se derivan de los puestos que registran y no de la dotación total. Se instalan seis Botoneras de Pala, una por pala y una de repuesto, operables con guantes, con lector de credencial y memoria para 48 horas de declaraciones. Las tabletas pertenecen al puesto y se comparten entre turnos, por lo que se compran 33 para los 30 puestos que registran al mismo tiempo en mina, mantenimiento, planta, seguridad y acopio de puerto, con un 10 \% de repuesto. Se eligió una tableta rugerizada, legible al sol, operable con guantes y con baterías intercambiables, porque esas condiciones descartaron los dispositivos que la compañía instaló antes (Dell Technologies, s.f.).

En portería, la Consola de Portería se ejecuta en dos PC industriales, uno por sentido de circulación. Los lectores faciales, sujetos al acuerdo con el sindicato, se dimensionan por el relevo de una detención mayor, que llega a unos 2.064 pasos por hora y exige cinco carriles, se habilitan cinco carriles más uno de reserva, con un lector por carril y uno de repuesto en el sitio, siete lectores en total. Se instalan además tres caudalímetros, más uno de repuesto, en la alimentación de agua de la planta, el campamento y la carga de camiones aljibe, y se entrega un estudio de cobertura de la red LTE del rajo.

El caso describe una red LTE privada instalada en 2021, con cobertura en la plataforma principal, las rampas superiores y los accesos, y zonas de sombra permanentes en el fondo del rajo, en el botadero norte y en dos sectores de rampa, sin reevaluación desde su instalación (CMAA, 2026b, cap. 11). Para SITCO esa red basta en las zonas cubiertas, porque cada Botonera envía eventos pequeños y guarda 48 horas sin señal. Lo que falta es cobertura en los frentes de carguío del fondo del rajo y en el botadero norte, donde nace la declaración del material y donde hoy el dato llega recién al recuperar señal, y capacidad para el cambio de turno, cuando se concentran los registros de todos los equipos. El estudio de cobertura (paquete 1.3.1.5, meses 2 a 4, con actualizaciones en los meses 8, 14 y 20) mide ambas cosas, proyecta las sombras a tres años según el avance del rajo (RT-03.24) e identifica qué celdas o repetidores adicionales y qué reubicaciones necesitaría el cliente. La ampliación de la red queda fuera del alcance de esta propuesta y se escala al cliente como decisión de inversión (riesgo R-08 del Formulario T-16).

\subsubseccion{Licencias, garantías y reposición}
El software de terceros que se licencia es acotado, porque la plataforma usa componentes de código abierto con soporte vigente y servicios administrados de la nube. Las suscripciones de GitHub Enterprise Cloud, GitHub Code Security, GitHub Secret Protection y SonarQube Cloud van a nombre del cliente (Art. 84.3). Las licencias propietarias restantes se limitan a Ubuntu Pro, Windows Server, PTC Kepware, Windows 11 IoT Enterprise y los servicios de Microsoft. Todo el equipamiento es nuevo, con cinco años de soporte 24x7 en sitio para servidores, red, seguridad y respaldo, que cubren los 56 meses del contrato sin renovación tecnológica.

\subseccion{Emplazamiento y Servicios en la Nube}
La decisión de dónde se ejecuta cada componente se toma con los seis criterios: latencia, criticidad operacional, volumen de datos, restricciones regulatorias, disponibilidad de conectividad y costo total de propiedad. Esta sección aplica estos criterios componente por componente y describe después los servicios de Azure que ejecutan la parte en nube, con su configuración.

\subsubseccion{Criterios de emplazamiento entre la faena y la nube}
Las funciones que requieren baja latencia, comportamiento determinista, continuidad durante cortes, procesamiento de alto volumen local o conexión directa con equipos físicos se ejecutan en la faena. Las que atienden usuarios externos, presentan demanda variable o pueden diferirse durante la pérdida del enlace se ejecutan en la nube. Debido a los cortes de conectividad registrados en Aranda, la continuidad operacional constituye el criterio principal de emplazamiento.

El Nodo de Faena y los gabinetes de borde mantienen los módulos, conectores, identidades y datos necesarios para sostener la operación local. El historiador se recibe inicialmente en la faena y se transmite comprimido por su volumen. El conector de radio y megafonía también permanece en el Nodo y se utiliza de manera síncrona únicamente para las alertas de evacuación. Su interfaz y protocolo se confirmarán durante el levantamiento técnico, sin considerar una nueva red de radiofrecuencia. El detalle de los componentes se presenta en las tablas \ref{tab:conectores-integrador-sistemas-etapa-lugar-ejecucion} y \ref{tab:modo-ventana-comportamiento-ante-falla-integraciones}, el Anexo B y el Formulario T-11.

Chile Central concentra las instancias maestras, los servicios compartidos, los portales, las integraciones externas y los almacenes autoritativos. Los módulos que funcionan exclusivamente en la nube mantienen en la faena las copias de lectura indispensables para la continuidad definida en 4.1.2.1.

Los datos personales y sensibles se alojan en Chile Central. Brazil South se utiliza solamente para recuperación ante desastres y no atiende solicitudes durante la operación normal.

La arquitectura limita los componentes instalados en la faena debido a sus exigencias de energía, climatización, mantenimiento y reposición. En la nube se privilegian servicios administrados, las excepciones y sus alternativas se justifican en el Anexo F. Los módulos conservan portabilidad mediante contenedores, infraestructura como código y contratos OpenAPI y AsyncAPI.

\subsubseccion{Servicios de Azure y su configuración}
La nube principal se presenta en las Figuras \ref{fig:nube-principal-chile-central-acceso-produccion} y \ref{fig:nube-principal-chile-central-operacion-produccion}. La Figura \ref{fig:nube-principal-chile-central-acceso-produccion} muestra por dónde entra cada población a Producción y la Figura \ref{fig:nube-principal-chile-central-operacion-produccion} muestra cómo se ejecutan los servicios y dónde se guardan los datos.

\begin{horizontal}
\figura{22.40}{fig-4-26-nube-principal-chile-central-acceso}{Z02 · Nube principal en Chile Central: acceso a producción.}{fig:nube-principal-chile-central-acceso-produccion}{Fuente: elaboración propia.}
\end{horizontal}
La figura muestra que la nube principal tiene dos caminos de entrada, uno para la población externa y otro para el personal propio, y que ninguno llega directamente a un almacén. En la fila superior, los contratistas, compradores e integraciones pasan primero por el DNS público, que resuelve los dominios de los portales hacia Front Door, y después por Front Door Premium, que aplica la red de distribución de contenidos, el WAF con reglas administradas y propias, la protección contra denegación de servicio en las capas 3, 4 y 7 del modelo OSI y la terminación TLS 1.3. El DNS va antes que Front Door porque ninguna dirección de la solución se publica directamente.

\begin{horizontal}
\figura{22.40}{fig-4-27-nube-principal-chile-central-operacion}{Z02 · Nube principal en Chile Central: operación de producción.}{fig:nube-principal-chile-central-operacion-produccion}{Fuente: elaboración propia.}
\end{horizontal}
La figura muestra el clúster AKS como núcleo de la operación en la nube. A él llegan el tráfico de producción y las colas de sincronización de la faena y del puerto. El Integrador de Sistemas conecta cada sistema externo por separado y entrega sus eventos a Service Bus Premium. Azure Monitor y Sentinel reciben la telemetría de los componentes. Los datos se distribuyen entre cinco almacenes: PostgreSQL para transacciones y auditoría, TimescaleDB para series de tiempo, ClickHouse para analítica, OpenSearch para búsqueda y Blob Storage para documentos y evidencia. OpenSearch se reconstruye desde PostgreSQL mediante interfaces versionadas.

Producción y los ambientes de Desarrollo, QA y Preproducción se ubican en Chile Central, Brazil South mantiene la recuperación. El tráfico público pasa por Azure DNS y Front Door Premium, con WAF en modo de prevención, protección contra bots y TLS 1.3. Las integraciones usan API Management Premium: Producción opera con dos unidades redundantes, una por zona y Brazil South mantiene una unidad para recuperación. En total se ofertan cinco unidades Premium. Desarrollo y QA comparten una instancia Developer. El personal propio usa Entra ID P2, contratistas, compradores y terceros usan Entra External ID. Entra Private Access permite acceder al Portal Web SITCO sin VPN mediante dos conectores.

El clúster AKS de producción distribuye sus nodos entre tres zonas: tres nodos D4s\_v5 para el sistema, entre cuatro y nueve D8s\_v5 para servicios, y tres E8s\_v5 para datos. Los servicios de negocio parten con tres réplicas, una por zona, y escalan antes de los cambios de turno. Las llamadas internas usan TLS mutuo. Las imágenes firmadas se guardan en Azure Container Registry, replicado en Brazil South, donde se mantiene un clúster mínimo que escala al activarse la recuperación.

Service Bus Premium replica los eventos a Brazil South con un retraso máximo de cinco minutos (Microsoft, s.f.-b). PostgreSQL tiene espera síncrona en otra zona, recuperación a un punto de los últimos 35 días y una réplica de lectura en Brazil South que se promueve durante la conmutación (Microsoft, s.f.-j). TimescaleDB, ClickHouse y OpenSearch operan en el grupo de datos, con réplicas entre zonas y discos ampliables. Blob Storage usa redundancia de zona, versionado e inmutabilidad según la retención donde los objetos también se replican a Brazil South. Redis se recrea desde el código en caso de recuperación porque almacena estado transitorio.

Cada región cuenta con una bóveda Key Vault Premium, claves en HSM y rotación de certificados cada 90 días. Log Analytics recibe la telemetría de la nube, la región de recuperación, el Nodo y los gabinetes para Sentinel conserva doce meses de eventos de seguridad en línea. La conexión con la faena usa Azure VPN Gateway activo-activo, con puertas redundantes en Chile Central y Brazil South. Desarrollo y QA operan con recursos de menor capacidad y no usan secretos ni datos reales de producción.

La red sigue una topología centralizada, con rangos separados por ambiente y región. La red central es 10.100.0.0/24, Producción usa 10.110.0.0/16, los ambientes no productivos, 10.120.0.0/16, y Brazil South, 10.130.0.0/16. En Producción, las subredes privadas se asignan por función: 10.110.1.0/27 para API Management, 10.110.2.0/24 para puntos de conexión privados, 10.110.3.0/27 para PostgreSQL y 10.110.16.0/20 para AKS. La exposición pública queda limitada a Front Door.

\subseccion{Arquitectura de Despliegue, Redes y Continuidad}
El despliegue se describe en cinco partes: los ambientes por los que avanza cada versión, los enlaces y la red de la sala técnica, los gabinetes de borde con sus zonas de seguridad, la alta disponibilidad frente a los puntos únicos de falla y la estrategia de respaldos. Cada parte responde la misma pregunta: qué pasa con la operación cuando falla una conexión o un equipo.

\subsubseccion{Ambientes de despliegue}
La solución opera cinco ambientes: Desarrollo, QA, Preproducción, Producción y Recuperación ante Desastres. Todos tienen la misma arquitectura, con los mismos servicios y contratos de interfaz, y se diferencian en sus datos, su escala y quién puede modificarlos. La Figura \ref{fig:entrega-software-cinco-ambientes-despliegue} muestra cómo una versión recorre los ambientes.

\begin{horizontal}
\figura{22.36}{fig-4-28-entrega-software-cinco-ambientes-despliegue}{Entrega de software y cinco ambientes de despliegue.}{fig:entrega-software-cinco-ambientes-despliegue}{Fuente: elaboración propia.}
\end{horizontal}
El código y sus contratos versionados producen una imagen firmada que pasa por Desarrollo, QA, Preproducción y Producción sin recompilarse. Cada ambiente usa el tipo de datos que le corresponde, y el paso a Producción requiere la aprobación registrada del Comité de Proyecto. Recuperación en Brazil South mantiene una réplica continua, no forma parte de la cadena de promoción.

Las herramientas de integración, análisis, pruebas y despliegue se describen en el Subdocumento 6, secciones 6.2.2.4 y 6.2.2.5, y en las Tablas \refext{T10}{tab:herramientas-evidencia-integracion} y \refext{T10}{tab:herramientas-evidencia-entrega} del Formulario T-10. Los ambientes tienen sus propias redes y secretos, sin conexión entre los no productivos y Producción. Desarrollo y QA comparten un clúster aislado que usa simuladores y datos versionados. Preproducción replica la configuración de Producción en tres zonas y permite probar hasta 300 transacciones por segundo. La sala técnica dispone de un banco de pruebas para el Nodo. Brazil South mantiene un clúster mínimo de tres nodos de servicios y tres de datos, que escala al activarse.

En la nube, la actualización usa despliegue azul-verde. Se descarta la actualización en el mismo entorno porque mezcla durante el cambio la versión vigente con la nueva y dificulta una reversión inmediata, el esquema azul-verde permite validar el entorno nuevo antes de conmutar el tráfico y volver al anterior si falla. En la faena, Azure Arc y Flux entregan la actualización por anillos, desde relaves hasta el Nodo, en la ventana de 03:00 a 05:00 y fuera de las restricciones operacionales indicadas. Se descarta el despliegue simultáneo en todos los sitios porque una falla afectaría a toda la operación, los anillos permiten validar cada grupo antes de avanzar al siguiente. La versión anterior queda disponible para revertir el cambio.

\subsubseccion{Enlaces y red de la sala técnica}
Las dos fibras cumplen la exigencia de caminos físicos y proveedores distintos, con ingreso al recinto por puntos separados. La fibra A la mantiene el proveedor actual por su ruta existente. La fibra B la contrata Overfore a otro proveedor, por un trazado distinto en todo su recorrido, incluido el camino privado de 26 km, y llega a la sala técnica por una canalización independiente. La diversidad se acredita con los planos de ruta de ambos proveedores y con una prueba de corte trimestral.

El tercer camino es un enlace satelital de órbita baja con un tercer proveedor, conectado al router SD-WAN de la fibra B. Mientras alguna fibra está disponible, el satélite sólo transporta las sondas que miden su estado. Si caen ambas fibras, el SD-WAN traslada el tráfico al satélite en 5 segundos o menos, con prioridad para la operación en vivo, las alertas, el tráfico interactivo y la bandeja de salida, las series del historiador se guardan en el Nodo y se envían al volver una fibra. Se exige al proveedor una capacidad de subida sostenida de al menos 20 Mbps, que cubre la punta de la faena, y una latencia de decenas de milisegundos, que se mide en la puesta en marcha. Con el satélite, la faena queda sin salida sólo si fallan a la vez los tres caminos, un escenario del orden de minutos al año si el satélite está disponible el 99 \% del tiempo. El satélite geoestacionario existente se mantiene sólo para voz y correo.

La Figura~\ref{fig:accesos-externos-enlaces-wan} muestra la conexión física de la faena con la nube y con los usuarios externos.

\begin{horizontal}
\csname ov@hoja\endcsname{330.2mm}{215.9mm}
\expandafter\def\csname ov@altofig\endcsname{0.88}
\figura{27.52}{fig-4-29-accesos-externos-enlaces-wan}{Z01 · Accesos externos y enlaces WAN.}{fig:accesos-externos-enlaces-wan}{Fuente: elaboración propia.}
\end{horizontal}
Las personas y los sistemas externos llegan a la nube sin pasar por la faena, y que la faena llega a la nube por tres caminos independientes. En la parte superior izquierda, el CIO y el personal propio usan el Portal Web SITCO por acceso de confianza cero, y los contratistas y compradores usan los portales externos a través del DNS público y el WAF. En la parte superior derecha, SAP corporativo y los sistemas de terceros que están fuera de la faena, como los portales de telemetría, el control de fatiga y los laboratorios externos, se integran directamente con la nube. En la parte inferior se dibujan los caminos de la faena. El satélite geoestacionario existente se conserva para voz y correo que sólo lleva, como último recurso, el correo de respaldo de las alertas de fatiga (4.1.6.3). La fibra actual pasa a ser el enlace A y la fibra adicional el enlace B, y el satélite de órbita baja es el enlace C, de respaldo. Los tres terminan en la sala técnica. La Figura \ref{fig:sala-tecnica-nodo-faena} muestra la sala técnica y el Nodo de Faena.

\begin{horizontal}
\csname ov@hoja\endcsname{330.2mm}{215.9mm}
\expandafter\def\csname ov@altofig\endcsname{0.88}
\figura{27.52}{fig-4-30-sala-tecnica-nodo-faena}{Z04 · Sala técnica y Nodo de Faena.}{fig:sala-tecnica-nodo-faena}{Fuente: elaboración propia.}
\end{horizontal}
La figura muestra el recorrido desde las fibras y el enlace satelital hasta el Nodo de Faena y los gabinetes. Cada fibra llega a un router Cisco Catalyst 8300 independiente. Ambos routers se conectan al par de firewalls Cisco Secure Firewall 1230 en alta disponibilidad, donde se aplican las reglas entre la faena y la nube y entre las zonas de la red local. Desde allí, un par de switches Cisco Catalyst 9500 distribuye el tráfico hacia los servidores del Nodo, los gabinetes por la fibra de faena y el segmento dedicado de SITCO en la red LTE del cliente.

La salida a Azure usa túneles IPsec con cifrado AES-256-GCM, enrutamiento dinámico y puertas VPN en modo activo-activo. SD-WAN mide la pérdida y la latencia de los caminos y conmuta al disponible en cinco segundos o menos. El enlace satelital entra por el router asociado a la fibra B.

El conector de radio y megafonía sale desde la sala técnica hacia el sistema existente y sólo transporta alertas de evacuación. El radioenlace de relaves es independiente: comunica su gabinete con la red de faena y reenvía las mediciones cuando vuelve el enlace.

La red de la faena usa 10.10.0.0/16, dividido en subredes para el Nodo (10.10.10.0/24), los gabinetes (10.10.20.0/24), SITCO en LTE (10.10.30.0/23), los sistemas integrados (10.10.40.0/24), la DMZ industrial (10.10.50.0/24) y la gestión fuera de banda (10.10.99.0/24). El gabinete del puerto usa 10.11.0.0/24. Los switches distribuyen la hora NTP sincronizada con Azure para cada evento registra la deriva de su reloj para ordenar los datos capturados durante un corte. Si fallan tres veces seguidas las sondas de los tres caminos, el Nodo registra el inicio del modo desconectado. La puerta operacional sigue respondiendo desde la faena, por lo que las personas pueden continuar registrando.

\subsubseccion{Gabinetes de borde y zonas de seguridad}
Los gabinetes protegen el dato en el punto donde nace cuando el sitio pierde su camino hacia el Nodo o hacia la nube. La Figura \ref{fig:gabinetes-borde-porteria-planta-relaves} muestra los tres gabinetes de la faena, de portería, planta y relaves. El cuarto gabinete está en el acopio de puerto, en Antofagasta.

\begin{horizontal}
\figura{22.44}{fig-4-31-gabinetes-borde-porteria-planta-relaves}{Z05 · Gabinetes de borde de portería, planta y relaves.}{fig:gabinetes-borde-porteria-planta-relaves}{Fuente: elaboración propia.}
\end{horizontal}
La Figura compara los gabinetes de portería, planta y relaves. Los tres reciben la conexión del Nodo mediante dos switches, que separan las redes locales, y tienen un agente en servidores activo-pasivo, una cola local en PostgreSQL y una UPS con al menos 30 minutos de autonomía. Los equipos y modelos se detallan en el T-11 y el Anexo A.

Cambian las fuentes de datos: portería envía los lectores y la consola al agente, que valida el acceso con la lista local en 1,5 segundos o menos, incluso sin enlace y relaves entrega mediciones ambientales y geotécnicas, que el agente conserva cuando se interrumpe el radioenlace. En relaves, ese radioenlace conecta el gabinete con la red de la faena. La Figura \ref{fig:planta-red-control-lectura-solo-sentido} muestra la lectura en un solo sentido desde la red de control de la planta.

\begin{horizontal}
\csname ov@hoja\endcsname{330.2mm}{215.9mm}
\expandafter\def\csname ov@altofig\endcsname{0.88}
\figura{27.52}{fig-4-32-planta-red-control-lectura-solo}{Z07 · Planta y red de control: lectura en un solo sentido.}{fig:planta-red-control-lectura-solo-sentido}{Fuente: elaboración propia.}
\end{horizontal}
La Figura \ref{fig:planta-red-control-lectura-solo-sentido} muestra una lectura unidireccional desde la red de control hacia SITCO, sin canal de retorno. El diodo de datos separa físicamente ambas redes: su transmisor está en control y su receptor en la zona desmilitarizada industrial (International Electrotechnical Commission, s.f.). El historiador espejo y los servidores Kepware reciben los datos mediante OPC UA restringido a lectura y suscripción. Los firewalls ISA3000 controlan el único conducto entre la zona desmilitarizada y el agente del gabinete de planta. El caudalímetro de alimentación de la planta también entrega datos por esta ruta; los del campamento y de carga de aljibes llegan al Nodo por la red de faena, porque están fuera de la red de control. La Figura \ref{fig:gabinete-borde-acopio-puerto} muestra el gabinete de borde del acopio de puerto.

\begin{horizontal}
\csname ov@hoja\endcsname{330.2mm}{215.9mm}
\expandafter\def\csname ov@altofig\endcsname{0.88}
\figura{27.52}{fig-4-33-gabinete-borde-acopio-puerto}{Z06 · Gabinete de borde del acopio de puerto.}{fig:gabinete-borde-acopio-puerto}{Fuente: elaboración propia.}
\end{horizontal}
En Z06, el gabinete del acopio de puerto se conecta directamente con Azure, sin pasar por la faena. Usa como camino principal la red del recinto portuario y, como respaldo, un servicio 4G de otro operador a través del router Cisco IR1101. Ambos caminos usan el túnel IPsec que el router establece por la conexión disponible. El switch local conecta los servidores de borde activo-pasivo, que registran la recepción, el acopio y el embarque, y mantienen una cola cuando se pierden ambos enlaces. Las tabletas leen por esa red los códigos QR de las guías. El gabinete usa la energía del recinto y tiene protección contra la corrosión salina. Las especificaciones de los equipos se detallan en el T-11 y sus anexos.

La red se divide en zonas y conductos según IEC 62443, con permisos limitados al flujo que necesita cada zona. Los dispositivos de terreno sólo acceden a la puerta operacional y la red de control sólo puede enviar datos a través del diodo. La gestión fuera de banda exige doble factor y sesiones grabadas, mientras la red corporativa accede al Portal Web SITCO por HTTPS. Esta separación limita el efecto de un dispositivo comprometido y evita que las zonas de menor confianza alcancen los almacenes o la red de control.

\subsubseccion{Alta disponibilidad y puntos únicos de falla}
La alta disponibilidad se diseña para la meta interna de 99,95 \% mensual de cada componente de infraestructura, superior al 99,9 \% del RNF-DIS-002, y de 99,9 \% de la transacción de negocio crítica de extremo a extremo. En la nube se obtiene por zona: los tres grupos de nodos de AKS, PostgreSQL con espera síncrona, Service Bus, Managed Redis, Blob Storage, API Management y la puerta VPN se reparten en al menos dos zonas, y la caída de una zona completa no interrumpe ningún servicio. En la faena se obtiene por duplicación: dos routers, dos firewalls, dos switches centrales, tres servidores del Nodo, dos servidores por gabinete, dos switches en cada gabinete de faena y uno en el de puerto, dos UPS con líneas A y B, dos unidades de climatización y dos generadores en N+1. En terreno se obtiene con el registro local de 48 horas de cada dispositivo y con los repuestos en la faena.

Subsisten cuatro puntos de falla común. El primero es la sala técnica como recinto: su pérdida detiene el Nodo de Faena. Se acota con un recinto protegido y equipos en N+1, si aun así se pierde, los dispositivos y gabinetes registran hasta 48 horas y el Nodo se reinstala de forma provisoria en la sala de servidores habilitada en 2014. El segundo son los enlaces, cuya caída corta la salida de la faena, se acota con dos trazados de fibra y el satélite, y la faena sigue en operación local hasta reconciliar. El tercero es la energía, la red LTE o el diodo de datos, cuya falla afecta un sitio o una fuente que se acota con UPS, generadores, registro local y repuestos. El cuarto es Chile Central o Entra ID, cuya caída deja sin portales o sin sesiones nuevas, se acota con la conmutación a Brazil South, las identidades locales y las cuentas de emergencia.

\subsubseccion{Estrategia de respaldos 3-2-1-1-0}
La estrategia mantiene tres copias en dos medios, una fuera del sitio, otra inmutable o fuera de línea y cero errores en las pruebas de restauración. Brazil South sirve para recuperación ante desastres que su réplica no cuenta como respaldo. En la nube se conservan producción en Chile Central, respaldos de servicio y una copia inmutable en Brazil South, con autorización múltiple y claves separadas (Microsoft, s.f.-d). En la faena, los datos del Nodo se replican entre tres servidores y se respaldan en la sala de 2014: incrementales cada hora y copias diarias durante 30 días. Las imágenes y configuraciones también se copian semanalmente a cartuchos cifrados fuera de línea.

La retención depende del dato: los respaldos transaccionales y de series incluyen copias diarias inmutables por doce meses y mensuales por cinco o diez años, según su uso, los tributarios se conservan siete años en SAP. Al quinto año, restaurar los datos transaccionales estimados toma una hora y las series, dos horas a 400 MB/s. Los documentos se restauran por versión, la analítica y OpenSearch se reconstruyen. La infraestructura se recupera desde el código y las imágenes firmadas. Cada mes se prueba una restauración en Preproducción, se verifica la integridad y se registran tiempos y correcciones.

\subseccion{Dimensionamiento y Plan de Capacidad}
El dimensionamiento convierte la volumetría operacional del caso en carga sobre cada recurso, con supuestos declarados. La carga transaccional de referencia es de 40 TPS en régimen y 200 TPS en la punta, la aceptación se verifica con una prueba de 300 TPS y la arquitectura crece hasta 600 TPS sin cambiar su topología. Las lecturas continuas de sensores se dimensionan por separado, porque forman un flujo de ingesta y no transacciones. El cálculo de cada cifra de esta sección se presenta en el Anexo D, y la capacidad que resulta para cada implemento, en el Formulario T-11.

\subsubseccion{Volumetría operacional}
La operación registra unos 8.100 hechos de negocio diarios, entre viajes, ciclos de carguío, pozos, muestras, despachos a puerto, órdenes de trabajo, registros de seguridad, alertas de fatiga y pasos de portería, que generan unos 24.300 eventos de dominio. Las series de tiempo son de otro orden de magnitud. Las 5.000 señales del historiador que SITCO incorpora, la telemetría de 52 equipos y los 189 instrumentos y caudalímetros suman unos 94 millones de lecturas diarias, cerca de 1.090 por segundo.

El criterio de todo el dimensionamiento. Los hechos de negocio, menos de uno por segundo, no fijan la capacidad, la fijan la concurrencia del cambio de turno y el volumen de las series. La proyección del caso a tres años es moderada, con 7 \% más viajes, 10 \% más etiquetas, 35 \% más equipos con telemetría y 24 \% más instrumentos, y queda cubierta por el factor de diseño de tres veces.

\subsubseccion{Usuarios, concurrencia y transacciones}
SITCO se dimensiona para 8.500 personas con acceso en una detención mayor de planta, con 245 personas concurrentes en régimen, que actúan cada 60 segundos, y una referencia de 505 personas concurrentes y 37 dispositivos de terreno en el relevo de las 08:00, el de mayor carga.

En el relevo, con una acción cada 30 segundos y seis llamadas por acción, las 505 personas producen unas 101 TPS, que con las 34 de tabletas, botoneras, consolas, portería, cierre de turno y conectores suman unas 135 TPS. Se adopta 200 TPS como punta declarada, un 48 \% sobre las 135 TPS calculadas, para absorber las simultaneidades que una estimación promedio no representa, como el cierre mensual de conciliación. La detención de marzo y septiembre, cuando la dotación llega a 8.500 personas (Compañía Minera Altos de Aranda [CMAA], 2026b, cap. 14, p. 23), no duplica a los 505 usuarios concurrentes: las cerca de 4.300 personas adicionales (8.500 − 4.180) son trabajadores de contratistas que llegan a SITCO como pasos de portería, carga que se prueba aparte (Subdocumento 5, sección 5.4.6). Aun si se duplicaran los 90 administradores de contratistas, la carga subiría en 18 TPS, a unas 153 TPS, dentro de la punta y de la prueba de 300 TPS. La Tabla \ref{tab:volumetria-referencia-capacidad-diseno} resume la volumetría de referencia y la capacidad de diseño.

\begin{tabla}{I{0.235}I{0.223}I{0.542}}{Volumetría de referencia y capacidad de diseño.}{tab:volumetria-referencia-capacidad-diseno}
  {\cab{Magnitud} & \cab{Valor de diseño} & \cab{Cálculo o criterio}}
  {Fuente: elaboración propia, cálculo detallado en el Anexo D de Overfore-Subdocumento4-Anexos.}
Carga transaccional estimada en régimen & 40 TPS & 245 personas × 6 llamadas ÷ 60 s = 24,5 TPS, más 9 de conectores y 6,5 de terreno, portería y eventos \\
Punta transaccional declarada & 200 TPS & Reserva sobre las 135 TPS calculadas para cambios de turno y detenciones \\
Prueba de carga & 300 TPS & 1,5 veces la punta, con ingesta de series simultánea \\
Capacidad de crecimiento & 600 TPS & Tres veces la punta declarada mediante escalado \\
Usuarios y dispositivos concurrentes & 505 personas y 37 dispositivos & 30 tabletas, 5 botoneras y 2 consolas en servicio \\
Hechos y eventos operacionales & 8.100 hechos y 24.300 eventos diarios & Promedio de tres eventos por hecho \\
Señales continuas del historiador & 5.000 cada 5 s & 86,4 millones de muestras diarias de 64 bytes \\
Series totales & 94,16 millones de lecturas diarias & Historiador, 52 equipos y 189 instrumentos o caudalímetros \\
Series almacenadas & 2,19 TB al año sin comprimir & 64 bytes por muestra del historiador y 60 bytes por otra lectura \\
Series a cinco años & 2,8 TB lógicos & Compresión de diseño 4:1 y siete días recientes sin comprimir \\
Corte de 48 horas & 188,3 millones de lecturas & 12,0 GB sin comprimir o 3,0 GB con compresión 4:1 \\
Reconciliación en cuatro horas & 13.100 lecturas por segundo & Incluye todas las muestras diferidas, flujo corriente en paralelo \\
\end{tabla}
La volumetría muestra que la carga transaccional y la ingesta de series son flujos distintos, por lo que el almacenamiento y la sincronización se dimensionan para las series sin restar capacidad transaccional al relevo.

\subsubseccion{Capacidad de cómputo en la nube y en la faena}
En la nube, los nueve servicios de negocio parten con tres réplicas cada uno, una por zona, que con el supuesto de 25 TPS por vCPU cubren la prueba de 300 TPS. Con los portales, los conectores y la malla de servicios, el grupo solicita unas 21 vCPU, que caben en cuatro nodos D8s\_v5, para 600 TPS se proyectan al menos seis nodos, dentro del máximo configurado de nueve. A 200 TPS, PostgreSQL procesa unas 1.000 sentencias por segundo, y la prueba de carga confirmará si el tamaño D4ds\_v5 se mantiene.

En la faena, el Nodo se dimensiona con un servidor fuera de servicio. Su carga solicita unos 189 GB, el 37 \% de la memoria de los dos servidores restantes, y llega a 56 \% con tres veces la volumetría, porque sólo crece la memoria de PostgreSQL, que se duplica (Anexo D, D.3.3), bajo el umbral de ampliación de 60 \%. Durante un corte, la punta local de unas 120 TPS usa cerca del 42 \% de sus hilos de procesador, lo que se verificará con la prueba de operación aislada de 48 horas.

\subsubseccion{Almacenamiento a cinco años}
Se generan unos 6 GB diarios sin comprimir, que con la compresión de diseño de 4:1 llegan a unos 2,8 TB al quinto año, frente a 71 GB del almacén transaccional, 0,68 TB de documentos y evidencia, y 0,75 TB entre analítica e índice. Una copia completa de SITCO ocupa así unos 4,3 TB al quinto año, y el volumen histórico a migrar es de unos 205 GB, casi todo formularios de seguridad escaneados.

Cada nodo del grupo de datos guarda una réplica de las series, la analítica y el índice, unos 3,5 TB al quinto año. Por eso los discos iniciales de 3 TB se amplían en línea a 6 TB antes de alcanzar el 80 \% de ocupación, y hasta 14 TB si se materializa el crecimiento de tres veces. La compresión efectiva se medirá con datos reales antes de fijar la capacidad contratada. En la faena, cada servidor del Nodo guarda unos 295 GB, cifra que se confirmará en la prueba de corte antes de fijar la alarma de capacidad.

\subsubseccion{Ancho de banda y sincronización después de un corte}
La faena requiere 18 Mbps en punta y 8 Mbps en régimen. Con tres veces la volumetría, la demanda de diseño llega a 54 Mbps, capacidad que puede sostener una sola fibra de 100 Mbps. Si fallan ambas fibras, el enlace satelital de 20 Mbps mantiene la conexión. Los caudales de los demás sitios y su cálculo están en el Anexo D.3.5, Tabla \refext{A4}{tab:ancho-banda-sincronizacion-tras-corte}, las capacidades de los enlaces se especifican en el T-11.

El Nodo almacena los datos de un corte de hasta 48 horas en una cola de 150.000 eventos, con alerta al 80 \%. Al restablecerse la conexión, la sincronización debe completarse en cuatro horas, con caudal reservado y prioridad para el tráfico en vivo, el envío de eventos acumulados se limita a 40 por segundo. Los supuestos y cálculos de reconciliación están en la Tabla \refext{A4}{tab:ancho-banda-sincronizacion-tras-corte}. Este comportamiento se comprobará con una prueba de desconexión y reconexión de 48 horas.

\subsubseccion{Crecimiento y umbrales de ampliación}
El primer recurso que alcanza su límite es el disco del grupo de datos, por el crecimiento de las series, y su ampliación en línea se activa cuando la proyección a seis meses supera el 80 \% de ocupación. El grupo de servicios crece de cuatro a seis o más nodos según el rendimiento medido, y PostgreSQL aumenta de tamaño si la prueba de 300 TPS muestra saturación. El Nodo recibe más memoria o un cuarto servidor al superar el 60 \% de uso con un servidor fuera de servicio, la segunda fibra admite ampliar su capacidad por contrato y el servicio satelital admite un plan de mayor capacidad sin cambiar el terminal. Una unidad operacional adicional agrega su propio Nodo y sus gabinetes y comparte los servicios de la nube.

\seccion{Data center}
SITCO Aranda distribuye sus componentes entre tres emplazamientos. La sala técnica de la faena aloja el Nodo de Faena y mantiene la operación local cuando se interrumpe el enlace con Antofagasta. Azure Chile Central aloja los nueve módulos. Azure Brazil South recibe la réplica y permite recuperar los servicios críticos si falla la región primaria. La sala de servidores de 2014 se destina al personal de soporte, el enrolamiento y la custodia del respaldo local, los cuatro gabinetes de borde se ubican donde se originan o protegen los datos.

El CIO no se habilita como centro de datos secundario porque depende del mismo enlace de fibra que la solución debe tolerar. Los sitios pueden cumplir un RTO de cuatro horas y un RPO de quince minutos, pero la recuperación en la nube usa procedimientos automatizados. También evita que el equipo de TI, compuesto por once personas, tenga que operar una segunda sala. Aunque el CIO permitiría mantener los datos en Chile, esa ventaja no compensa las diferencias de continuidad, operación y costo, la transferencia internacional a Brazil South requiere la aprobación correspondiente. El CIO conserva su función como sitio de usuarios: sus 24 posiciones acceden a SITCO y, durante un corte, coordinan con la faena por radio y teléfono satelital.

\subseccion{Especificaciones Data Center Primaria}
El centro de datos primario tiene dos partes complementarias. En la nube, el proveedor es Microsoft Azure y la región primaria es Chile Central, en Santiago (Microsoft, s.f.-f), en las que se reparten los servicios de producción de la plataforma. En la faena, el sitio LOCAL es la sala técnica principal que exige el requisito RT-06.01, que obliga a habilitar o reemplazar la sala de servidores de 40 m² habilitada en 2014 y descrita en el capítulo 6 del caso (CMAA, 2026b, cap. 6, p. 11), porque no cumple el estándar de un recinto técnico: tiene climatización simple, alimentación ininterrumpida de 15 minutos y control de acceso por llave. La solución la reemplaza por un módulo nuevo en lugar de remodelarla, porque remodelarla obligaría a trasladar temporalmente lo que hoy funciona en ella. El módulo prefabricado lo provee e instala Overfore (Formulario T-11, T11-29), la obra civil de base, la acometida eléctrica y las canalizaciones son de cargo del cliente (SUP-34), y su especificación técnica y su coordinación son de cargo de Overfore.

\subsubseccion{Recinto y distribución interna}
La sala técnica es un módulo prefabricado de 8x5 metros, 40 m², que llega armado de fábrica y se instala junto al edificio de administración, donde ya llega la fibra actual. Es de uso exclusivo de la solución y tiene acceso independiente del resto de las dependencias. Sus muros son paneles de acero con núcleo de lana mineral, con resistencia al fuego de 120 minutos, y su piso técnico elevado de 30 cm aloja bandejas separadas para energía y datos. La Figura \ref{fig:plano-distribucion-sala-tecnica-principal} presenta el plano de distribución interna.

\begin{horizontal}
\csname ov@hoja\endcsname{330.2mm}{215.9mm}
\expandafter\def\csname ov@altofig\endcsname{0.88}
\figura{27.52}{fig-4-34-plano-distribucion-sala-tecnica-principal}{Z04 · Plano de distribución de la sala técnica principal.}{fig:plano-distribucion-sala-tecnica-principal}{Fuente: elaboración propia.}
\end{horizontal}
La sala técnica nueva es un módulo exclusivo junto al edificio de administración, con acceso independiente, envolvente resistente al fuego, separación de energía y datos, y respaldo alojado en la sala de 2014. La Figura \ref{fig:plano-distribucion-sala-tecnica-principal} muestra la distribución y capacidad de crecimiento.

Los 3 gabinetes rack de 42U se identifican como RK-1 a RK-3. RK-1 aloja comunicaciones y utiliza el 36 \% de su capacidad, RK-2 aloja los servidores y utiliza el 24 \%, con espacio para un cuarto servidor del Nodo. RK-3 se instalan con PDU A y B y quedan disponibles para crecimiento y segregación futura de equipos, sin alterar la topología de la sala. El cableado estructurado se instala conforme a ANSI/TIA-568 e ISO/IEC 11801, se certifica enlace por enlace y su informe se entrega al cliente.

El ingreso se controla con verificación facial de cada persona autorizada que ingresa (RT-06.20), con su consentimiento y con huella dactilar (AFIS) como respaldo, y con bitácora de cada ingreso y egreso. Como sólo alcanza al personal de la sala, su aplicación no depende del acuerdo con el sindicato, que condiciona sólo la verificación facial de portería (SUP-13). La sala de 2014 aloja las funciones que no deben estar junto a los equipos, la custodia de los respaldos y el puesto del personal de soporte, de modo que la operación habitual no requiere entrar a la sala técnica.

\subsubseccion{Climatización a 3.180 metros}
Para la parte de climatización se tienen dos unidades de precisión en N+1, cada una de las cuales cubre la carga corregida por altura. El suministro, humedad, sobrepresión y detección de agua se monitorean desde la plataforma de observabilidad.

La altura también reduce la temperatura admisible: la guía de ASHRAE admite hasta 32 °C de entrada en la clase A1, y a esta altura el máximo baja a unos 24,4 °C (ASHRAE, 2021). Por eso el suministro del pasillo frío se fija en 21 °C, con alerta a 24 °C. Como el aire del altiplano es muy seco, las unidades incluyen humidificador para evitar descargas electrostáticas, y el recinto se mantiene en leve sobrepresión para que el polvo del rajo no ingrese. El pasillo frío confinado permite una eficiencia en el uso de la energía (PUE) estimada de 1,6, y la temperatura, la humedad y la presencia de agua se integran a la misma plataforma de observabilidad que la nube.

\subsubseccion{Energía}
La carga de TI es de 3,65 kW y, con un PUE de 1,6, la sala consume 5,84 kW. Se instalan dos UPS de doble conversión de 10 kVA en una configuración en que cada una puede sostener por sí sola la carga de TI, con margen de crecimiento y ajuste por altura. Su autonomía se dimensiona en 30 minutos a plena carga, el mínimo obligatorio de RT-06.07 de las Bases técnicas transversales (CMAA, 2026c, cap. 6), que cubre el arranque y la toma de carga de los generadores; las UPS de los gabinetes de borde tienen 60 minutos, el doble de ese mínimo. Sus baterías de litio entregan 30 minutos de autonomía, suficientes para cubrir el arranque de los generadores o apagar los equipos de forma ordenada si ambos fallan.

Los generadores también alimentan los gabinetes de portería y planta, para una carga total de 7,34 kW. Se instalan dos generadores diésel en N+1, dimensionados para esa carga. Como la faena está a 145 km de Antofagasta y la ruta puede quedar cortada entre uno y tres días en invierno, el estanque de doble pared es de 1.000 litros. Esa capacidad entrega unos 19 días de autonomía a la carga actual, con recargas cada 9,5 días al llegar a la mitad. Los generadores se prueban mensualmente bajo carga. La Figura \ref{fig:energia-respaldo-sala-tecnica-gabinetes} muestra la energía de respaldo de la sala técnica y de los gabinetes.

\begin{horizontal}
\figura{22.32}{fig-4-35-energia-respaldo-sala-tecnica-gabinetes}{Z09 · Energía de respaldo de la sala técnica y de los gabinetes.}{fig:energia-respaldo-sala-tecnica-gabinetes}{Fuente: elaboración propia.}
\end{horizontal}
La figura anterior muestra que la energía sigue un orden fijo, desde las fuentes hasta los equipos protegidos, y que ningún equipo de TI depende de una sola línea. A la izquierda están las dos fuentes: la red eléctrica principal y los generadores A y B, alimentados por el estanque de combustible. Ambas llegan al tablero de transferencia automática, que es el único punto donde se decide de qué fuente se alimenta la sala, por eso se dibuja entre las fuentes y las cargas. Desde el tablero salen dos líneas independientes, A y B, hacia las UPS A y B de la sala, y cada UPS alimenta una de las dos fuentes de cada equipo, de modo que la falla de una UPS, de una línea o de una fuente de poder no apaga ningún equipo. Desde el mismo tablero sale la distribución hacia las UPS locales de los gabinetes de portería y planta; los gabinetes de relaves y de puerto se alimentan de su red local, porque están a 6 km y a 145 km de la sala, respectivamente.

La instalación eléctrica del recinto es independiente de la del edificio de administración, con circuitos separados para las líneas A y B, la climatización y los gabinetes, y cumple la normativa eléctrica chilena vigente, incluida la NCh Elec. 2777 de puesta a tierra. Cada semestre se revisa y se mide la instalación, con un informe entregable al cliente.

\subsubseccion{Detección y extinción de incendios}
La detección temprana usa aspiración de aire con tecnología láser sobre el pasillo caliente y la zona de energía, con dos niveles de alarma. El primero avisa a la consola operativa y el segundo, confirmado por un detector de otra zona, inicia la extinción. Así el polvo del rajo no dispara el agente, pero el calentamiento de un cable o de una batería se detecta antes de que haya llama.

La extinción es automática, con agente limpio FM-200 conforme a NFPA 2001 (National Fire Protection Association, 2022). A 3.180 metros la menor presión hace que el agente se expanda más, por lo que la carga se corrige por altura y se reduce a unos 30 kg en un solo cilindro. La descarga detiene la climatización, cierra las compuertas y espera 30 segundos con alarma, y una prueba de hermeticidad verifica que la concentración se mantenga diez minutos. Se complementa con extintores portátiles de CO₂, y toda alarma o descarga se notifica a la consola operativa y al cliente, con una prueba semestral del sistema.

\subseccion{Especificaciones Data Center Secundario}
El sitio secundario es la región Azure Brazil South, en São Paulo, en modalidad activo-pasivo. Recibe una réplica continua de los datos de producción, mantiene una capacidad mínima en espera y asume los servicios críticos si la región primaria deja de estar disponible. Una región secundaria dentro de Santiago no cubriría el evento de fuerza mayor que justifica un sitio secundario, porque compartiría con la primaria la zona sísmica, la red eléctrica y el entorno urbano.

La elección se registra en dos decisiones. ADR-09, región de recuperación: el sitio secundario debe cubrir la fuerza mayor y restablecer los servicios críticos en cuatro horas con una pérdida máxima de 15 minutos de datos, y Chile Central es la única región de Azure en Chile. Se evaluaron tres alternativas: (a) sólo las tres zonas de disponibilidad de Chile Central, que protegen contra la falla de un centro de datos pero comparten la región, su zona sísmica y su red eléctrica; (b) el CIO de Antofagasta como sitio propio, que se descarta en ADR-10, y (c) Brazil South en activo-pasivo, que se adopta por ser la región general de Azure más cercana fuera de esos riesgos comunes. Sus consecuencias son dos: la réplica de datos sensibles fuera de Chile requiere la autorización escrita del cliente (Ley N° 21.719; riesgo R-16 del Formulario T-16), con una adaptación condicionada de 400 HH si no se obtiene, y el desfase de la réplica asíncrona debe quedar dentro de los 15 minutos, lo que se mide en la prueba de recuperación del paquete 1.5.5.1 antes del paso a producción afectado. ADR-10, el CIO no es sitio de recuperación: el CIO es un sitio de usuarios con 24 posiciones, sin recinto, climatización ni energía de respaldo para servidores, comparte con la faena la red eléctrica y la zona sísmica del norte, y convertirlo en sitio secundario obligaría a construir y operar 24x7 un segundo centro de datos. Por eso conserva su función de usuarios y, durante un corte, coordina con la faena por radio y teléfono satelital. Ambas decisiones tienen fecha 4 de noviembre de 2026 y están aceptadas; un cambio posterior se registra como una decisión nueva.

\subsubseccion{Región, modalidad y replicación}
La modalidad activo-pasivo se elige frente al costo, al objetivo de recuperación y a la complejidad operacional. Una modalidad activo-activo obligaría a escribir el mismo dato en dos regiones a la vez, lo que contradice la regla de un solo sitio escritor por agregado que ordena la arquitectura, y duplicaría el cómputo en forma permanente. La modalidad activo-pasivo mantiene en espera sólo lo necesario para cumplir cuatro horas: un clúster AKS de tres nodos de servicios y tres de datos que escala al activarse, una unidad de API Management, la puerta VPN con los túneles de la faena preconfigurados, el registro de imágenes replicado y los almacenes que reciben la réplica.

La distancia entre ambas regiones es de unos 2.600 km, y las amenazas comunes analizadas no las afectan a la vez. Un terremoto en la zona de subducción chilena no alcanza a São Paulo, que está sobre una placa estable; una falla eléctrica extendida del sistema chileno no afecta al sistema brasileño, y la cordillera de los Andes separa ambos entornos. La falla del propio proveedor en una región no se propaga a la otra, porque cada región tiene su plano de control y la réplica viaja por la red troncal del proveedor.

La Figura \ref{fig:recuperacion-brazil-south} presenta la recuperación en Brazil South (zona Z03).

\figura{16.05}{fig-4-36-recuperacion-brazil-south}{Z03 · Recuperación en Brazil South.}{fig:recuperacion-brazil-south}{Fuente: elaboración propia.}
La Figura \ref{fig:recuperacion-brazil-south} separa la protección de los datos de la operación durante una contingencia. En la franja superior, la réplica continua protege las transacciones, las series de tiempo, los eventos y los objetos. Una ruta independiente lleva los respaldos a la copia inmutable fuera del sitio. En la franja inferior se muestra Brazil South, que activa sus servicios críticos después de la conmutación, usa los datos recuperados y envía telemetría al sistema de monitoreo. La analítica y OpenSearch se reconstruyen una vez recuperados los servicios principales. El procedimiento también coordina el retorno a Chile Central y la restauración desde el respaldo si la réplica contiene datos corruptos.

El retraso de las réplicas se monitorea frente al RPO de 15 minutos establecido en la sección 4.3.2.2. PostgreSQL suele presentar retrasos de segundos (Microsoft, s.f.-j) Service Bus tiene un máximo configurado de cinco minutos (Microsoft, s.f.-b).

Los datos personales se alojan en Chile Central. Su réplica en Brazil South constituye una transferencia internacional, sujeta a la aprobación del cliente y a cláusulas contractuales de protección de datos. Los almacenes usan cifrado con claves bajo control del cliente, conforme a la Ley N.º 21.719 (CMAA, 2026a, p. 11, Ley N.º 21.719, 2024).

\subsubseccion{Conmutación, retorno y pruebas}
Para los servicios críticos, la recuperación en Brazil South tiene un RPO máximo de 15 minutos y un RTO de 4 horas, (Subdocumento 3, sección 3.2.2.6). La operación de la faena se mantiene en el Nodo de Faena durante la contingencia. La región secundaria, la réplica y los recursos de recuperación se especifican en el Formulario T-11, ítem T11-02, y en las fichas correspondientes del Anexo A.

La conmutación procede si el proveedor declara indisponible la región primaria o si Front Door no alcanza su origen por más de 30 minutos. El jefe del incidente confirma la decisión, después, el procedimiento documentado activa los recursos de Brazil South, cambia el tráfico y valida los servicios antes de reabrirlos. El retorno a Chile Central se programa entre las 03:00 y las 05:00, fuera del cierre contable y de las 48 horas previas a un embarque, incluye reconciliar los datos y verificar su integridad.

La conmutación se prueba dos veces al año durante las detenciones de planta de marzo y septiembre, según el ítem T11-02. Cada prueba registra los tiempos y la pérdida de datos observados, junto con las correcciones necesarias. Esta prueba de continuidad complementa la restauración mensual de respaldos descrita en 4.2.3.5.

\seccionsinnumero{Referencias}
\begin{referencias}
\item Angular. (s.f.). Versioning and releases. \href{https://angular.dev/reference/releases}{https://angular.dev/reference/release}s
\item ASHRAE. (2021). Thermal guidelines for data processing environments (5.ª ed.). ASHRAE.
\item Canonical. (s.f.). Ubuntu release cycle. https://ubuntu.com/about/release-cycle
\item ClickHouse. (s.f.). Which ClickHouse version to use in production? https://clickhouse.com/docs/faq/operations/production
\item Compañía Minera Altos de Aranda. (2026a). \emph{Bases administrativas}. Licitación N° TFEP-01/2026.
\item Compañía Minera Altos de Aranda. (2026b). \emph{Bases técnicas del caso Minería}. Licitación N° TFEP-01/2026.
\item Compañía Minera Altos de Aranda. (2026c). \emph{Bases técnicas transversales}. Licitación N° TFEP-01/2026.
\item Dell Technologies. (s.f.). Dell Latitude 7030 Rugged Extreme Tablet. https://www.dell.com/en-us/shop/dell-laptops/latitude-7030-rugged-extreme-tablet/spd/latitude-10-7030-rugged-tablet
\item Eclipse Adoptium. (s.f.). Support. https://adoptium.net/support
\item Envoy Project. (s.f.). Release schedule. https://github.com/envoyproxy/envoy/blob/main/RELEASES.md
\item Evans, E. (2003). Domain-driven design: Tackling complexity in the heart of software. Addison-Wesley.
\item Grafana Labs. (s.f.). Strategies for upgrading your self-managed Grafana instance. https://grafana.com/docs/grafana/latest/upgrade-guide/when-to-upgrade/
\item International Electrotechnical Commission. (s.f.). IEC 62443 background. https://syc-se.iec.ch/?p=4375
\item International Organization for Standardization. (2020). ISO 15143-3:2020 Earth-moving machinery and mobile road machinery. Worksite data exchange. Part 3: Telematics data. ISO.
\item International Organization for Standardization. (2022). ISO/IEC/IEEE 42010:2022 Software, systems and enterprise — Architecture description. ISO.
\item International Society of Automation. (2010). ANSI/ISA-95.00.01-2010 (IEC 62264-1 Mod) Enterprise-control system integration. Part 1: Models and terminology. ISA.
\item Kubernetes. (s.f.). Releases. https://kubernetes.io/releases/
\item Ley N° 19.799. (2002). Sobre documentos electrónicos, firma electrónica y servicios de certificación de dicha firma. Diario Oficial de la República de Chile.
\item Ley N° 21.719. (2024). Regula la protección y el tratamiento de los datos personales y crea la Agencia de Protección de Datos Personales. Diario Oficial de la República de Chile.
\item Martin, R. C. (2017). Clean architecture: A craftsman's guide to software structure and design. Prentice Hall.
\item Microsoft. (s.f.-a). Azure Cache for Redis retirement: Frequently asked questions. https://learn.microsoft.com/en-us/azure/azure-cache-for-redis/retirement-faq
\item Microsoft. (s.f.-b). Azure Service Bus Geo-Replication. https://learn.microsoft.com/en-us/azure/service-bus-messaging/service-bus-geo-replication
\item Microsoft. (s.f.-c). Distribute your app and the WebView2 Runtime. https://learn.microsoft.com/en-us/microsoft-edge/webview2/concepts/distribution
\item Microsoft. (s.f.-d). Immutable vault for Azure Backup. https://learn.microsoft.com/en-us/azure/backup/backup-azure-immutable-vault-concept
\item Microsoft. (s.f.-e). Istio-based service mesh add-on for Azure Kubernetes Service support policy. https://learn.microsoft.com/en-us/azure/aks/istio-support-policy
\item Microsoft. (s.f.-f). List of Azure regions. https://learn.microsoft.com/en-us/azure/reliability/regions-list
\item Microsoft. (s.f.-g). Message sessions. Azure Service Bus. https://learn.microsoft.com/en-us/azure/service-bus-messaging/message-sessions
\item Microsoft. (s.f.-h). .NET and .NET Core support policy. https://dotnet.microsoft.com/en-us/platform/support/policy/dotnet-core
\item Microsoft. (s.f.-i). Operating systems and browsers supported by Microsoft Intune. https://learn.microsoft.com/en-us/intune/fundamentals/ref-supported-platforms
\item Microsoft. (s.f.-j). Read replicas in Azure Database for PostgreSQL. https://learn.microsoft.com/en-us/azure/postgresql/flexible-server/concepts-read-replicas
\item Microsoft. (s.f.-k). Windows 11 IoT Enterprise LTSC 2024. Microsoft Lifecycle. https://learn.microsoft.com/en-us/lifecycle/products/windows-11-iot-enterprise-ltsc-2024
\item Microsoft. (s.f.-l). Windows Server 2025. Microsoft Lifecycle. https://learn.microsoft.com/en-us/lifecycle/products/windows-server-2025
\item National Fire Protection Association. (2022). NFPA 2001: Standard on clean agent fire extinguishing systems. NFPA.
\item Nygard, M. T. (2018). Release it! Design and deploy production-ready software (2.ª ed.). Pragmatic Bookshelf.
\item OPC Foundation. (s.f.). Unified Architecture. https://opcfoundation.org/about/opc-technologies/opc-ua/
\item OpenSearch Project. (s.f.). Release schedule and maintenance policy. https://opensearch.org/releases/
\item PostgreSQL Global Development Group. (s.f.). Versioning policy. https://www.postgresql.org/support/versioning/
\item Prometheus. (s.f.). Long-term support. https://prometheus.io/docs/introduction/release-cycle/
\item PTC. (s.f.). Support and subscription renewal: Kepware. https://www.ptc.com/en/products/kepware/standard-support
\item Richardson, C. (2018). Microservices patterns: With examples in Java. Manning.
\item Rose, S., Borchert, O., Mitchell, S. y Connelly, S. (2020). Zero trust architecture (NIST Special Publication 800-207). National Institute of Standards and Technology. https://doi.org/10.6028/NIST.SP.800-207
\item Spring. (s.f.). Spring Boot: Support. https://spring.io/projects/spring-boot\#support
\item Zetetic. (s.f.). SQLCipher. https://www.zetetic.net/sqlcipher/
\end{referencias}

\seccionsinnumero{Declaración de uso de IA}
Se utilizó IA generativa, principalmente como apoyo en la redacción, en la organización del contenido y en los diagramas; el detalle por sección está en la Tabla \ref{tab:declaracion-uso-ia-seccion}.

\begin{tabla}{I{0.167}I{0.149}I{0.273}I{0.113}I{0.131}I{0.167}}{Declaración de uso de IA por sección.}{tab:declaracion-uso-ia-seccion}
  {\cab{Sección} & \cab{Herramienta} & \cab{Finalidad del uso} & \cab{Nivel en\newline texto} & \cab{Nivel en\newline diagramas} & \cab{Revisión humana\newline (Quien y que verificó)}}
  {Fuente: elaboración propia.}
Introducción y capítulo 4 & Claude & Redacción de enlace entre arquitectura lógica, física y continuidad. & Medio & Ninguno & Gabriel Sanzana: comprobó la coherencia del capítulo con los Subdocs. 3 y 5. \\
4.1 & claude & Base narrativa de la arquitectura lógica y sus criterios. & Medio & Medio & Patricio Figueroa: revisó la separación de responsabilidades entre capas. \\
4.1.1 & claude & Organización de la selección de tecnologías de software. & Medio & Ninguno & Gabriel Sanzana: comprobó las tecnologías elegidas frente a las necesidades del caso. \\
4.1.1.1 a 4.1.1.5 & claude & Apoyo para comparar tecnologías, sistemas operativos y vigencia. & Medio & Ninguno & Gabriel Sanzana: contrastó las versiones, alternativas y fechas de soporte declaradas. \\
4.1.2 & claude & Apoyo en la síntesis de la justificación del estilo arquitectónico. & Medio & Ninguno & Patricio Figueroa: comprobó las decisiones de diseño y sus consecuencias operativas. \\
4.1.2.1 a 4.1.2.4 & claude & Desarrollo de criterios, decisiones ADR y patrones. & Medio & Ninguno & Patricio Figueroa: contrastó los criterios de evaluación y las decisiones registradas. \\
4.1.3 & claude & Apoyo para explicar la vista general en ocho capas. & Medio & Medio & Patricio Figueroa: revisó las dependencias y la lectura del esquema por capas. \\
4.1.3.1 a 4.1.3.4 & claude & Organización de capas, componentes, emplazamiento y ambientes. & Medio & Medio & Lucas Zamora: revisó la ubicación de cada componente y su relación con la faena. \\
4.1.4 & claude & Base de redacción de la descripción de las capas. & Medio & Ninguno & Lucas Zamora: revisó que cada capa tenga una responsabilidad reconocible. \\
4.1.4.1 a 4.1.4.8 & claude & Desarrollo del texto sobre presentación, servicios, datos y controles transversales. & Medio & Medio & Gabriel Sanzana: verificó los canales, servicios y controles de las ocho capas frente a la Figura \ref{fig:vista-general-sitco-aranda-ocho-capas} y a la Tabla \ref{tab:correspondencia-canales-modulos-plataformas-esquema-solucion}. \\
4.1.5 & claude & Apoyo comparativo para justificar los canales de terreno. & Medio & Ninguno & Lucas Zamora: revisó la adecuación de los canales al uso durante un corte. \\
4.1.5.1 & claude & Ajuste de claridad en la decisión de botonera y aplicación de terreno. & Medio & Ninguno & Lucas Zamora: revisó la decisión frente a los equipos y al alcance del Subdoc. 3. \\
4.1.6 & claude & Organización de la explicación de módulos y límites de contexto. & Medio & Ninguno & Lucas Zamora: validó la autonomía de cada módulo y sus dependencias. \\
4.1.6.1 & claude & Base explicativa de los módulos de la cadena del mineral. & Medio & Medio & Patricio Figueroa: verificó el recorrido desde la operación hasta la conciliación. \\
4.1.6.1.1 a 4.1.6.1.3 & claude & Redacción de las funciones de Operación Mina, Trazabilidad y Conciliación. & Medio & Medio & Gabriel Sanzana: validó los límites y los datos intercambiados entre los tres módulos. \\
4.1.6.2 a 4.1.6.5 & claude & Apoyo para describir mantenimiento, cumplimiento, módulos transversales e interfaces. & Medio & Ninguno & Lucas Zamora: comprobó las funciones y los contratos entre módulos. \\
4.1.7 & claude & Estructuración de la explicación de integración y eventos. & Medio & Ninguno & Lucas Zamora: contrastó la consistencia de las integraciones con los sistemas existentes. \\
4.1.7.1 a 4.1.7.9 & claude & Desarrollo de conectores, contratos y escenarios de integración. & Medio & Medio & Gabriel Sanzana: revisó los conectores, eventos y comportamiento ante fallas. \\
4.1.8 & claude & Apoyo en la síntesis de la explicación de la arquitectura de seguridad. & Medio & Ninguno & Gabriel Sanzana: comprobó la cobertura de amenazas y la coherencia con los accesos. \\
4.1.8.1 a 4.1.8.5 & claude & Apoyo para ordenar identidad, exposición, cifrado y controles. & Medio & Ninguno & Gabriel Sanzana: comprobó las poblaciones usuarias, excepciones y controles de seguridad. \\
4.1.9 & claude & Base descriptiva de la operación sin enlace. & Medio & Ninguno & Gabriel Sanzana: verificó las funciones que siguen disponibles en la faena. \\
4.1.9.1 a 4.1.9.2 & claude & Desarrollo de las reglas de continuidad y sincronización. & Medio & Ninguno & Lucas Zamora: revisó la recuperación de datos y la resolución de conflictos. \\
4.2 & codex y claude & Apoyo para explicar la arquitectura física y comparar el inventario. & Medio & Medio & Benjamín Rodríguez: verificó que los equipos y servicios del inventario coincidan con el diseño físico de 4.2. \\
4.2.1 & codex y claude & Organización de los implementos necesarios para la propuesta. & Medio & Ninguno & Marcelo Flores: comprobó el inventario frente al Formulario T-11. \\
4.2.1.1 a 4.2.1.4 & codex y claude & Apoyo en la descripción de nube, faena, terreno y licencias. & Medio & Medio & Kavon Kermani: contrastó las cantidades, garantías y condiciones de reposición. \\
4.2.2 & codex y claude & Comparación del emplazamiento y de la configuración de servicios de nube. & Medio & Medio & Kavon Kermani: comprobó la ubicación y la configuración frente al esquema físico. \\
4.2.2.1 a 4.2.2.2 & codex y claude & Apoyo para justificar el emplazamiento en faena, nube y servicios Azure. & Medio & Medio & Kavon Kermani: verificó los servicios Azure y los criterios de emplazamiento. \\
4.2.3 & codex y claude & Apoyo para comparar diagramas de despliegue, redes y continuidad. & Medio & Medio & Benjamín Rodríguez: validó la edición manual en draw.io y los enlaces representados. \\
4.2.3.1 a 4.2.3.5 & codex y claude & Organización de ambientes, enlaces, gabinetes, redundancia y respaldos. & Medio & Medio & Benjamín Rodríguez: revisó las rutas, redundancias y estrategia de respaldo del diseño. \\
4.2.4 & codex y claude & Apoyo de redacción para explicar el plan de capacidad. & Medio & Ninguno & Kavon Kermani: revisó la relación entre demanda y capacidad ofertada. \\
4.2.4.1 a 4.2.4.6 & codex y claude & Contraste de cálculos de volumetría, cómputo, almacenamiento y ancho de banda. & Bajo & Ninguno & Benjamín Rodríguez: verificó las fórmulas y los resultados frente al Anexo D de OVERFORE-Subdocumento4-Anexos. \\
4.3 & codex & Comparación del plano y del texto sobre el data center. & Medio & Medio & Kavon Kermani: validó la distribución y su correspondencia con los equipos. \\
4.3.1 & codex & Apoyo para describir el data center primario. & Medio & Medio & Kavon Kermani: verificó los servicios de la sala técnica y su redundancia. \\
4.3.1.1 a 4.3.1.4 & codex & Apoyo para explicar recinto, climatización, energía y protección contra incendios. & Medio & Medio & Marcelo Flores: comprobó el plano en draw.io y los requisitos físicos de la sala. \\
4.3.2 & codex & Base explicativa del data center secundario. & Medio & Medio & Benjamín Rodríguez: revisó la función de recuperación y la ubicación de la réplica. \\
4.3.2.1 a 4.3.2.2 & codex & Organización de replicación, conmutación, retorno y pruebas. & Medio & Bajo & Benjamín Rodríguez: revisó la secuencia de recuperación y sus criterios de prueba. \\
Referencias & claude & Apoyo para ordenar y uniformar las referencias citadas. & Medio & Ninguno & Patricio Figueroa: revisó la existencia de las fuentes y su correspondencia con las citas. \\
Anexo A & codex y claude & Apoyo para ordenar fichas y comparar especificaciones técnicas. & Medio & Ninguno & Marcelo Flores: comprobó las fichas frente a cantidades y componentes del T-11. \\
Anexo B & codex y claude & Organización de la correspondencia entre componentes lógicos y físicos. & Medio & Ninguno & Kavon Kermani: verificó el emplazamiento físico de los componentes lógicos. \\
Anexo C & codex y claude & Apoyo de redacción en responsabilidades y fuentes técnicas. & Medio & Ninguno & Marcelo Flores: revisó las responsabilidades y referencias de los equipos. \\
Anexo D & codex y claude & Contraste de fórmulas y cálculos de volumetría y dimensionamiento. & Bajo & Ninguno & Benjamín Rodríguez: revisó las unidades, operaciones y cantidades frente al T-11. \\
Anexo E & codex y claude & Apoyo en la síntesis de la matriz de controles de seguridad. & Medio & Ninguno & Lucas Zamora: validó las amenazas, controles y componentes cubiertos. \\
Anexo F & codex y claude & Apoyo comparativo para justificar tecnologías de software. & Medio & Ninguno & Patricio Figueroa: comprobó las alternativas, criterios y decisión de selección. \\
T-11 & codex y claude & Comparación del inventario de componentes, equipos y servicios ofertados. & Bajo & Ninguno & Marcelo Flores: verificó las cantidades, denominaciones y vínculos con las fichas técnicas. \\
Figuras revisadas de 4.1 a 4.3 & Codex & Adaptación de ilustraciones a carta u oficio, revisión de legibilidad y enlaces entre zonas y figuras. & Medio & Medio & Autorrevisión técnica y visual; revisión independiente del equipo aún no registrada. \\
\end{tabla}

\end{document}
